% macro conflict, first definition kept: \newcommand{\Ob}{\operatorname{Ob}}
% macro conflict, first definition kept: \newcommand{\Fl}{\mathrm{Fl}}
% macro conflict, first definition kept: \newcommand{\Fl}{\mathrm{Fl}}
% macro conflict, first definition kept: \newcommand{\Fl}{\mathrm{Fl}}
% macro conflict, first definition kept: \newcommand{\Ob}{\mathrm{Ob}}
\documentclass[11pt]{article}
\usepackage{fontspec}
\usepackage[french]{babel}
\usepackage{amsmath}
\usepackage{amssymb}
\usepackage{amsthm}
\usepackage{mathrsfs}
\usepackage{tikz-cd}
\usepackage{enumitem}

\newcommand{\printedpage}[1]{\marginpar{\footnotesize p.~#1}}
\newcommand{\msp}[1]{\marginpar{\footnotesize [#1]}}
\newcommand{\circled}[1]{\tikz[baseline=(c.base)]\node[draw,circle,inner sep=1pt](c){#1};}
\newcommand{\Cat}{\mathrm{Cat}}
\newcommand{\Ob}{\mathrm{Ob}\,}
\newcommand{\hot}{\mathrm{hot}}
\newcommand{\Loc}{\mathrm{Loc}}
\newcommand{\Hom}{\operatorname{Hom}}
\newcommand{\Fl}{\operatorname{Fl}}
\newcommand{\Hot}{\mathrm{Hot}}
\newcommand{\Ens}{\mathrm{Ens}}
\newcommand{\id}{\mathrm{id}}
\newcommand{\Cof}{\mathrm{Cof}}
\newcommand{\Mono}{\mathrm{Mono}}
\newcommand{\univ}{\mathrm{univ}}
\newcommand{\Fib}{\mathrm{Fib}}
\newcommand{\Fibf}{\mathrm{Fibf}}
\newcommand{\Parf}{\mathrm{Parf}}
\newcommand{\Catfib}{\mathrm{Catfib}}
\newcommand{\Fibsc}{\mathrm{Fibsc}}
\newcommand{\HOT}{\mathrm{HOT}}
\newcommand{\cF}{\mathcal{F}}
\newcommand{\cM}{\mathcal{M}}
\newcommand{\Comp}{\mathrm{Comp}}
\newcommand{\dottedimplies}{\mathrel{\tikz[baseline=-0.6ex]{\draw[dotted, double, double distance=1.5pt, -{Implies}] (0,0) -- (0.9,0);}}}
\newcommand{\inc}{\mathrm{inc}}
\newcommand{\dotimp}{\overset{\bullet}{\Longrightarrow}}
\newcommand{\dottedimpl}{\mathrel{\begin{tikzcd}[ampersand replacement=\&, column sep=small, baseline=-0.5ex]{}\arrow[r,Rightarrow,dotted]\&{}\end{tikzcd}}}
\newcommand{\cercle}[1]{\tikz[baseline=(c.base)]\node[draw,circle,inner sep=1pt](c){#1};}
\newcommand{\fourche}{\begin{tikzcd}[ampersand replacement=\&, row sep=0.5em, column sep=1em, baseline=(\tikzcdmatrixname-2-1.base)]
\newcommand{\carre}{\begin{tikzcd}[ampersand replacement=\&, row sep=0.5em, column sep=1em, baseline=(\tikzcdmatrixname-2-1.base)]
\newcommand{\uFl}{\underline{\mathrm{Fl}}}
\newcommand{\uCat}{\underline{\mathrm{Cat}}}
\begin{document}
% ---- from der-ch16_p001-005.tex
% ===================== page 1 (PDF) =====================
\printedpage{[non numérotée]}
% head: (none printed)

\begin{center}
{\sffamily\LARGE LES DÉRIVATEURS}

\bigskip

{\large\textbf{ALEXANDRE GROTHENDIECK}}

\vspace{3em}

{\large\textbf{Chapitre XVI}}

\medskip

{\large\textbf{Les localiseurs fondamentaux}}
\end{center}

\clearpage

% ===================== page 2 (PDF) =====================
\printedpage{[non numérotée]}
% head: (none printed)

Ce texte a été déchiffré et transcrit en \LaTeXe{} par M. Künzer. Il a été édité par M. Künzer, J. Malgoire, et G. Maltsiniotis. La transcription est aussi fidèle que possible au manuscrit. Seules quelques corrections mineures évidentes ont été effectuées. Pour les rares ajouts ou commentaires des éditeurs, ainsi que pour la numérotation originale des pages du manuscrit, les caractères de machine à écrire \texttt{[typewriter]} entre crochets sont utilisés.

Cette édition est provisoire. Les remarques, commentaires et corrections sont bienvenus. Envoyer un message à :

\textbf{[e-mail address removed]}

\hfill G. Maltsiniotis\hspace*{2cm}

\clearpage

% ===================== page 3 (PDF) =====================
\printedpage{1}
% head: A. Grothendieck : Dérivateurs, Ch. XVI, Édition des Universités de Montpellier II et Paris VII
% foot: Version 18 mars 2019

\msp{page 1}

\section*{1\quad Axiomes Loc(1) à Loc(3) (et variantes) : localiseurs fondamentaux}

J’ai envie de reprendre l’exposé des axiomes utiles pour un localiseur fondamental
\[
W \subset \mathrm{Fl}(\Cat),
\]
en me guidant cette fois plus sur la signification géométrique des axiomes, que sur les alias $\pm$ techniques qu’on rencontre (ou ne rencontre pas) pour les vérifier. Je révise donc dans cet esprit la présentation de XV §9.

\begin{description}
\item[Loc (1)] (\textit{saturation}).
\begin{itemize}
\item[\textit{a})] Iso $\subset W$.
\item[\textit{b})] Si deux parmi $f$, $g$, $gf$ sont dans $W$, aussi le troisième.
\item[\textit{c})] Si $gf$ et $fg$ sont dans $W$, $f$ et $g$ aussi.
\end{itemize}
\end{description}

Les conditions (\textit{a}) et (\textit{b}) toutes seules, je vais les appeler (pour un $W \subset \mathrm{Fl}\,\mathcal{M}$ général) \textit{saturation faible} du localiseur. Les conditions (\textit{a}), (\textit{b}), (\textit{c}), \textit{saturation modérée} (“mild saturation”). Autres variantes :

\begin{description}
\item[Loc (1 bis)] Les conditions (\textit{a}), (\textit{b}), (\textit{c}) du haut, plus
\begin{itemize}
\item[\textit{d})] $W$ stable par facteurs directs.
\end{itemize}
\end{description}

C’est bien sûr impliqué par :

\begin{description}
\item[Loc (1 ter)] $W$ est \textit{fortement saturé, i.e.} si $f \in \Cat$ est telle que $\hot_W(f)$ est un isomorphisme, alors $f \in W$.
\item[Loc (2)] (\textit{Objet final}). Si $X$ dans Cat a un objet final, alors $X$ est $W$-asphérique, \textit{i.e.} $X \longrightarrow e$ est dans $W$.
\end{description}

Parfois, cet axiome est conséquence de la variante affaiblie suivante.

\begin{description}
\item[Loc (2$_0$)] (Axiome du segment). $\boldsymbol{\Delta}^1$ est $W$-asphérique.
\end{description}

\msp{page 2}

\begin{description}
\item[Loc (3)] (\textit{localisation,} forme faible). Soit $f : X \longrightarrow S$ une flèche dans Cat $W$\textit{-asphérique, i.e.} telle que $\forall\ s \in \Ob S$, la catégorie localisée $X/s$ sur $S$ en $s$, $X/s = X \times_S S/s$, soit $W$-asphérique. Alors $f \in W$.\footnote[1]{Dans une version finale, il faudrait intervertir les notations pour Loc (3) et pour Loc (3 bis), car ce dernier m’apparaît comme le plus important des deux - il est satisfait pour tous les localiseurs qui ont été utilisées (à ma connaissance), et est valable pour tout localiseur de la forme $W_{\mathbf{D}}$, où $\mathbf{D}$ est un dérivateur (ou seulement un prédérivateur satisfaisant Der (4,5) \texttt{[il faut aussi Der (2) ; la numérotation des axiomes est celle du chapitre I, p. 68]}).} % mark printed as superscript (1); [sic] utilisées -> utilisés, agreement with "localiseurs"
\end{description}

\textbf{Rappel 1.} Loc (1, 2, 3) impliquent que l’on a
\[
W = W^\circ \qquad \textit{i.e. } f \in W \iff f^\circ \in W,
\]
\printedpage{2}
% head: A. Grothendieck : Dérivateurs, Ch. XVI, Édition des Universités de Montpellier II et Paris VII
% foot: Version 18 mars 2019
où on dénote par $X \longmapsto X^\circ$, $f \longmapsto f^\circ$ le foncteur de passage à la catégorie opposée, foncteur
\[
X \longmapsto X^\circ : \Cat \xrightarrow{\ \sim\ } \Cat
\]
qui est un \textit{automorphisme} de Cat (donc commute à toutes $\varinjlim$ et $\varprojlim$ ). Il s’ensuit que l’ensemble d’axiomes Loc (1, 2, 3) équivaut à Loc (1, 2$^{\mathrm{opp}}$, 3$^{\mathrm{opp}}$), où Loc (2$^{\mathrm{opp}}$) est l’axiome “dual” de Loc (2), savoir que tout objet de Cat ayant un objet initial est $W$-asphérique, et de même Loc (3$^{\mathrm{opp}}$) l’axiome que toute $f : X \longrightarrow S$ $W$\textit{-coasphérique, i.e.} telle que les catégories colocalisées $s\backslash X$ de $X$ sur $S$ soient $W$-asphériques, est dans $W$.

\textbf{Définition 1.1.} On appelle \textit{localiseur fondamental} une partie $W$ de Fl(Cat) satisfaisant les axiomes Loc (1, 2, 3) (\textit{i.e.} saturation modérée, axiome de l’objet final, axiome de localisation).

\msp{page 3}

\textbf{Rappel 1.2.} Si $W$ est un localiseur fondamental, alors tout \textit{homotopisme} $f : X \longrightarrow Y$ de Cat est dans $W$ (c’est pour cela essentiellement que l’on a imposé (\textit{c}) dans la saturation modérée Loc (1)). Si $f, g : X \rightrightarrows Y$ sont homotopes (\textit{i.e.} appartiennent à une même composante connexe de $\underline{\mathrm{Hom}}(X, Y)$), alors $\hot_W(f) = \hot_W(g)$ et $f \in W \iff g \in W$.

Voici la forme forte de Loc (3).

\begin{description}
\item[Loc (3 bis)] (\textit{Localisation,} forme forte). Soit $S$ dans Cat et $f : X \longrightarrow Y$ un morphisme dans $\Cat/S$ tel que “$f$ soit une $W$-équivalence localement sur $S$”, \textit{i.e.} pour tout $s \in \Ob S$, le morphisme induit
\[
f/s : X/s \longrightarrow Y/s
\]
pour les catégories localisées de $X$, $Y$ sur $S$ en $s$, soit dans $W$. Alors $f \in W$.
\end{description}

L’axiome Loc (3) n’est autre que le cas particulier de Loc (3 bis) relatif au cas où $Y = S$.

\textbf{Proposition 1.3.} \textit{Soit $W \subset \mathrm{Fl}(\mathcal{M})$. Alors on a les équivalences suivantes}
\[
\begin{aligned}
\text{Loc (1, 2, 3)} &\iff \text{Loc (1, $2_0$, $3'$)}\\
\text{Loc (1, 2, 3 bis)} &\iff \text{Loc (1, $2_0$, $3'$ bis),}
\end{aligned}
\]
\textit{où les axiomes} Loc ($3'$) \textit{et} Loc ($3'$ bis) \textit{se formulent ainsi :}

\msp{page 4}

\begin{description}
\item[Loc (3$'$)] \circled{\textit{a}} (Axiome des Cat-fibrations, forme directe faible). Soient $S$ dans Cat, $X$ une catégorie fibrée scindée sur $S$, telle que les fibres $X_s$ ($s \in \Ob S$) soient $W$-asphériques. Alors $X \longrightarrow S$ est dans $W$.

\circled{\textit{b}} Soit $X$ dans $\Cat/S$, alors $i_X : X \longrightarrow \Phi_S(X)$ est dans $W$\footnote[2]{(\textit{b}) concerne la \textit{comparaison} d’une $X$ quelconque sur $S$, avec la catégorie fibrée scindée $\Phi_S(X)$ sur $S$.}. % mark printed as superscript (2)
\end{description}

\clearpage

% ===================== page 5 (PDF) =====================
\printedpage{3}
% head: A. Grothendieck : Dérivateurs, Ch. XVI, Édition des Universités de Montpellier II et Paris VII
% foot: Version 18 mars 2019

\begin{description}
\item[Loc (3$'$ bis)] \circled{\textit{a}} (Axiome des Cat-fibrations, forme directe forte). Soient $S$ dans Cat, $X$, $Y$ deux catégories fibrées scindées sur $S$, $f : X \longrightarrow Y$ un $S$-foncteur cartésien compatible aux scindages, tel que $\forall\ s \in \Ob S$, le foncteur induit sur les fibres
\[
f_s : X_s \longrightarrow Y_s
\]
soit dans $W$. Alors $f \in W$.

\circled{\textit{b}} Comme Loc ($3'$) (\textit{b}).
\end{description}

\textbf{N.B.} En fait, moyennant Loc (1, 2), les énoncés Loc ($3'$), Loc ($3'$ bis) impliquent les énoncés similaires sans scindages, et sans même supposer $f$ cartésien dans Loc ($3'$ bis), \textit{cf.} démonstration ci-dessous.

\textsc{Démonstration.} \textit{a}) Loc (1, 2, 3) $\Longrightarrow$ Loc (1, $2_0$, $3'$). Il suffit de prouver Loc ($3'$). Prouvons Loc ($3'$ $a$). Pour ceci, sous les conditions de Loc ($3'$ $a$), prouvons que $X \longrightarrow S$ est $W$-coasphérique, \textit{i.e.} que les catégories
\[
s\backslash X = X \times_S s\backslash S
\]
sont $W$-asphériques. Il en résultera bien que $f \in W$ en vertu de Loc (3) appliqué à $f^\circ$ et de $W = W^\circ$. Mais
\msp{page 5}
considérons l’inclusion canonique
\[
i_s : X_s \lhook\joinrel\longrightarrow s\backslash X,
\]
il résulte du fait que $p : X \longrightarrow S$ est Cat-fibrant que ce foncteur est $W$-asphérique (il a même un adjoint à droite, \textit{i.e.} les catégories $X_s/u$, où $u \in \Ob(s\backslash X)$, ont un objet final). Donc par Loc (3) on a $i_s \in W$, et comme $X_s$ est $W$-asphérique par hypothèse, de même $s\backslash X$, q.e.d.

Prouvons Loc ($3'$ $b$) : On note que $i_X : X \longrightarrow \Phi_S(X)$ est $W$\textit{-coasphérique} (il a même un adjoint à gauche), donc $i_X^\circ$ est $W$-asphérique, donc dans $W$, donc $i_X^\circ \in W^\circ$ car $W = W^\circ$. % [sic] i_X^\circ \in W^\circ -> i_X \in W (or i_X \in W^\circ), apparently intended: from i_X^\circ \in W one gets i_X \in W^\circ = W

$a$ bis) Loc (1, 2, 3 bis) $\Longrightarrow$ Loc (1, $2_0$, $3'$ bis). Cette fois il faut prouver seulement Loc ($3'$ bis). Pour (\textit{b}), c’est déjà fait. Reste à prouver (\textit{a}). Si $f$ est comme dans l’énoncé Loc ($3'$ bis $a$), prouvons que pour tout $s$ dans $S$, le morphisme induit
\[
s\backslash f : s\backslash X \longrightarrow s\backslash Y
\]
est dans $W$. Comme $W = W^\circ$, on aura donc $(s\backslash f)^\circ \in W$, et appliquant Loc (3 bis) à $f^\circ$, on trouve $f^\circ \in W$, donc $f \in W$, et on aura fini. Mais $s\backslash f$ s’insère dans le carré commutatif
\msp{page 6}
\[
\begin{tikzcd}
X_s \arrow[r, "f_s"] \arrow[d, "i_s^X"'] & Y_s \arrow[d, "i_s^Y"] \\
s\backslash X \arrow[r, "s\backslash f"] & s\backslash Y,
\end{tikzcd}
\]

% ---- from der-ch16_p006-010.tex
\setcounter{footnote}{2}

\printedpage{4}
% head: A. Grothendieck : Dérivateurs, Ch. XVI, Édition des Universités de Montpellier II et Paris VII
où, on l'a vu dans $(a)$, $i_s^X$ et $i_s^Y$ sont dans $W$. Par hypothèse on a $f_s \in W$, donc $s\backslash f \in W$, q.e.d.

\textit{b}) Loc (1, $2_0$, $3'$) $\Longrightarrow$ Loc (1, 2, 3). Prouvons d'abord Loc (2). Appliquant Loc ($3'$ \textit{a}) à la projection
\[
\boldsymbol{\Delta}^1 \times X \longrightarrow X \qquad \text{(où $X \in \operatorname{Ob}\mathrm{Cat}$),}
\]
qui est fibrée scindée à fibres isomorphes à $\boldsymbol{\Delta}^1$, donc $W$-asphériques par Loc ($2_0$), on voit que cette projection est dans $W$. Cela permet d'appliquer les arguments standard d'homotopie (\textit{cf.} XII, p. 1 \ldots) pour conclure qu'un homotopisme est dans $W$, donc une catégorie contractile est $W$-asphérique, a fortiori une catégorie avec objet final (ou avec objet initial) est $W$-asphérique, \textit{i.e.} on a Loc (2).

Prouvons Loc (3). Comme Loc (1, 2, 3) $\Longleftrightarrow$ Loc (1, $2^\circ$, $3^\circ$), il suffit de prouver Loc ($3^\circ$), vu qu'on a prouvé Loc ($2^\circ$). Soit
\[
f : X \longrightarrow S \qquad \text{$W$-\textit{co}asphérique,}
\]
prouvons $f \in W$. Introduisons la catégorie \textit{fibrée} sur $S$ dont les fibres sont les $s\backslash X$, soit $\Phi_S(X)$, et considérons l'inclusion canonique
\[
i_X : X \longrightarrow \Phi_S(X).
\]
\msp{page 7}Par Loc ($3'$ \textit{b}) elle est dans $W$. D'autre part, les fibres $s\backslash X$ de $\Phi_S(X)$ sont $W$-coasphériques par hypothèse, donc par Loc ($3'$ \textit{a}) $\Phi_S(X) \longrightarrow S$ est dans $W$. Comme $i_X \in W$, il s'ensuit $(f : X \longrightarrow S) \in W$, q.e.d.

\textit{b} bis) Loc (1, $2_0$, $3'$ bis) $\Longrightarrow$ Loc (1, 2, 3 bis). On a déjà prouvé Loc (1, 2, 3) (vu que Loc ($3'$ bis) contient Loc ($3'$) comme cas particulier). Prouvons donc Loc (3 bis) en admettant déjà tout ça. Il suffira encore de le prouver sous forme duale, puisque
\[
\text{Loc (1, 2, 3 bis)} \iff \text{Loc(1, $2^\circ$, 3 bis$^\circ$),}
\]
\msp{page 8}et que Loc ($2^\circ$) est connu par Loc (1, 2, 3) $\Longleftrightarrow$ Loc (1, $2^\circ$, $3^\circ$). Il reste donc à prouver que si
\[
f : X \longrightarrow Y \qquad \text{dans } \mathrm{Cat}/S
\]
est telle que
\[
s\backslash f : s\backslash X \longrightarrow s\backslash Y
\]
est dans $W$ pour tout $s \in \operatorname{Ob} S$, alors $f \in W$. Considérons alors le carré commutatif
\[
\begin{tikzcd}
X \arrow[r, "i_X"] \arrow[d, "f"'] & \Phi_S(X) \arrow[d, "\Phi_S(f)"] \\
Y \arrow[r, "i_Y"] & \Phi_S(Y),
\end{tikzcd}
\]
où $i_X, i_Y \in W$ par Loc ($3'$ bis \textit{b}). L'hypothèse sur $f$ signifie que $\Phi_S(f)$ induit des $W$-équivalences sur les fibres, donc par Loc ($3'$ bis) $\Phi_S(f) \in W$, donc $f \in W$ par Loc (1), q.e.d.
% foot: Version 18 mars 2019

\printedpage{5}
% head: A. Grothendieck : Dérivateurs, Ch. XVI, Édition des Universités de Montpellier II et Paris VII
\section*{2\quad Axiomes Loc (4) et Loc (5) (des sommes directes, de connexité)\protect\footnote{Il s'agit ici d'axiomes qu'on peut considérer comme étant de nature plus ou moins “triviale”. Dans tous les cas à ma connaissance, ils sont bel et bien trivialement vérifiés.}}

\begin{description}[leftmargin=7em, labelwidth=6em]
\item[Loc (4)] (Axiome des \textit{sommes directes}). Pour une famille $(f_i)_{i\in I}$ de flèches dans Cat, si $f = \coprod f_i$, on a $f \in W$ si et seulement si $f_i \in W$ pour tout $i \in I$.
\end{description}

Cet axiome se décompose en deux, en le “si” et le “seulement si”, que je vais formuler sous la forme la plus économique :

\textbf{Proposition 2.1.} \textit{L'axiome} Loc (4) \textit{des sommes équivaut à la conjonction des deux axiomes suivants :}

\begin{description}[leftmargin=7em, labelwidth=6em]
\item[Loc (4 a)] \msp{page 9}Soit $(f_i)_{i\in I}$ une famille de flèches dans Cat, si les $f_i \in W$, alors $f = \coprod f_i \in W$.
\item[Loc (4 b)] Soit $S$ dans Cat, et $S_0 \hookrightarrow S$ un sommand direct (sous-catégorie à la fois ouverte et fermée). Soit $X$ dans Cat$/S$. Si $X \longrightarrow S$ est dans $W$, alors $X_0 = X|S_0 \longrightarrow S_0$ aussi. (En d'autres termes, l'inclusion $S_0 \hookrightarrow S$ est une “$W$-fibration”, \textit{i.e.} est dans $\mathrm{Fib}_W$.)
\end{description}

\textbf{Corollaire 2.2.} \textit{Considérons la variante de} Loc (4 a) \textit{obtenue en se bornant aux familles finies.} (\textit{Elle équivaut à} Loc (4) \textit{lorsque $W$ est stable par} $\varinjlim$ \textit{filtrantes,} cf. \textit{axiome} Loc (8) \textit{plus bas}). \textit{Cette variante équivaut à}

\begin{description}[leftmargin=7.5em, labelwidth=6.5em]
\item[Loc (4 a$_0$)] Tout objet $X$ de Cat est “$W$-cofibrant”, \textit{i.e.} $\emptyset \longrightarrow X$ est dans $\mathrm{Cof}_W$, \textit{i.e.} pour toute $f : Y \longrightarrow Y'$ qui est dans $W$, le morphisme correspondant $\mathrm{id}_X \amalg f : X \amalg Y \longrightarrow X \amalg Y'$ est dans $W$.
\end{description}

\textbf{Rappel 2.3.} De toutes façons, si $W$ est un localiseur fondamental (ce que nous supposerons désormais, sauf mention du contraire), l'analogue \textit{fini} de Loc (4 a) pour les produits est vrai, \textit{i.e.} si on a une famille \textit{finie} de flèches $(f_i)_{i\in I}$ dans Cat, si les $f_i$ sont dans $W$, alors $f = \prod f_i \in W$. \texttt{[Pour cela, il faut supposer que $W$ satisfait à l'axiome de localisation, forme \textit{forte}, Loc 3 bis.]} L'analogue pour une famille infinie n'est vérifié pratiquement pour aucun localiseur fondamental.

\msp{page 10}\textbf{Proposition 2.4.} \textit{Soit $W$ un localiseur fondamental. Alors les produits finis existent dans} $\mathrm{Hot}_W$\textit{, et le foncteur} $\mathrm{hot}_W : \mathrm{Cat} \longrightarrow \mathrm{Hot}_W \overset{\textit{déf}}{=} (\mathrm{Cat})W^{-1}$ \textit{y commute. Si $W$ satisfait à} Loc (4 a) (resp. Loc (4 a$_0$)), \textit{alors les sommes quelconques} (resp. \textit{finies}) \textit{existent dans} $\mathrm{Hot}_W$\textit{, et le foncteur} $\mathrm{hot}_W$ \textit{y commute. De plus, dans ce cas les sommes sont distributives par rapport aux produits (vu qu'il en est ainsi dans} Cat).

La démonstration résulte simplement du résultat général :

\textbf{Lemme 2.5.} \textit{Soit $(\mathcal{M}, W)$ une catégorie de modèles ($W$ partie quelconque de} $\mathrm{Fl}\,\mathcal{M}$). \textit{Supposons que les produits} (resp. \textit{les produits finis,} resp. \textit{les produits binaires}) \textit{existent dans $\mathcal{M}$, et que de plus pour toute famille $(f_i)_{i\in I}$ de flèches de $\mathcal{M}$} (resp. \textit{toute famille finie,}
% foot: Version 18 mars 2019
\printedpage{6}
% head: A. Grothendieck : Dérivateurs, Ch. XVI, Édition des Universités de Montpellier II et Paris VII
resp. \textit{famille binaire}) \textit{$f_i \in W$ pour tout $i \in I$ implique $f = \prod f_i \in W$. Alors les produits du type envisagé existent dans} $\mathrm{Ho}_W = \mathcal{M}W^{-1}$\textit{, et le foncteur} $\mathrm{ho}_W : \mathcal{M} \longrightarrow \mathcal{M}W^{-1}$ \textit{y commute. Enfin, si $\mathcal{M}$ admet un objet final $e$, alors $e$ est aussi objet final de} $\mathrm{Ho}_W$.

\texttt{[Il n'est pas clair que l'assertion soit vraie dans cette généralité pour les produits \textit{infinis} ; elle est vraie si $W$ est la classe des équivalences faibles d'une \textit{catégorie de modèles de Quillen.}]}

\msp{page 11}On en conclut les énoncés duaux concernant les sommes directes et les objets initiaux.

La démonstration est un cas particulier du résultat bien utile suivant, appliqué au couple de foncteurs adjoints
\[
\begin{tikzcd}
\mathcal{M} \arrow[r, shift left, "\mathrm{diag}"] & \mathcal{M}^I \arrow[l, shift left, "\text{\scriptsize\shortstack{produits\\ dans $\mathcal{M}$\\ indexés par $I$}}"]
\end{tikzcd}
\qquad\qquad :
\]

\textbf{Lemme 2.6}\footnote{\textit{cf.} XII, lemme page 248.}. \textit{Soient $(\mathcal{M}, W)$, $(\mathcal{M}', W')$ deux catégories de modèles, et}
\[
\begin{tikzcd}
\mathcal{M} \arrow[r, shift left, "\varphi"] & \mathcal{M}' \arrow[l, shift left, "\psi"]
\end{tikzcd}
\]
\textit{un couple de foncteurs adjoints compatibles aux localiseurs,} i.e. \textit{tels que}
\[
\varphi(W) \subset W', \qquad \psi(W') \subset W.
\]
\textit{Alors les foncteurs localisés}
\[
\begin{tikzcd}
\mathcal{M}W^{-1} \arrow[r, shift left, "\bar{\varphi}"] & \mathcal{M}'W'^{-1} \arrow[l, shift left, "\bar{\psi}"]
\end{tikzcd}
\]
\textit{sont également adjoints, les flèches d'adjonction pour $\bar{\varphi}\bar{\psi}$ et $\bar{\psi}\bar{\varphi}$ étant déduites de celles pour $\varphi\psi$ et $\psi\varphi$ par composition avec les foncteurs de localisation $\mathcal{M} \longrightarrow \mathcal{M}W^{-1}$, $\mathcal{M}' \longrightarrow \mathcal{M}'W'^{-1}$.}

\begin{description}[leftmargin=7em, labelwidth=6em]
\item[Loc (5)] \msp{page 12}(Axiome de \textit{connexité} ou de \textit{non-trivialité}).

Soit $f \in W$, alors $\pi_0(f)$ est bijectif.
\end{description}

Cette propriété de $W$ est d'une utilisation constante en pratique, quand on travaille avec $W$ au lieu de se borner à construire des $W$ et faire joujou avec les axiomes et leur stabilité par passage à des $\underline{\mathrm{Hom}}(I, M)$ (où ici $\mathcal{M} = \mathrm{Cat}$). D'autre part, on peut aussi le regarder comme un axiome de \textit{non-trivialité} pour $\mathrm{Hot}_W$. Je rappelle à ce sujet :

\textbf{Proposition 2.7.} \textit{Soit $W$ un localiseur fondamental. Conditions équivalentes :}
\begin{itemize}[leftmargin=3em]
\item[a)] $W$ \textit{satisfait l'axiome de connexité} Loc (5).
\item[b)] \textit{Le foncteur} $\pi_0 : \mathrm{Cat} \longrightarrow \mathrm{Ens}$ \textit{se factorise par un foncteur}
% foot: Version 18 mars 2019
\printedpage{7}
% head: A. Grothendieck : Dérivateurs, Ch. XVI, Édition des Universités de Montpellier II et Paris VII
\[
\begin{tikzcd}
\mathrm{Cat} \arrow[r, dotted] \arrow[rr, dotted, bend right, "\pi_0"'] & \mathrm{Hot}_W \arrow[r, "\pi_0"] & \mathrm{Ens}.
\end{tikzcd}
\]
\item[c)] \textit{La catégorie} $\mathrm{Hot}^*_W = \mathrm{Hot}_W\backslash\{\emptyset\}$ \textit{n'est pas équivalente à la catégorie ponctuelle.}
\item[c$'$)] \textit{La catégorie} $\mathrm{Hot}_W$ \textit{n'est pas équivalente à $e$ ni à $\boldsymbol{\Delta}^1$.}\footnote{\textbf{NB} Si $W$ est fortement saturée (condition Loc (1 ter), alors les conditions ci-contre (\textit{i.e.} Loc (5)) équivalent à ceci : % [sic] "(condition Loc (1 ter)," -> "(condition Loc (1 ter)),", parenthèse ouvrante non refermée
\begin{itemize}[leftmargin=2.5em]
\item[d)] \textit{$W$ n'est égal ni à Fl}(Cat) \textit{tout entier, ni à Fl}(Cat$^*$) $\cup$ $\{\mathrm{id}_\emptyset\}$ (i.e. \textit{Fl}(Cat)$\backslash W$ \textit{n'est ni vide, ni égal à l'ensemble des flèches $\emptyset \longrightarrow X$ avec $X \neq \emptyset$}).
\end{itemize}
Si $W$ satisfait Loc (7 D$_{W_\omega}$), alors $W$ satisfait Loc (5) si et seulement si $W \neq \mathrm{Fl}(\mathrm{Cat})$ (\textit{cf.} XII, th. 2, p. 20).}
\end{itemize}

Pour la démonstration, et pour d'autres formes équivalentes de l'axiome de connexité, \msp{page 13}je renvoie à XV 9.12 (p. 133 bis). Pour le nom “axiome de connexité”, je peux offrir l'explication suivante 2.8 et 2.9.

\textbf{Corollaire 2.8.} \textit{Si l'axiome de connexité} Loc (5) \textit{est satisfait, alors les objets $W$-asphériques sont} 0\textit{-connexes} (\textit{tautologie}), \textit{et la réciproque est vraie si $W$ est fortement saturé.} (\textit{Ou aussi, si $W$ satisfait l'axiome des fibrations} Loc (7) \textit{ci-dessous, puisque par} XII, th. 2, p. 20, \textit{dans ce cas,} Loc (5) \textit{équivaut à $W \neq \mathrm{Fl}(\mathcal{M})$ \ldots}).

Car si l'axiome de connexité n'est pas satisfait, alors pour toute $X$ non vide dans Cat, $X \longrightarrow e$ est un isomorphisme, donc ($W$ étant fortement saturé) $X \longrightarrow e$ est dans $W$, \textit{i.e.} $X$ est $W$-asphérique, donc il y a des objets $W$-asphériques non 0-connexe. Notons aussi : % [sic] "non 0-connexe" -> "non 0-connexes", accord avec "des objets"

\textbf{Corollaire 2.9.} \textit{Supposons que $W$ satisfait à l'axiome} Loc (4) \textit{des sommes directes} (\textit{ou seulement à} Loc (4 b)). \textit{Pour que $W$ satisfasse l'axiome de connexité, il faut et il suffit que pour toute flèche $f : X \longrightarrow Y$ avec $Y$} 0\textit{-connexe, si $f$ est dans $W$, alors $X$ est} 0\textit{-connexe.} (\textit{C'est nécessaire bien sûr sans supposer} Loc (4 b) \textit{ni rien du tout.})

\msp{page 14}\texttt{[Les numéros 2.7, 2.8 et 2.9 apparaissent deux fois.]}

\textbf{2.7.} Considérons maintenant le foncteur composé
\[
\mathrm{Ens} \xrightarrow{I \longmapsto I_{\mathrm{Cat}}} \mathrm{Cat} \xrightarrow{\quad \mathrm{hot}_W \quad} \mathrm{Hot}_W, \qquad I \longmapsto I_{\mathrm{Hot}_W} \text{ ou } I_W, \tag*{(2.7.1)}
\]
où le premier foncteur est celui qui associe à l'ensemble $I$ la catégorie discrète $I_{\mathrm{Cat}}$ ayant $I$ comme ensemble d'objets. (Le plus souvent, on identifiera $I_{\mathrm{Cat}}$ avec $I$ dans les notations, tout comme pour un ensemble ordonné $J$, on identifie $J$ et la “catégorie ordonné” % [sic] "catégorie ordonné" -> "catégorie ordonnée", accord
définie par $J$. Ici $I$ peut être considéré comme muni de l'ordre discret.) On a d'ailleurs
\[
I_{\mathrm{Cat}} \simeq I \times e \ \ (\overset{\text{déf}}{=} \text{somme de $I$ copies de $e$}), \tag*{(2.7.2)}
\]
donc quand Loc (4 a) est satisfait, donc $\mathrm{hot}_W$ commute aux sommes directes, on aura un isomorphisme canonique
\[
\begin{aligned}
I_W \simeq I \times e_W \quad &\text{(où $e_W$ est l'objet final de $\mathrm{Hot}_W$,} \\
&\text{\textit{i.e.} $e$ considéré comme objet de $\mathrm{Hot}_W$).}
\end{aligned} \tag*{(2.7.3)}
\]
% foot: Version 18 mars 2019

\printedpage{8}
% head: A. Grothendieck : Dérivateurs, Ch. XVI, Édition des Universités de Montpellier II et Paris VII
\textbf{Proposition 2.8.} \textit{Supposons l'axiome de connexité} Loc (5) \textit{satisfait, d'où un foncteur}
\[
\pi_0 : \mathrm{Hot}_W \longrightarrow \mathrm{Ens}. \tag*{(2.8.1)}
\]
\textit{Ce foncteur est} adjoint à gauche \textit{du foncteur $I \longmapsto I_W$} (2.7.1), i.e. \textit{on a une bijection fonctorielle en $(X, I)$ dans} $\mathrm{Hot}_W \times \mathrm{Ens}$ :
\msp{page 15}
\[
\mathrm{Hom}_{\mathrm{Ens}}(\pi_0(X), I) \simeq \mathrm{Hom}_W(X, I_W). \tag*{(2.8.2)}
\]
\textit{Le foncteur $\pi_0$ est également} adjoint à droite \textit{du} même \textit{foncteur $I \longmapsto I_W$,} i.e. \textit{on a un bijection fonctorielle} % [sic] "un bijection" -> "une bijection", accord
\[
\mathrm{Hom}_{\mathrm{Ens}}(I, \pi_0(X)) \simeq \mathrm{Hom}_W(I_W, X). \tag*{(2.8.3)}
\]

Prenant $I = e$ (ensemble ponctuel), on trouve le :

\textbf{Corollaire 2.9.} \textit{On a l'isomorphisme fonctoriel en $X$ dans} $\mathrm{Hot}_W$
\[
\boxed{\mathrm{Hom}_W(e_W, X) \simeq \pi_0(X).} \tag*{(2.8.4)}
\]

\textbf{Commentaire 2.10.} La formule d'adjonction (2.8.2) est d'ailleurs familière, chaque fois qu'on a une catégorie $\mathcal{H}$ (telle Cat, où $\mathrm{Hot}_W$ \textit{etc.}) % [sic] "où" -> "ou", conjonction de coordination
munie d'un foncteur $\pi_0$, $\pi_0(X)$ correspondant à l'intuition d'une décomposition d'un objet en “composantes connexes”. (Je ne vais pas expliciter ici les conditions générales sous lesquelles on définit un tel foncteur.) Par contre, l'adjonction en sens inverse de ces deux mêmes foncteurs est un phénomène tout à fait spécial à des catégories \msp{page 16}du type $\mathrm{Hot}_W$. Si on tient compte de la description de $I_W$ comme $I \times e_W$, la formule d'adjonction inhabituelle (2.8.3) équivaut, bien sûr, à la formule (2.8.4) - elle s'en déduit en “élevant à la puissance $I$”. On peut renverser la vapeur, et dans une catégorie générale $\mathcal{H}$, \textit{définir} un foncteur $\pi_0 : \mathcal{H} \longrightarrow \mathrm{Ens}$ par le formule (2.8.4), % [sic] "par le formule" -> "par la formule", accord
\[
\pi_0(X) \overset{\text{déf}}{=} \mathrm{Hom}(e_{\mathcal{H}}, X) \qquad (= \Gamma_{\mathcal{H}}(X), \text{ avec une notation familière}),
\]
où $e_{\mathcal{H}}$ est un objet final de $\mathcal{H}$. On aura donc, en posant $I_{\mathcal{H}} = I \times e_{\mathcal{H}}$, la formule
\[
\mathrm{Hom}(I_{\mathcal{H}}, X) \xrightarrow{\ \sim\ } \mathrm{Hom}(I, \pi_0(X))
\]
(\textit{i.e.} la formule (2.8.3)), disant que l'on a un couple $(I_{\mathcal{H}}, \pi_0)$ de foncteurs adjoints. Supposons d'autre part que pour tout ensemble $I$, l'application canonique d'adjonction
\[
I \longrightarrow \pi_0(I_{\mathcal{H}}) = \mathrm{Hom}_{\mathcal{H}}(e_{\mathcal{H}}, e_{\mathcal{H}} \times I)
\]
soit \textit{bijective}\footnote{ça a tendance à être ultrafaux si $\mathcal{H}$ est additive. Par contre, vrai si $\mathcal{H}$ est une catégorie à sommes distributives par rapport aux produits (voire universelles), dont l'objet final est \textit{connexe} \ldots}, \textit{i.e.} que $I \longmapsto I_{\mathcal{H}}$ soit \textit{pleinement fidèle.} On définit alors, pour tout $I$, une application canonique
% foot: Version 18 mars 2019

% ---- from der-ch16_p011-015.tex
\printedpage{9}
% head: A. Grothendieck : Dérivateurs, Ch. XVI, Édition des Universités de Montpellier II et Paris VII
% foot: Version 18 mars 2019

\[
\begin{tikzcd}[column sep=small]
\Hom(X,I_{\mathcal{H}}) \arrow[rr] \arrow[dr,dotted,"\pi_0^{X,I_{\mathcal{H}}}"'] & & \Hom(\pi_0(X),I) \arrow[dl,dotted,"\sim"'] \\ % ?? uncertain: label of the dotted right arrow, a tilde printed along the arrow -- \sim
& \Hom(\pi_0(X),\pi_0(I_{\mathcal{H}})). &
\end{tikzcd}
\]

La “bi-adjonction” des foncteurs $\pi_0$ et $I\longmapsto I_{\mathcal{H}}$ revient à dire que cette flèche est bijective,
\msp{17}
ou encore, que la flèche $\pi_0^{X,I_{\mathcal{H}}}$ :
\[
\Hom(X,I_{\mathcal{H}})\longrightarrow\Hom(\Gamma_{\mathcal{H}}(X),\Gamma_{\mathcal{H}}(I_{\mathcal{H}})\simeq I)
\]
soit bijective.

\textbf{Corollaire 2.11.} \textit{Le foncteur} $I\longmapsto I_W$ \textit{de} Ens \textit{dans} $\Hot_W$ \textit{est pleinement fidèle,} i.e. \textit{l'homomorphisme d'adjonction}
\[
\pi_0(I_W)\longrightarrow I
\]
\textit{est bijectif.} (\textit{Où encore,} % [sic] Où encore -> Ou encore, the accent makes no sense here
\textit{$\pi_0$ est un foncteur le localisation.}) % [sic] le localisation -> de localisation, wrong word

En effet, cette bijectivité est évidente quand on travaille dans Cat, et s'en déduit aussitôt dans $\Hot_W$.

\textsc{Démonstration de} 2.8. La formule d'adjonction (2.8.2) résulte de la formule correspondante pour Cat, évidente, et du lemme 2.6 sur les foncteurs adjoints entre catégories de modèles. L'adjonction dans l'autre sens se réduit à la formule (2.8.4). On a des flèches canoniques
\[
\begin{tikzcd}
\Hom_W(e_W,X) \arrow[r,shift left,"\alpha"] & \pi_0(X), \arrow[l,shift left,"\beta"]
\end{tikzcd}
\]
où $\alpha$ est $\pi_0^{e_W,X}$, compte tenu que $\pi_0(e_W)=\pi_0(e_{\Cat})=e_{\Ens}$, et $\Hom(e_{\Ens},\pi_0(X))\simeq\pi_0(X)$. La flèche $\beta$ est la flèche composée
\[
\Ob X\xrightarrow{\ \sim\ }\Hom_{\Cat}(e_{\Cat},X)\xrightarrow{\mathrm{hot}_W}\Hom_W(e_W,X)
\]
\msp{18}
en notant que cette flèche passe au quotient en $\pi_0(X)\longrightarrow\Hom_W(e_W,X)$, vu que deux applications homotopes de $e_{\Cat}$ dans $X$ sont égales dans $\Hot_W$. Je dis que
\[
\beta\alpha=\id\ ,\qquad\qquad \alpha\beta=\id\ .
\]
La relation $\alpha\beta=\id$ revient à dire que la flèche $i_x : e_{\Cat}\longrightarrow X$ définie par $x\in\Ob X$, donne $\pi_0(i_x) : e_{\Ens}\longrightarrow\pi_0(X)$ dont la valeur est la classe de $x$ dans $\pi_0(X)$, c'est tautologique. Pour prouver que $\beta\alpha=\id$, il suffit alors de prouver que $\beta$ est surjective, \textit{i.e.} le

\textbf{Corollaire 2.12.} \textit{Pour tout $X$ dans} Cat, \textit{l'application}
\[
\underbrace{\Hom_{\Cat}(e_{\Cat},X)}_{\simeq\Ob X}\longrightarrow\Hom_W(e_W,X)
\]
\printedpage{10}
% head: A. Grothendieck : Dérivateurs, Ch. XVI, Édition des Universités de Montpellier II et Paris VII
% foot: Version 18 mars 2019
\textit{est} surjective (\textit{moyennant l'axiome de connexité} Loc (5)).

\textsc{Démonstration de} 2.12. Si $X=\emptyset$, l'assertion revient à dire que $\Hom_W(e_W,\emptyset)=\emptyset$, mais plus généralement, on a le

\textbf{Corollaire 2.13.} (Moyennant Loc (5).) \textit{Soit $X$ non vide dans} Cat, \textit{alors} $\Hom_W(X,\emptyset)=\emptyset$. (\textit{Donc $\emptyset$ est un objet initial} strict \textit{de} $\Hot_W$.)

En effet, si on avait $X\longrightarrow\emptyset$ dans $\Hot_W$,
\msp{13$'$ (il y a un décalage dans la numérotation)}
on en déduirait $\pi_0(X)\longrightarrow\emptyset$, ou $\pi_0(X)=\emptyset$, absurde\,!

Reste le cas où $X$ est non vide. Cela résulte du fait général :

\textbf{Lemme 2.14.} \textit{Soit $\mathcal{M}$ une catégorie ayant un objet initial} strict (\textit{si} $\Hom_{\mathcal{M}}(X,\emptyset_{\mathcal{M}})\neq\emptyset$, \textit{on a} $X=\emptyset_{\mathcal{M}}$), \textit{un objet final $e$, et telle que pour tout $X\neq\emptyset_{\mathcal{M}}$, $X\longrightarrow e$ ait une section,} i.e. $\Hom(e,X)\neq\emptyset$. \textit{Soit} $W\subset\Fl\mathcal{M}$
\begin{itemize}[leftmargin=3em,labelwidth=2.5em,align=right]
\item[1$^{\circ}$)] \textit{ne contenant pas de flèches} $\emptyset_{\mathcal{M}}\longrightarrow X$ \textit{avec} $X\neq\emptyset_{\mathcal{M}}$, \textit{et}
\item[2\,$^{\circ}$)] \textit{tel que pour tout diagramme}
\begin{tikzcd}[baseline=(\tikzcdmatrixname-1-1.base)]
X_1 & X_2 \arrow[l,"u_2"'] \\
e \arrow[u,"u_1"] \arrow[ur,dotted,"s"'] &
\end{tikzcd}
\textit{avec} $u_2\in W$ $\exists$ $s$ dans $\mathcal{M}$ \textit{tel que} $u_2s=u_1$

\textit{dans} $\mathcal{M}W^{-1}$.
\end{itemize}
\textit{Alors}
\begin{itemize}[leftmargin=3em,labelwidth=2.5em,align=right]
\item[a)] $\emptyset_{\mathcal{M}}$ \textit{est un objet initial} strict \textit{de} $\mathcal{M}W^{-1}$.
\item[b)] \textit{Pour tout} $X\neq\emptyset_{\mathcal{M}}$, $\Hom_{\mathcal{M}}(e,X)\longrightarrow\Hom_W(e,X)$ \textit{est surjectif.}
\end{itemize}

\textsc{Démonstration}. Que $\emptyset_{\mathcal{M}}$ soit objet initial de $\mathcal{M}W^{-1}$ a été vu (dernière assertion du lemme 2.5, forme duale), prouvons qu'il est strict, \textit{i.e.} que si $\Hom_W(X,\emptyset_{\mathcal{M}})\neq\emptyset$, on a $X=\emptyset_{\mathcal{M}}$. En effet, si on avait $u : X\longrightarrow\emptyset_{\mathcal{M}}$ dans $\mathcal{M}W^{-1}$, $u$ serait représenté par
\[
X=X_0\xleftarrow{\ W\ }X_1\longrightarrow X_2\xleftarrow{\ W\ }\cdots X_{2n-2}\longrightarrow X_{2n-1}\xleftarrow{\ W\ }X_{2n}=\emptyset_{\mathcal{M}}.
\]
On procède par récurrence sur la demi-longueur de cette chaîne pour voir que $X=\emptyset_{\mathcal{M}}$. En effet, l'hypothèse sur $W$ implique que $X_{2n-1}=\emptyset_{\mathcal{M}}$, puis celle que $\emptyset_{\mathcal{M}}$ est strict implique
$X_{2n-1}=\emptyset_{\mathcal{M}}$, d'où la conclusion. % [sic] X_{2n-1} -> X_{2n-2}, the strictness of the initial object gives X_{2n-2}

\msp{14$'$}
Prouvons la surjectivité de $\Hom_{\mathcal{M}}(e,X)\longrightarrow\Hom_W(e,X)$ pour $X\neq\emptyset_{\mathcal{M}}$, en montrant que tout chemin réduit de longueur 2n+1 est “homotope” à un chemin réduit de longueur $2n-1$ :
\[
e=X_0\longrightarrow X_1\xleftarrow{\ W\ }X_2\longrightarrow X_3\longleftarrow\ \cdots\ \longleftarrow X_{2n}\longrightarrow X_{2n+1}=X.
\]
En effet, on peut remplacer la section
\[
\begin{tikzcd}
e=X_0 \arrow[r,"u_1"'] \arrow[rr,dotted,bend right,"s"'] & X_1 & X_2 \arrow[l,"W"',"u_2"] \arrow[r,"u_3"'] & X_3 & X_4 \arrow[l,"W"']
\end{tikzcd}
\]
\printedpage{11}
% head: A. Grothendieck : Dérivateurs, Ch. XVI, Édition des Universités de Montpellier II et Paris VII
% foot: Version 18 mars 2019
par
\[
\begin{tikzcd}
e \arrow[r,"u_3s"'] & X_3 & \cdots, \arrow[l,"W"']
\end{tikzcd}
\]
où $s$ est une section de $X_2$ sur $e$ telle que
\[
u_2s=u_1\qquad\quad\textit{dans}\qquad\quad\mathcal{M}W^{-1}\ .
\]
Cela prouve donc par récurrence sur $n$ que le chemin est homotope à un chemin de longueur 1, $e\longrightarrow X$, q.e.d.

Il faut donc, dans le cas actuel, vérifier 1$^{\circ}$) (c'est essentiellement le corollaire 2.12.), et 2$^{\circ}$)
\[
\begin{tikzcd}
X_1 & X_2 \arrow[l,"u_2\in W"'] & \\
e \arrow[u,"u_1"] \arrow[ur,dotted,"s"'] & & .
\end{tikzcd}
\]
Mais comme $u_2\in W$, on a $\pi_0(u_2)$ isomorphisme, donc
\msp{15$'$}
surjectif, donc $\exists$ $s$ tel que $u_2s$ soit homotope à $u_1$ dans Cat, ce qui implique $u_2s=u_1$ dans $\Hot_W$, q.e.d.

\textbf{Commentaire final.} L'axiome de connexité Loc (5) diffère de tous les autres axiomes Loc (1) à Loc (8) (mais à l'exception de Loc (8 bis) d'accessibilité) par le fait que ce n'est pas un axiome de stabilité. Les autres axiomes Loc (1) à Loc (8) (et aussi Loc (8 bis) d'accessibilité et même les axiomes ultérieures % [sic] ultérieures -> ultérieurs, agreement with axiomes
d'exactitude~\footnote[7]{il n'y en aura en fait pas besoin, grâce à la providentielle théorie de K. S. \textsc{Brown}\,!}) sont satisfaits si $W=\Fl\mathcal{M}$ - l'axiome de connexité est donc le seul qui s'oppose à la \textit{trivialité} de la théorie homotopique définie par $W$. C'est dire qu'il ne peut être conséquence de l'un des autres. Par contre, l'axiome Loc (4) des sommes directes n'est pas indépendant des autres. Signalons à ce propos (rappel) :
\[
\left\{
\begin{array}{rl}
\text{Loc (4 a$_0$)} & \text{est conséquence de l'axiome Loc (6 a) du carré cocartésien (forme directe).}\\[1ex]
\text{Loc (4 a)} & \text{est conséquence de Loc (6 a) précédent}\\
 & \text{+ l'axiome des limites inductives filtrantes Loc (8).}\\[3ex]
\text{Loc (4 b)} & \text{est conséquence de l'axiome Loc (7) des fibrations. Donc}\\[1ex]
\text{Loc (4)} & \text{est conséquence de Loc (6 a, 7, 8).}
\end{array}
\right.
\]

\printedpage{12}
% head: A. Grothendieck : Dérivateurs, Ch. XVI, Édition des Universités de Montpellier II et Paris VII
% foot: Version 18 mars 2019
\msp{16$'$}

\section*{3\quad Axiome Loc (6) du carré cocartésien (ou de “Mayer-Vietoris”, ou des $W$-cofibrations)~\footnote[8]{C'est le nom “\textit{axiome des $W$-cofibrations}” qui m'apparaît à présent comme le meilleur, -- c'est ce nom qui me paraît le mieux saisir la signification conceptuelle et non technique de l'axiome.}}

Je vois plusieurs variantes, et n'ai pas encore l'esprit totalement net sur les implications mutuelles. Je commence par la variante plus compliquée, mais dont on aura besoin en tous cas :

\begin{itemize}[leftmargin=5em,labelwidth=4em,align=left]
\item[\textbf{Loc 6}] (Axiome des $W$-cofibrations, ou lemme des 5, ou axiome des homomorphismes de carrés cocartésiens). Considérons dans Cat un diagramme commutatif cubique
\[
\begin{tikzcd}[row sep=small,column sep=small]
X_0 \arrow[rr] \arrow[dr] \arrow[dd,"f_0"'] & & X_2 \arrow[dr] \arrow[dd,"f_2" near start] & \\
& X_1 \arrow[rr,crossing over] & & X_3 \arrow[dd,"f_3"] \\
Y_0 \arrow[rr] \arrow[dr] & & Y_2 \arrow[dr] & \\
& Y_1 \arrow[rr] \arrow[from=uu,crossing over,"f_1"' near start] & & Y_3,
\end{tikzcd}
\]
\end{itemize}
homomorphisme $f=(f_0,f_1,f_2,f_3)$ d'un diagramme carré $X$ dans un autre $Y$. On suppose que $X$ provient d'un objet $X_3$ de Cat de deux ouverts $X_1$, $X_2$ de $X_3$ \textit{recouvrant} $X_3$, et de leur intersection $X_0=X_1\cap X_2$, et de même pour $Y$. (Appelons ces diagrammes dans Cat des diagrammes MV (“\textit{ouverts}”) - “\textit{Mayer-Vietoris}”.) On a alors ce qui suit~\footnote[9]{Il est bon d'introduire des propriétés du diagramme
\begin{itemize}[leftmargin=3em,labelwidth=2.5em,align=right]
\item[1$^{\circ}$)] $f_0\in W$.
\item[2$^{\circ}$)] $f_1$ et $f_2$ dans $W$.
\item[3$^{\circ}$)] $f_3$ dans $W$.
\end{itemize}
et de formuler les parties a) b) c) de Loc (6) ainsi
\begin{itemize}[leftmargin=3em,labelwidth=2.5em,align=right]
\item[a)] 1$^{\circ}$ + 2$^{\circ}$ $\Longrightarrow$ 3$^{\circ}$.
\item[b)] 2$^{\circ}$ + 3$^{\circ}$ $\Longrightarrow$ 1$^{\circ}$.
\item[c)] 1$^{\circ}$ + 3$^{\circ}$ $\Longrightarrow$ 2$^{\circ}$.
\end{itemize}} :
\begin{itemize}[leftmargin=3em,labelwidth=2.5em,align=right]
\item[a)] $f_0,f_1,f_2\in W\Longrightarrow f_3\in W$~\footnote[10]{N.B. il semble que dans Loc (6 a), il suffit de l'énoncé en supposant $f_0$, $f_2$ des \textit{isomorphismes,} \textit{cf.} p. 22.}.
\item[b)] $f_1,f_2,f_3\in W\Longrightarrow f_0\in W$.
\item[c)] $f_0,f_3\in W\Longrightarrow f_1,f_2\in W$~\footnote[11]{Il ne faudrait pas inclure c dans Loc 6 \textit{cf.} p. 26.}.
\end{itemize}
\noindent\footnote[12]{N.B. On verra plus bas (3.3) que Loc 6 c implique déjà Loc (6 a, b), \textit{i.e.} Loc (6) tout entier. Et si $W$ satisfait Loc (1 bis), \textit{i.e.} stable par facteurs directs, alors Loc (6 a, b) $\Longrightarrow$ Loc (6 c).}.

\printedpage{13}
% head: A. Grothendieck : Dérivateurs, Ch. XVI, Édition des Universités de Montpellier II et Paris VII
% foot: Version 18 mars 2019
\msp{17$'$}
\noindent\footnote[13]{Je vais désigner par Loc (6) l'axiome Loc (6 a, b), sans y inclure c.}.

\textbf{N.B.} Je distingue ces conditions comme Loc (6 a), Loc (6 b), Loc (6 c). Je présume qu'on n'aura jamais à faire usage à Loc (6 b), sans supposer également Loc (6 a), mais l'inverse n'est sûrement pas le cas. Ainsi, $W_0=\{f\in\Fl\Cat\mid\pi_0(f)\text{ bijectif}\}$ satisfait a), mais non b). Je présume (sans avoir vérifié) que c'est pareil pour les $W_n$, $n\in\mathbf{N}$, de la “\textit{$n$-équivalence homotopique}”.

Il se pourrait que Loc (6 a) et Loc (6 a, b) soient équivalentes à la forme plus faible, et plus simple, suivante. (Si oui, c'est sous la forme la plus simple, bien sûr, qu'il convient de les énoncer.)

\begin{itemize}[leftmargin=5em,labelwidth=4em,align=left]
\item[\textbf{Loc 6$'$}] (Axiome des carrés cocartésiens, ou de Mayer-Vietoris). Soit
\[
\begin{tikzcd}[row sep=small,column sep=small]
X_0 \arrow[rr,"i_2"] \arrow[dr,"i_1"] & & X_2 \arrow[dr,"j_1"] & \\
& X_1 \arrow[rr,"j_2"] & & X_3
\end{tikzcd}
\]
un diagramme MV dans Cat. Alors
\begin{itemize}[leftmargin=3em,labelwidth=2.5em,align=right]
\item[a)] Si $i_1\in W$, $j_1\in W$.
\item[b)] Si $j_1\in W$, $i_1\in W$.
\end{itemize}
\end{itemize}

\msp{18$'$}
\textbf{Proposition 3.1.}
\[
\begin{aligned}
\text{Loc (6 a)}\quad &\Longrightarrow\quad \text{Loc (6$'$ a),}\\
\text{Loc (6 b)}\quad &\Longrightarrow\quad \text{Loc (6$'$ b),}\\
\text{Loc (6 c)}\quad &\Longrightarrow\quad \text{Loc (6$'$ a, b)}\ \footnotemark[14].
\end{aligned}
\]
\footnotetext[14]{N.B. on a même Loc (6 c) $\Longrightarrow$ $\underbrace{\text{Loc 6 (a, b)}}_{\text{Loc (6)}}$.}

Ça se voit sur le diagramme cubique
\[
\text{(Q)}\qquad\qquad
\begin{tikzcd}[row sep=small,column sep=small]
X_0 \arrow[rr,"i_2"] \arrow[dr,"i_1"] & & X_2 \arrow[dr,"j_1"] \arrow[from=dd,"f_2=\id"' near start] & \\
& X_1 \arrow[rr,"j_2",crossing over] & & X_3 \\
X_0 \arrow[uu,"f_0=\id"] \arrow[rr,"i_2"] \arrow[dr,"\id"] & & X_2 \arrow[dr,"\id"] & \\
& X_0 \arrow[rr,"i_2"] \arrow[uu,crossing over,"f_1=i_1" near end] & & X_2, \arrow[uu,"f_3=j_1"']
\end{tikzcd}
\]
où les $f_0$, $f_1$, $f_2$, $f_3$ de l'énoncé de Loc (6) deviennent $\id_{X_0}$, $i_1$, $\id_{X_2}$, $j_1$. Loc (6 a) donne donc $i_1\in W\Longrightarrow j_1\in W$, \textit{i.e.} Loc (6$'$ a), et Loc (6 c) $j_1\in W\Longrightarrow i_1\in W$, \textit{i.e.} Loc (6$'$ b). Pour Loc (6 b) $\Longrightarrow$ Loc (6$'$ b), on utilise (\textit{cf.} XII, p. 78) le cube

% ---- from der-ch16_p016-020.tex
\noindent\printedpage{14}% head: A. Grothendieck : Dérivateurs, Ch. XVI, Édition des Universités de Montpellier II et Paris VII
% foot: Version 18 mars 2019
\[
\begin{tikzcd}[ampersand replacement=\&, row sep=small, column sep=small]
X_0 \arrow[rr, "i_2"] \arrow[dr, "i_1"] \arrow[dd, "i_1"'] \& \& X_2 \arrow[dr, "j_1"] \arrow[dd, "j_1" near start] \& \\
\& X_1 \arrow[rr, "j_2" near start, crossing over] \& \& X_3 \arrow[dd, "\id"] \\
X_1 \arrow[rr, "j_2" near end] \arrow[dr, "\id"'] \& \& X_3 \arrow[dr, "\id"'] \& \\
\& X_1 \arrow[rr, "j_2"] \arrow[from=2-2, "\id"' near start, crossing over] \& \& X_3,
\end{tikzcd}
\tag*{(Q$'$)} % ?? uncertain: placement of the label "id" printed above the left part of the back lower edge X_1 -> X_3 -- label of the front vertical X_1 -> X_1, or of the edge X_1 -> X_3
\]
où $f_0 = i_1$, $f_1 = \id_{X_1}$, $f_1 = j_1$, $f_3 = \id_{X_3}$, donc % [sic] $f_1 = j_1$ -> $f_2 = j_1$, $f_1$ is already given as $\id_{X_1}$ and $f_2$ is missing
\[
j_1 \in W \Longrightarrow i_1 \in W, \ \ \textit{i.e.}\ \text{Loc (6 b)} \Longrightarrow \text{Loc (6$'$ b)},
\]
et si on a Loc (6 c), on voit que $i_1 \in W \Longrightarrow j_1 \in W$, \textit{i.e.} Loc (6$'$ a).

\textbf{Commentaire 3.2.} Loc (6 a) m'apparaît comme le “minimum vital” pour établir l'existence d'une théorie des cofibres $W$-homotopique dans Cat. Pour cela, il faut en effet disposer d'un critère serviable assurant qu'une flèche $i : A \longrightarrow B$ dans Cat est dans $\Cof_W$.

\msp{page 19}On trouve, quand on dispose de Loc (6 a) (sous forme duale, \textit{i.e.} pour les \textit{fermés,} c'est kif-kif), que les \textit{immersions ouvertes de Dwyer} sont dans $\Cof_W$, ce qui suffit pour avoir \textit{une bonne théorie des cofibres $W$-homotopiques, dans laquelle les carrés MV sont $W$-cocartésiens} (chose a laquelle on tient absolument). % [sic] a laquelle -> à laquelle, missing accent

Quant à Loc (6 b), si je dispose à la place de Loc (6$'$ b) + Loc (6 a), je dois m'en tirer pour prouver Loc (6). Dans XII Appendice (pages 265 ff.) je prouve, paraît-il, que (avec les notations actuelles) Loc (6$'$) implique Loc (6), et même Loc (6 c) si $W$ est stable par facteurs directs (condition Loc (1 bis)). Mais la démonstration est incomplète, puisque j'admets tacitement que les immersions de Dwyer sont dans $\Cof_W$, chose que je n'avais établie que moyennant l'axiome W (7), en fait moyennant Loc (6 a) (ce qui est une partie de W (7)).

\msp{page 20}Mais ici me vient l'idée du lemme de K. S. \textsc{Brown}, cité par \textsc{Baues}, que je vais redonner ici :

\textbf{Lemme 3.2.1.} (K. S. \textsc{Brown}) \textit{Soit $(\mathcal{M}, W)$ un localiseur} ($W$ \textit{modérément saturé}), \textit{muni d'un ensemble $C \subset \Fl\,\mathcal{M}$ satisfaisant les conditions :}
\begin{enumerate}
\item[1)] \textit{$C$ stable par composition et par cochangement de base} (N.B. \textit{on n'a à supposer l'existence des sommes amalgamées}
$\begin{tikzcd}[ampersand replacement=\&, baseline=(\tikzcdmatrixname-1-1.base), row sep=small, column sep=small]
X_0 \arrow[rr, "i_1"] \arrow[dr, "i_2"'] \& \& X_1 \arrow[dr] \& \\
\& X_2 \arrow[rr] \& \& X_3
\end{tikzcd}$
\textit{que si $i_1 \in C$}), \textit{contient les isomorphismes.}
\item[2)] \textit{$\mathcal{M}$ contient un objet initial, $\emptyset_{\mathcal{M}}$, et pour tout objet $X$ de $\mathcal{M}$, $\emptyset_{\mathcal{M}} \longrightarrow X$ est dans $C$.} (N.B. (1) \textit{et} (2) \textit{impliquent que $\mathcal{M}$ est stable par sommes finies.})
\item[3)] \textit{Toute flèche $f$ de $\mathcal{M}$ se factorise en $pi$, avec $i \in C$, $p \in W$.}
\item[4)] \textit{$C \cap W$ est stable par cochangement de base,} i.e. \textit{si}
\[
\begin{tikzcd}[ampersand replacement=\&, row sep=small, column sep=small]
X_0 \arrow[rr, "i_1"] \arrow[dr, "i_2"'] \& \& X_1 \arrow[dr] \& \\
\& X_2 \arrow[rr, "j_1"'] \& \& X_3
\end{tikzcd}
\tag*{$(*)$}
\]
\printedpage{15}% head: A. Grothendieck : Dérivateurs, Ch. XVI, Édition des Universités de Montpellier II et Paris VII
% foot: Version 18 mars 2019
\textit{est cocartésien et $i_1 \in C \cap W$, alors $j_1 \in C \cap W$.} (N.B. \textit{$j_1 \in C$ est déjà contenu dans} (1), \textit{il reste à exiger $j_1 \in W$.})
\end{enumerate}
\textit{Sous ces conditions, on a $C \subset \Cof_W$.}

On peut aussi énoncer le lemme ainsi

\textbf{Corollaire 3.2.2.} \textit{Supposons que $C$ satisfait} (1), (2), (3). \textit{Alors}
\msp{page 21}\textit{les conditions suivantes sur $C$ sont équivalentes.}
\begin{enumerate}
\item[a)] \textit{Condition} (4) \textit{ci-dessus,} i.e. \textit{stabilité de $C \cap W$ par cochangement de base.}
\item[a$'$)] \textit{$C \cap W \subset W^{\univ}$.}
\item[b)] \textit{$C \subset \Cof_W$.}
\item[b$'$)] \textit{Pour tout carré cocartésien}
\[
\begin{tikzcd}[ampersand replacement=\&, row sep=small, column sep=small]
X_0 \arrow[rr, "i"] \arrow[dr, "p"'] \arrow[to=2-4, phantom, "\textit{coc.}" description] \& \& X_1 \arrow[dr, "q"] \& \\
\& X'_0 \arrow[rr, "j"] \& \& X'_1
\end{tikzcd}
\]
\textit{avec $i \in C$, $p \in W$, on a $q \in W$.}
\end{enumerate}
\textsc{Démonstration}. On a
\[
\begin{aligned}
&\text{b} \iff \text{b}'\\
&\text{a} \iff \text{a}'
\end{aligned}
\]
car $C$ est stable par cochangement de base. On a b $\Longrightarrow$ a$'$ car $\Cof_W \cap W = W^{\univ}$ grâce à l'hypothèse 3°). On a a $\Longrightarrow$ b$'$ par le lemme de Brown.

On va appliquer ceci au cas $\mathcal{M} = \mathrm{Cat}$, $C =$ immersions ouvertes de Dwyer. Les conditions 1), 2) ne dépendent pas de $W$, elles sont vérifiées, et 3) aussi dès que $W$ est un localiseur fondamental. Quant à 4), c'est une conséquence de l'ancien W (7 bis), \textit{i.e.} Loc (6$'$) (on a
\msp{page 22}besoin apparemment de Loc (6$'$ a \textit{et} b).) On en conclut $C \subset \Cof_W$ - c'est merveilleux! (Mais je ne sais pas prouver le lemme \dots)

Cela légitime la démonstration faite dans XII \textit{loc. cit.}, et on trouve la

\textbf{Proposition 3.3.} \textit{La condition} Loc (6$'$) (\textit{axiome du carré cocartésien, forme forte} (a) \textit{et} (b)) équivaut \textit{à} Loc (6 a, b) (\textit{axiome des homomorphismes des carrés cocartésiens}) \textit{et est} impliqué \textit{par} Loc (6 c). \textit{Si $W$ satisfait à} Loc (1 bis), i.e. \textit{est stable par facteurs directs, alors} Loc (6$'$) équivaut \textit{à} Loc (6 c) \textit{et à} Loc (6 a b c).

\textbf{Commentaire 3.4.} Ainsi, on trouve une théorie des cofibres $W$-homotopiques \textit{soit} sous l'hypothèse Loc (6$'$) (équivalente à Loc (6) plus compliquée), soit sous l'hypothèse plus faible (et plus compliquée) Loc (6 a).

% left margin: a vertical line is printed alongside the following paragraph
Dans le premier cas, c'est une théorie de nature très particulière, puisque les carrés $W$-cocartésiens sont “réversibles” (\textit{i.e.} on a non seulement $i_1 \in W \Longrightarrow j_1 \in W$, mais aussi l'implication inverse).

Il reste donc la question si Loc (6$'$ a) suffit pour une théorie des cofibres $W$-homotopiques, ou ce qui revient sûrement au même, si ça implique déjà Loc (6 a) (donc lui est équivalent). C'est entièrement suspendu à la \textit{question} :

\printedpage{16}% head: A. Grothendieck : Dérivateurs, Ch. XVI, Édition des Universités de Montpellier II et Paris VII
% foot: Version 18 mars 2019
\msp{page 23}% left margin: a question mark followed by a double vertical line alongside the first three lines of the following paragraph
\marginpar{\footnotesize ?\,$\Vert$}l'axiome Loc (6$'$ a) (axiome du carré cocartésien, forme directe) implique-t-il que l'ensemble des immersions ouvertes de Dwyer qui sont dans $W$ (immersions de Dwyer $W$-triviales) est stable par cochangement de base ? En fait, on voit qu'il suffirait pour cela de vérifier l'axiome du cube (ou des homomorphismes de carrés cocartésiens), forme directe Loc (6 a), dans le cas où $f_0$ et $f_2$ sont des identités, donc à prouver qu'alors $f_1 \in W \Longrightarrow f_3 \in W$ (ce qui donnerait directement $\mathrm{Dw} \subset \Cof_W$). Mais je ne vois aucun moyen de le prouver \dots

Supposons maintenant que la condition Loc (6 a) est vérifiée, donc qu'on dispose d'une “bonne théorie des cofibres $W$-homotopiques”, donc aussi d'une théorie des carrés $W$-cocartésiens (\textit{cf.} XII, p. 258-262). Il y a deux questions qui se posent, concernant la validité, avec ces hypothèses plus faibles, des développements de XII :
\begin{enumerate}
\item[1°)] Soit un carré commutatif
\[
\begin{tikzcd}[ampersand replacement=\&, row sep=small, column sep=small]
X_0 \arrow[rr] \arrow[dr] \& \& X_2 \arrow[dr] \& \\
\& X_1 \arrow[rr] \& \& X_3
\end{tikzcd}
\]
dans Cat, est-il vrai que ce carré est $W$-cocartésien si et seulement si la flèche
\[
\text{canonique}\ \displaystyle\int
\begin{tikzcd}[ampersand replacement=\&, row sep=small, column sep=small]
\& X_1 \\
X_0 \arrow[ur] \arrow[dr] \& \\
\& X_2
\end{tikzcd}
\quad \longrightarrow \ X_3 \text{ est dans } W \text{ ? }\text{\footnotemark[15]}
\]
\footnotetext[15]{Mieux, dans tous les cas, quand on a un diagramme
$\begin{tikzcd}[ampersand replacement=\&, baseline=(\tikzcdmatrixname-2-1.base), row sep=tiny, column sep=small]
\& X_1 \\
X_0 \arrow[ur] \arrow[dr] \& \\
\& X_2
\end{tikzcd}$, son intégrale $\displaystyle\int$ est une somme amalgamée $W$-homotopique dudit diagramme (moyennant Loc (6 a) sans plus, et même seulement Loc (6$'$ a)).}
\end{enumerate}
\msp{page 24}
\begin{enumerate}
\item[2°)] Est-il encore vrai que si $A$ est dans Cat, et $u : F \lhook\joinrel\longrightarrow G$ un homomorphisme dans $A^{\wedge}$, alors $u \in \Cof_{W_A}$, et si $u \in W$, $u \in W_A^{\univ}$, donc a-t-on
\[
\Mono_{A^{\wedge}} \subset \Cof_{W_A}, \qquad \Mono_{A^{\wedge}} \cap W_A \subset W_A^{\univ} \text{ ? }\text{\footnotemark[16]}
\]
\footnotetext[16]{N.B. Loc (6$'$ a) suffit au lieu de ( Loc (6 a)).}
\end{enumerate}
Mais 1°) est prouvé dans XII, prop. 20, p. 262 - je viens de vérifier que la démonstration est valable dans les hypothèses actuelles \footnote[17]{Mais elle utilise de plus Loc (3 bis). Mais avec une autre démonstration plus simple que je viens de rajouter (p 264), on peut s'en passer (et on obtient une vison plus complète de la situation).}, pourvu que je vérifie que le corollaire 3, p. 149 dans XII est valable, \textit{i.e.} % [sic] (footnote 17) vison -> vision, missing letter
\begin{enumerate}
\item[1° bis)] Si
$\begin{tikzcd}[ampersand replacement=\&, baseline=(\tikzcdmatrixname-1-1.base)]
X_0 \arrow[r] \arrow[d] \& X_1 \arrow[d] \\
X_2 \arrow[r] \& X
\end{tikzcd}$
est un carré MV dans Cat, alors
$\displaystyle\int
\begin{tikzcd}[ampersand replacement=\&, row sep=small, column sep=small]
\& X_1 \\
X_0 \arrow[ur] \arrow[dr] \& \\
\& X_2
\end{tikzcd}
\quad \longrightarrow \quad X$ est dans $W$ (donc, moyennant 1°), le carré est $W$-cocartésien).
\end{enumerate}

\printedpage{17}% head: A. Grothendieck : Dérivateurs, Ch. XVI, Édition des Universités de Montpellier II et Paris VII
% foot: Version 18 mars 2019
Pour le vérifier, je suis amené à vérifier que le théorème 4 dans loc. cit. (p. 146) est encore valable.

Je viens de le faire, et ai annoté en conséquence le théorème 4 et ses corollaires (certains résultats (cor. 4 et cor. 5) utilisant W (7 bis) dans toute sa force, \textit{i.e.} Loc (6$'$), \textit{i.e.} Loc (6 a \textit{et} b). Mais le corollaire 3 du théorème 4 est O.K.

Je viens aussi de vérifier la validité de 2°), en vérifiant la démonstration du corollaire du théorème 5 (XII, p. 168) et en l'annotant en conséquence.

\msp{page 25}3.5. Résumons : on a
\[
\text{Loc (6 c)} \ \Longrightarrow \ \text{Loc (6)} \ \iff \ \text{Loc (6$'$)}
\tag*{(3.5.1)}
\]
et
\[
\text{Si } \ W = W^{\natural}\,, \ \text{ on a } \ \text{Loc (6 c)} \ \iff \ \text{Loc (6)} \ \iff \ \text{Loc (6$'$)}.
\tag*{(3.5.2)}
\]
D'autre part
\[
\text{Loc (6 a)} \ \Longrightarrow \ \text{Loc (6$'$ a)}\,,
\tag*{(3.5.3)}
\]
et j'ignore si l'implication inverse
\[
\text{Loc (6$'$ a)} \ \overset{?}{\Longrightarrow} \ \text{Loc (6 a)}
\]
est valable \footnote[18]{C'est le cas en présence de Loc (8, 9).}. Pour établir Loc (6 a), il suffit de le faire en tous cas quand $f_0$ et $f_2$ sont des isomorphismes.

L'axiome Loc (6 a) semble le “minimum vital” pour que $W$ donne lieu à une “bonne théorie des cofibres $W$-homotopiques”, donc aussi des carrés $W$-cocartésiens. De façon plus précise, j'ai envie de prouver ceci :

\textbf{Proposition 3.6.} \textit{L'axiome} Loc (6 a) \textit{équivaut à la conjonction des deux propriétés suivantes.}

\printedpage{18}% head: A. Grothendieck : Dérivateurs, Ch. XVI, Édition des Universités de Montpellier II et Paris VII
% foot: Version 18 mars 2019
\begin{itemize}[leftmargin=7em, labelwidth=6em]
\item[\textbf{Loc (6 A)}] Cat a “assez de $W$-cofibrations”,
\msp{page 26}\textit{i.e.} toute flèche $f$ de Cat se factorise en $f = pi$, avec $i \in \Cof_W$ et $p \in W$. (Cela signifie aussi que Cat muni de $(W, \Cof_W)$ est une catégorie à cofibrations de K. S. \textsc{Brown} - \textsc{Anderson}, et tous les résultats de la théorie de \textsc{Brown-Quillen} sont applicables. En particulier, on a une bonne théorie des carrés $W$-cocartésiens.)
\item[\textbf{Loc (6 B)}] Si $X$ dans Cat est réunion de deux ouverts $U$, $V$, d'intersection $S$, alors le carré
\[
\begin{tikzcd}[ampersand replacement=\&, row sep=large, column sep=large]
S \arrow[r] \arrow[d] \& U \arrow[d] \\
V \arrow[r] \& X
\end{tikzcd}
\]
est $W$-cocartésien \footnote[19]{En fait, Loc (6 B) est conséquence de Loc (6 A), laquelle donc est \textit{équivalente} à Loc (6 a). (\textit{Cf.} esquisse de démonstration dans 4.4.15, pages 66-68.)}.
\end{itemize}
D'autre part, la condition Loc (6 A) peut être remplacée par la condition (en apparence plus faible)
\begin{itemize}[leftmargin=7em, labelwidth=6em]
\item[\textbf{Loc (6 A$'$)}] Tout morphisme de Dwyer est dans $\Cof_W$.
\end{itemize}
\textbf{Corollaire 3.7.} \textit{La condition} Loc (6), i.e. Loc (6 a, b) (\textit{ou encore la condition équivalente} Loc (6$'$)) \textit{équivaut à la conjonction des conditions} Loc (6 A) (\textit{ou} Loc (6 A$'$)) \textit{et} Loc (6 B) \textit{ci-dessus}), \textit{et de la condition}
\begin{itemize}[leftmargin=7em, labelwidth=6em]
\item[\textbf{Loc (6 C)}] Pour tout carré $W$-cocartésien
\[
\begin{tikzcd}[ampersand replacement=\&, row sep=large]
X_0 \arrow[r, "i"] \arrow[d] \& X_1 \arrow[d] \\
X_2 \arrow[r, "j"] \& X
\end{tikzcd}
\]
dans Cat, $j \in W \Longrightarrow i \in W$.
\end{itemize}
\msp{page 27}D'autre part, je soupçonne que Loc (6 c) n'est guère vérifiée que si $W = W^{\natural}$ (cas qui est le plus commun bien sûr), auquel cas Loc (6 c) \textit{équivaut} à Loc (6). Il n'y a sans doute pas lieu de l'inclure dans Loc (6), pas plus que $W = W^{\natural}$ dans Loc (1).

% ---- from der-ch16_p021-025.tex
\printedpage{19}
% head: A. Grothendieck : Dérivateurs, Ch. XVI, Édition des Universités de Montpellier II et Paris VII

\section*{4\quad Axiome Loc (7) des $W$-fibrations}

4.1. Rappelons la notion de \textit{$W$-fibration faible} (XII p. 51, 52) : la version la plus forte des 8 conditions équivalentes pour que $f : X \longrightarrow S$ soit une fibration faible est celle-ci.

\begin{itemize}[leftmargin=2cm]
\item[\textbf{(Ff)}] $\forall\ S'' \longrightarrow S'$ dans $\Cat/S$, avec $S'$, $S''$ \textit{contractiles,} le morphisme correspondant $X'' = X \times_S S'' \longrightarrow X' = X \times_S S'$ est dans $W$.
\end{itemize}

Quand $W$ est fortement saturé, alors ceci implique la condition, en apparence plus forte :

\begin{itemize}[leftmargin=2cm]
\item[\textbf{(Ff$'$)}] $\forall\ S'' \longrightarrow S'$ dans $\Cat/S$, telle que $S'' \longrightarrow S'$ soit un homotopisme, $X'' \longrightarrow X'$ est dans $W$.
\end{itemize}

Je désigne par $\Fibf_W$ l'ensemble des fibrations faibles. Je résume dans un diagramme d'inclusions les principales classes de flèches dans Cat que nous avons rencontrées, s'apparentant à des classes de “fibrations” (pour $W$, dans un sens convenable) :

\msp{page 28}

\[
\begin{tikzcd}
 & & \Fib_W \arrow[d] & & \\
\Parf_\omega \arrow[r] & \Parf_W \arrow[r] & \Fibf_W & \underbrace{\Catfib^W}_{\mathcal{F}_W} \arrow[l] & \underbrace{\Catfib^\omega}_{\mathcal{F}_\omega} \arrow[l] \ ,
\end{tikzcd}
\]
où $\Fib_W$ désigne les $W$-fibrations ($f : S' \longrightarrow S$ est dans $\Fib_W$ si et seulement si le foncteur changement de base $X \longmapsto X' = X \times_S S'$ de $\Cat/S$ dans $\Cat/S'$ est compatible avec les localiseurs $W_S$ et $W_{S'}$), Catfib les Cat-fibrations, $\Catfib^W$ (resp. $\Catfib^\omega$) les Cat-fibrations telles que les foncteurs changement de base soient dans $W$ (resp. soient dans $W_\omega$, \textit{cf.} déf. plus bas) \footnote[20]{N.B. $\Catfib_0^W = \Catfib \cap \Fibf_W$, $\Catfib_0^\omega = \Catfib \cap \Fibf_{W_\omega}$}, $\Parf_W$ les morphismes $W$-parfaits, \textit{i.e.} $W$-lisses et $W$-propres, enfin $\Parf_\omega$ les morphismes $W_\omega$-parfaits, où $W_\omega$ est le plus petit localiseur satisfaisant Loc (1, 2, 3). Parmi les espèces de “fibrations” envisagées, c'est cette dernière seulement, avec $\Catfib^\omega$, qui est indépendante de $W$.

À ce titre, elle me paraît jouer un rôle similaire, dual, à celui des morphismes de Dwyer, et je présume qu'on peut leur faire jouer un rôle similaire. En tous \msp{page 29} cas, on sait que toute flèche $f$ de Cat se factorise en $pi$, avec $i \in W_\omega$ et $p \in \Parf_\omega$, donc $i \in W$ pour tout $W$ satisfaisant Loc (1, 2, 3). On sait aussi que $\Catfib^\omega$ satisfait l'énoncé de factorisation forte, $f$ se factorise en $f = pi$, $i \in W_\omega$, $p \in \Catfib_0^\omega$. \footnote[21]{Donc la propriété de factorisation est valable pour toutes ces classes, sauf peut-être $\Fib_W$.}

Notons aussi que $\Parf_\omega$ et $\Parf_W$, ainsi que $\Fib_W$, sont stables non seulement par changement de base (c'est le cas des 5 classes envisagées, qui de plus contiennent tous les isomorphismes), mais aussi par composition. (De plus, $\Fib_W$ est stable par facteurs directs, et je présume qu'il en est de même de $\mathrm{Liss}_W$ et de $\mathrm{Prop}_W$, donc de $\Parf_W$, en particulier de $\Parf_{W_\omega}$, mais je n'ai pas pris la peine de le vérifier.) Par contre, a priori $\Fibf_W$ n'est pas en général stable par composition, et $\Catfib^W$, $\Catfib^\omega$ n'ont pas de raison
% foot: Version 18 mars 2019
\printedpage{20}
% head: A. Grothendieck : Dérivateurs, Ch. XVI, Édition des Universités de Montpellier II et Paris VII
non plus d'être stables par composition \footnote[22]{il faudrait expliciter par des exemples.}.

4.2. Finalement, il paraît plus “éclairant” d'introduire, au lieu du localiseur fondamental $W_\omega$ bien déterminée, un localiseur fondamental % [sic] déterminée -> déterminé, accord avec « localiseur » (masculin)
\[
W_0 \subset \mathrm{Fl}(\Cat),
\]
qui jouera le rôle d'une sorte de “localiseur de référence” pour formuler certaines propriétés de $W$. On supposera toujours que
\[
W \supset W_0.
\]

\msp{page 30}

Suivant les cas, $W_0$ pourra être $W_\omega$, ou quelque autre localiseur fondamental défini comme le “plus petit” satisfaisant à des axiomes de stabilité donnés. Ou ça pourra être $\mathbf{W}_\infty$ (les équivalences faibles habituelles), quand on s'intéresse aux localiseurs qui contiennent $\mathbf{W}_\infty$ (les seuls, apparemment, avec lesquels on ait travaillé jusqu'à présent en homotopie). Ça pourra être aussi $W$, et les énoncés obtenus en spécialisant les résultats obtenus au cas où $W = W_0$ ne sont nullement tautologiques. Le diagramme de la page 28 se remplace alors par le diagramme d'inclusions suivant, plus éclairant pour notre propos :
\[
\begin{tikzcd}
 & \Fib_W \arrow[d] & \\
\Parf_W \arrow[r] & \Fibf_W & \mathcal{F}_W \arrow[l] \\
\Parf_{W_0} \arrow[r] \arrow[u] & \Fibf_{W_0} \arrow[u] & \mathcal{F}_{W_0} \arrow[l] \arrow[u] \\
 & \Fib_{W_0}, \arrow[u] &
\end{tikzcd}
\tag*{(4.2.1)}
\]
où pour tout localiseur fondamental $W$, je pose pour abréger
\[
\mathcal{F}_W = \Catfib^W,
\]
ce qui définit donc également $\mathcal{F}_{W_0}$. La classe $\Fib_W$ \msp{page 31} doit être regardée comme le point de mire de notre attention ; la plupart des propriétés de $W$ que nous allons passer en revue ici, s'exprimeront par une inclusion de la forme
\[
\mathcal{F} \subset \Fib_W,
\]
où $\mathcal{F}$ est une des sept autres classes du diagramme (4.2.1) ; à l'exception tout au plus du cas $\mathcal{F} = \Fib_{W_0}$, \textit{i.e.} de l'inclusion $\Fib_{W_0} \subset \Fib_W$. Comme propriété sur $W$, cette inclusion paraît très faible et pour cela pas utile à grand chose, sauf bien sûr dans le cas où $\Fib_{W_0} \longrightarrow \Fibf_{W_0}$ est une égalité, de sorte que l'inclusion envisagée équivaut à $\Fibf_{W_0} \subset \Fib_W$ (et implique que $\Parf_{W_0}$ et $\mathcal{F}_{W_0}$ sont eux aussi contenus dans $\Fib_W$). C'est là par contre, comme nous allons voir, une propriété assez forte et fort utile de $W$, impliquant notamment le fait que $(\Cat, W, \Fib_W)$ est une catégorie à fibrations.
% foot: Version 18 mars 2019

\printedpage{21}
% head: A. Grothendieck : Dérivateurs, Ch. XVI, Édition des Universités de Montpellier II et Paris VII

\msp{page 32}

Aux fins de référence ultérieure, je vais formuler maintenant les divers aspects principaux de \textit{l'axiome des $W$-fibrations} Loc (7). La formulation à laquelle j'étais parvenu (\textit{cf.} I, VI, XII) utilisait la notion que j'appelle à présent “$W$-fibrations \textit{faible}” (l'ex $W$-fibration, désormais détrônée !), et sous sa forme la plus forte s'écrivait comme l'inclusion % [sic] W-fibrations faible -> W-fibration faible, accord du nombre
\[
\Fibf_W \subset \Fib_W.
\]
Cette notion de $W$-fibration faible, tout comme celle de morphisme $W$-parfait, si utiles soient-elles, m'apparaît cependant un peu technique, et je préfère donner une formulation qui ne fasse appel qu'à des notions plus familières et d'apparence plus élémentaires. C'est pourquoi je vais remplacer $\Fibf_W$ par la sous-classe $\mathcal{F}_W = \Catfib^W$, ce qui (on le verra) revient en fait au même.

\msp{page 33}

Cette classe $\mathcal{F}_W$ est extrêmement naturelle, et s'impose même dans le yoga des dérivateurs. En effet, dans la construction de $\HOT_W(S)$ pour $S$ dans Cat, on passe à une catégorie de fractions de la catégorie
\[
\Cat(S) = \underline{\mathrm{Hom}}(S, \Cat),
\]
par un localiseur $W(S)$ (que nous avions généralement noté $W_S^{\mathrm{g}}$ ou $W_S^{\mathrm{u}}$ \ldots). Cette catégorie $\Cat(S)$ est d'autre part isomorphe à la catégorie
\[
\Fibsc\, S
\]
des catégories fibrées et scindées sur $S$, avec comme flèches les $S$-foncteurs qui respectent le scindage (et en particulier cartésiens) (\textit{cf.} VI, et notamment le diagramme p. 175). On a vu (\textit{loc. cit.}) que la catégorie localisée $\HOT_W(S)$ est équivalente à la catégorie localisée de la catégorie très voisine
\[
\Fib\, S \qquad \text{(que je préférerais maintenant noter $\Catfib\, S$ ou $\Catfib_S$),}
\]
formé des catégories \textit{fibrées} sur $S$, avec comme \msp{page 34} flèches les seuls foncteurs cartésiens. (On fera attention cependant que le foncteur canonique $\Fibsc\, S \longrightarrow \Fib\, S$ est fidèle mais non pleinement, néanmoins il induit une équivalence sur les catégories localisées, en définissant de façon évidente, par la $W$-équivalence fibre par fibre, \textit{i.e.} comme $W_S^{\mathrm{u}}$, le localiseur pertinent sur $\Fib\, S$.) Donc pour cette raison de principe, et d'ailleurs tout au long du développement de la théorie de $\HOT_W$, la catégorie $\Fib\, S$ et par là, les foncteurs fibrants, ont joué un rôle essentiel. La classe de flèches que nous notons ici $\mathcal{F}_W$ ou $\Catfib^W$ (on pourrait aussi la noter $\Fib^W$, si ce n'était pour le risque de confusion avec $\Fib_W$) s'interprète alors très naturellement à partir de là. Une flèche $f : X \longrightarrow S$ est dans $\mathcal{F}_W$ si et seulement si elle est a) fibrante, et par là définit de façon particulièrement directe un objet de $\HOT_W(S)$, puisque à $S$-équivalence près, $X/S$ provient d'un objet de $\Fibsc\, S$ ou encore de $\Cat(S)$, \msp{page 35} et b) l'objet dans $\HOT_W(S)$ défini par $f$ est dans $\HOT_W^{\mathrm{lc}}(S)$, \textit{i.e.} dans l'image essentielle de la sous-catégorie strictement pleine $\Fibsc_0^W S$ de $\Fibsc\, S$ (par le foncteur canonique
% foot: Version 18 mars 2019
\printedpage{22}
% head: A. Grothendieck : Dérivateurs, Ch. XVI, Édition des Universités de Montpellier II et Paris VII
$\Fibsc \longrightarrow \HOT_W(S)$), ou ce qui revient au même (du moins si $W$ est fortement saturée) que le foncteur % [sic] saturée -> saturé, accord avec W (« W fortement saturé » ailleurs)
\[
\begin{aligned}
S^{\circ} \quad &\longrightarrow \quad \mathrm{Hot}_W \\
s \quad &\longmapsto \quad \mathrm{hot}_W(X_s)
\end{aligned}
\]
soit un foncteur “localement constant”. Ainsi, on considérera les $f \in \mathcal{F}_W$ comme les flèches les plus économiques et les plus naturelles dans Cat, qui définissent des $W$-types d'homotopie relative localement constants.

Notons aussi que toutes les classes de flèches (4.2.1) sont stables par changement de base (et elles contiennent les isomorphismes). Les catégories $\Fib_W$, $\Parf_W$ et leurs variantes $W_0$ sont aussi stables par composition, par contre, ce n'est pas le cas pour $\Fibf_W$ et $\mathcal{F}_W$ et leurs variantes $W_0$, sauf axiomes particuliers sur $W$ ou sur $W_0$). Il en résulte aussitôt \msp{page 36} que si $\mathcal{F}$ est une de ces classes, la condition % [sic] « sur W_0). » -> parenthèse fermante sans parenthèse ouvrante correspondante
\[
\mathcal{F} \subset \Fib_W
\]
équivaut à celle ci :

\begin{itemize}[label={},leftmargin=3.5cm]
\item Pour tout diagramme cartésien
\[
\begin{tikzcd}
X \arrow[d, "f"'] & X' \arrow[l, "q"'] \arrow[d, "f'"] \\
S & S', \arrow[l, "p"]
\end{tikzcd}
\tag*{(4.2.2)}
\]
avec $f \in \mathcal{F}$ et $p \in W$, on a $q \in W$.

(N.B. On a aussi $f' \in \mathcal{F}$, par la stabilité rappelée plus haut.)
\end{itemize}

4.3. Voici donc (\textit{cf.} plus bas) l'axiome pertinent, Loc (7) avec ses trois aspects a, b, c \footnote[23]{J'y ai rajouté finalement un quatrième F (comme “fibration”).}. On a vu (et on reverra) que l'on a les implications
\[
\text{a} \qquad \Longrightarrow \qquad \text{b} \qquad \Longrightarrow \qquad \text{c} \qquad (\Longrightarrow \qquad \text{F}).
\]
Quand on référera à Loc (7) sans autre précision, il sera entendu qu'il s'agit de la forme la plus forte Loc (7 a), qui implique donc les deux autres. Il n'en faut pas moins pour être totalement à l'aise. Mais je rappelle que la variante affaiblie Loc (7b) est suffisante pour faire la théorie des foncteurs $f_!$ pour les catégories $\HOT_W(S)$ \footnote[24]{Ici ma mémoire me joue un tour - on n'a nullement besoin pour cela de l'axiome déjà très fort Loc (7b), mais seulement de Loc (3 bis), dont il semble bien qu'il soit vérifiée pour tous les localiseurs fondamentaux “raisonnables”. % [sic] qu'il soit vérifiée -> qu'il soit vérifié, accord avec « il »
} (\textit{cf.} VI).

\msp{page 37}

C'est d'ailleurs une chose surprenante, puisque via le yoga général des dérivateurs associés aux catégories de modèles, on s'attendrait plutôt que pour avoir des adjoints à gauche $f_!$ des foncteurs $f^*$ d'image inverse, il faut faire sur $W$ des hypothèses du type \textit{$W$-cofibration}
% foot: Version 18 mars 2019
\printedpage{23}
% head: A. Grothendieck : Dérivateurs, Ch. XVI, Édition des Universités de Montpellier II et Paris VII
(essentiellement, Loc (6 a)). Il me semble peu probable que Loc (7 b), ni même Loc (7), implique Loc (6 a), ni même sa variante plus faible Loc (6$'$ a) (axiome du carré cocartésien, ou de Mayer-Vietoris, forme faible). Si d'autre part ces axiomes ne sont pas \textit{vérifiés,} alors que Loc (7) l'est, il y a des chances que cela donnera un exemple où il y a un formalisme des $f_!$, sans que l'axiome d'exactitude à gauche soit vérifié.

Quant à la forme Loc (7 c), plus faible encore, elle suffit pourtant (comme nous le verrons plus bas) à assurer que Cat admet suffisamment de $W$-fibrations \msp{page 38} (résultat qui n'a rien d'évident, et ne pouvait m'apparaître qu'après avoir pris connaissance du substantiel lemme de K. S. \textsc{Brown}, dans son beau travail où il introduit les catégories à fibrations et à cofibrations).

Nous donnerons aussi plus bas des cas fort simples, où la condition Loc (7 b) entraîne déjà Loc (7) ($W$ fortement saturé, \textit{i.e.} Loc (1 bis)), voire même où Loc (7 c) entraîne Loc (7) (cas où $W$ satisfait à l'axiome Loc (8) des limites inductives). Je signale en passant que je ne connais pas de localiseur fondamental satisfaisant l'axiome Loc (8) (qui implique déjà $W = W^{\natural}$, \textit{i.e.} Loc (1 bis)), sans que $W$ soit fortement saturé (\textit{i.e.} en satisfasse même Loc (1 ter)).

\begin{itemize}[label={},leftmargin=0.6cm]
\item \textbf{Loc (7)}
\begin{itemize}[leftmargin=1.2cm]
\item[a)] On a $\mathcal{F}_W \subset \Fib_W$, \textit{i.e.} pour tout diagramme cartésien (4.2.2) dans Cat, avec $f \in \mathcal{F}_W$ (\textit{i.e.} $f$ Cat-fibrant et faiblement $W$-fibrant) et $p \in W$, on a $q \in W$.
\item[b)] Pour tout diagramme commutatif
\[
\begin{tikzcd}[column sep=small]
X \arrow[rr, "f"] \arrow[dr, "p"'] & & Y \arrow[dl, "q"] \\
 & S &
\end{tikzcd}
\]
\msp{page 39}
avec $p, q \in \mathcal{F}_W$ et $f \in W$, on a $f \in W_S^{\mathrm{u}}$ (ou, ce qui revient au même comme $X$, $Y$ sont $W$-lisses sur $S$, $f_s \in W$ pour tout $s \in \mathrm{Ob}\, S$) \footnote[25]{\textbf{N.B.} La démonstration de $B(\mathcal{F}_W) \Longrightarrow B(\Fibf_W)$ montre en fait qu'il suffit de postuler b) lorsque $f$ est cartésien, et même lorsque $X$, $Y$ sont scindées sur $S$ et $f$ compatible avec le scindage.}.
\item[c)] On a
\[
\mathcal{F}_W \cap W \subset W^{\mathrm{u}},
\]
en d'autres termes, pour toute flèche $f : X \longrightarrow S$ dans Cat telle que $f \in \mathcal{F}_W$ (on pourrait appeler ces flèches des Cat-$W$-fibrations, étant entendu qu'elles sont aussi des Cat-fibrations, mais non nécessairement des $W$-fibrations), si $f \in W$, alors $f \in W_S^{\mathrm{u}}$ (ou, ce qui revient au même, les fibres $X_s$ de $f$ ($s \in \mathrm{Ob}\, S$) sont $W$-asphériques).
\end{itemize}
\end{itemize}
% foot: Version 18 mars 2019

% ---- from der-ch16_p026-030.tex
\printedpage{24}
% head: A. Grothendieck : Dérivateurs, Ch. XVI, Édition des Universités de Montpellier II et Paris VII
% foot: Version 18 mars 2019

\begin{itemize}[leftmargin=4em]
\item[F)] \footnote[26]{rajouté après coup.} Cat a suffisamment de $W$-fibrations, \textit{i.e.} Cat muni de $(W, \Fib_W)$ est une catégorie à cofibrations (propre, et même de Brown, \textit{i.e.} $X \longrightarrow e$ est dans $\Fib_W$ pour tout $X$ dans Cat. % [sic] (propre, et même de Brown, -> parenthesis never closed, closing ")" apparently missing
\end{itemize}

Il est clair que c est le cas particulier de b obtenu en prenant $Y = S$, $q = \mathrm{id}_S$, donc b $\Longrightarrow$ c. L'implication a $\Longrightarrow$ b est moins évidente, elle est établie d'abord dans XII, et nous y reviendrons ci-dessous (4.4.2 ci-dessous). Le reste du présent paragraphe est consacré à étudier diverses variantes des conditions a) b) c) précédentes et les implications entre celles-ci.

\msp{page 40}

4.4. Soit, plus généralement,
\[
\cF \subset \Fl(\Cat)
\]
un ensemble de flèches dans Cat, en pensant plus particulièrement au cas où $\cF$ est un des sept ensembles du diagramme (4.2.1) (p. 30) autres que $\Fib_W$. Nous pouvons formuler trois propriétés similaires pour $\cF$ à celles formulées dans Loc (7 a b c) dans le cas de $\cF_W$, notons ces propriétés
\[
\mathrm{A}(\cF)\,,\ \mathrm{B}(\cF)\,,\ \mathrm{C}(\cF)\,,\ \text{\footnotemark[27]}
\]
\footnotetext[27]{\textbf{N.B.} Ces conditions ont un sens si $\cF \subset \Fl(\cM)$, si $(\cM, W)$ est une catégorie de modèles quelconque.}
donc
\[
\left\{
\begin{array}{lcl}
\mathrm{A}(\cF) & : & \cF \subset \Fib_W. \\[1ex]
\mathrm{B}(\cF) & : & \Fl(\cF_S) \cap W_S \subset W_S^{\mathrm{u}} \\
 & & \text{(où $\cF_S$ désigne la sous-catégorie pleine de $\Cat/S$} \\
 & & \text{formée des $f : X \longrightarrow S$ avec $f \in \cF$).} \\[1ex]
\mathrm{C}(\cF) & : & \cF \cap W \subset W^{\mathrm{u}}.
\end{array}
\right.
\tag*{(4.4.1)}
\]
Il est clair qu'on a $\mathrm{B}(\cF) \Longrightarrow \mathrm{C}(\cF)$, et je vais donner une condition sur $\cF$ qui implique qu'on a aussi $\mathrm{A}(\cF) \Longrightarrow \mathrm{B}(\cF)$.

\msp{page 41}

(Nous l'appliquerons dans les cas suivants : % [sic] (Nous -> opening parenthesis never closed in this paragraph, closing ")" apparently missing
\[
\cF \;=\; \cF_W,\ \Fibf_W,\ \Parf_W,\ \cF_{W_0},\ \Fibf_{W_0}.
\tag*{(4.4.2)}
\]
Je signale d'ailleurs que l'on verra plus bas (cor. 4.4.10) que
\[
\mathrm{A}(\cF_W) \iff \mathrm{A}(\Fibf_W), \qquad \mathrm{A}(\cF_{W_0}) \iff \mathrm{A}(\Fibf_{W_0}),
\tag*{(4.4.3)}
\]
et de même pour les variantes B, C, de sorte que pour l'implication que nous voulons établir, il n'y a pas lieu de distinguer les cas de $\cF_W$ et $\Fibf_W$, de $\cF_{W_0}$ et $\Fibf_{W_0}$. D'ailleurs il suffira de regarder les ensembles $\cF_{W_0}$, $\Fibf_{W_0}$, $\Parf_{W_0}$, puisque le cas des ensembles relatifs à $W$ en sont un cas particulier (en prenant $W_0 = W$). \footnote[28]{Cela fera donc a priori $3 \times 6 = 18$ conditions différentes, mais à cause de (4.4.3) il n'y en a (à équivalence près) que $4 \times 3 = 12$, et en fait deux parmi ces quatre triples de conditions se réduisent à une seule (pour $\cF = \cF_W$ et $\cF = \cF_{W_0}$), de sorte qu'il ne reste au total que \textit{huit} conditions différentes, que nous résumons dans le diagramme récapitulatif (4.4.7.1) de la page 48. Il y en a même seulement \textit{quatre,} $\mathrm{A}_{W_0}$, $\mathrm{B}_{W_0}$, $\mathrm{C}_{W_0}$, $\mathrm{D}_{W_0}$, comme fonctions de $W_0$, les quatre autres A, B, C, D s'obtiennent en spécialisant au cas où $W_0 = W$.}

\printedpage{25}
% head: A. Grothendieck : Dérivateurs, Ch. XVI, Édition des Universités de Montpellier II et Paris VII
% foot: Version 18 mars 2019

\textbf{Lemme 4.4.1.} \footnote[29]{Vaut dans une catégorie de modèles quelconque.} \textit{Supposons que $\cF$ satisfasse la condition suivante : toute flèche $f$ de} Cat \textit{se factorise en $f = pi$, avec $i \in W$, $f \in \cF$. Alors la condition} $\mathrm{A}(\cF)$ \textit{sur $\cF$ implique que dans} Cat \textit{il y a “assez de $W$-fibrations”, et par suite} (\textit{par un argument formel bien} % [sic] f \in \cF -> p \in \cF, the factorisation is f = pi and lemme 4.4.3 3) restates this condition with p \in \cF
\msp{page 42}
\textit{connu}), \textit{que l'on a} $\mathrm{B}(\Fib_W)$, \textit{et a fortiori} (\textit{comme par hypothèse} $\cF \subset \Fib_W$) $\mathrm{B}(\cF)$.

\textbf{N.B.} Il est clair en effet que si on a $\cF \subset \cF'$, alors
\[
\mathrm{A}(\cF') \Longrightarrow \mathrm{A}(\cF) \qquad \text{et de même pour B, C.}
\tag*{$(*)$}
\]
Or la condition préliminaire sur $\cF$ est bien satisfaite dans le cas d'un ensemble de l'une des formes $\cF \;=\; \cF_W$, $\Fibf_W$, $\Parf_W$. Il suffit de le voir dans le cas des sous-ensembles $\cF_W$, $\Parf_W$ de $\Fibf_W$, \textit{i.e.} que toute flèche $f$ de Cat se factorise en
\[
f = pi,
\]
avec $i \in W$, et $p$ étant, au choix, dans $\Parf_W$, ou dans $\cF_W$. Or dans le cas de $\Parf_W$, \marginpar{$\|$}c'est ce qui est établie dans VII (Catégories de chemins (2)) grâce au formalisme des $\underline{\mathrm{Ch}}(S) \longrightarrow S \times S$. Pour trouver une factorisation avec $f = qj$, $j \in W$, $q \in \cF_W$, il suffit de voir que $p$ se factorise ainsi en $p = q_0 j_0$, on aura % [sic] établie -> établi, agreement with "ce qui"
\[
f = \underbrace{q_0}_{\in \cF_{W_0}} \underbrace{(j_0 i)}_{\in W},
\]
donc on peut supposer $f \in \Parf_W$. Mais si $X$ est $W$-parfait sur $S$, donc un
\msp{page 43}
$W$-fibré faible sur $S$, alors dans la factorisation
\begin{center}
\begin{tikzcd}[column sep=small]
X \arrow[rr, "i"] \arrow[dr, "f"'] & & \Phi_S(X) \arrow[dl, "p"] \\
 & S &
\end{tikzcd}
\end{center}
avec $i \in W$ et $p \in \Catfib$ (\textit{cf.} VI), on sait que $p$ est aussi une $W$-fibration faible, \textit{i.e.} $p \in \cF_W$, d'où la conclusion. Ainsi on trouve le

\textbf{Corollaire 4.4.2.} \textit{Supposons que $\cF$ soit l'un des six ensembles} (4.4.2) (i.e. \textit{l'un des ensembles du diagramme} (4.2.1), \textit{à l'exclusion de $\mathit{Fib}_W$ et $\mathit{Fib}_{W_0}$}). \textit{Alors on a}
\[
\mathrm{A}(\cF) \;\Longrightarrow\; \mathrm{B}(\cF) \;\Longrightarrow\; \mathrm{C}(\cF)\,.
\]

\textbf{Lemme 4.4.3.} \textit{Soit $(\cM, W)$ une catégorie de modèles, et soit $\cF \subset \Fl(\cM)$ satisfaisant les conditions suivantes}
\begin{itemize}[leftmargin=3em]
\item[1)] \textit{$\cF$ stable par changement de base, composition, contient les isomorphismes.}
\item[2)] \textit{$\cM$ a un objet final, et pour tout $X \in \mathrm{Ob}\,\cM$, $X \longrightarrow e$ est dans $\cF$.}
\item[3)] \textit{Toute flèche $f$ de $\cM$ se factorise en $f = pi$, avec $i \in W$, $p \in \cF$} (\textit{c'est la condition du lemme} 4.4.1).
\end{itemize}

\printedpage{26}
% head: A. Grothendieck : Dérivateurs, Ch. XVI, Édition des Universités de Montpellier II et Paris VII
% foot: Version 18 mars 2019
\noindent\textit{Alors les conditions} $\mathrm{A}(\cF)$, $\mathrm{B}(\cF)$, $\mathrm{C}(\cF)$ \textit{sur $\cF$ sont équivalentes.}

En effet, en vertu de 4.4.1, il reste à prouver l'implication
\[
\mathrm{C}(\cF) \;\Longrightarrow\; \mathrm{A}(\cF),
\]

\msp{page 44}

\textit{i.e.} ceci : si $\cF$ satisfait aussi la condition
\begin{itemize}[leftmargin=3em]
\item[4)] $\cF \cap W \subset W^{\mathrm{u}}$ (ou ce qui revient au même, comme $\cF$ est stable par changement de base, que $\cF \cap W$ est stable par changement de base),
\end{itemize}
alors on en conclut
\[
\cF \subset \Fib_W.
\]
Or ceci n'est autre que le lemme de K. S. \textsc{Brown}, sous la forme duale de celle donnée plus haut (3.2.1, page 20).

Parmi les six ensembles de flèches envisagés dans (4.4.2), les hypothèses 1) 2) 3) sont satisfaites pour $\Parf_W$ et $\Parf_{W_0}$, à l'exclusion des quatre autres. (Les autres satisfont à toutes les conditions, à la seule exception de la stabilité par composition \footnote[30]{\textbf{N.B.} Pour $\Fib_{W_0}$, ce sont les conditions 1) 2) qui sont vérifiées, et par contre, la condition 3) qui fait problème. Notons là $\mathrm{F}(W_0)$. Si elle est satisfaite (\textit{i.e.} si Cat a assez de $W_0$-fibrations), alors on a aussi équivalence des conditions $\mathrm{A}(\Fib_{W_0})$, $\mathrm{B}(\Fib_{W_0})$, $\mathrm{C}(\Fib_{W_0})$.}). On trouve donc :

\textbf{Corollaire 4.4.4.} \textit{Les conditions} $\mathrm{A}(\Parf_W)$, $\mathrm{B}(\Parf_W)$, $\mathrm{C}(\Parf_W)$ \textit{sont équivalentes} (nous désignons cette condition sur $W$ par D). \textit{De même les conditions} $\mathrm{A}(\Parf_{W_0})$, $\mathrm{B}(\Parf_{W_0})$, $\mathrm{C}(\Parf_{W_0})$ \textit{sont équivalentes.}

\msp{page 45}

(On désignera la condition obtenue sur le couple $W \supset W_0$ de localiseurs fondamentaux par $D_{W_0}$.)

Ainsi on a (comme $\Parf_{W_0} \subset \Parf_W$)
\[
\mathrm{D} = \mathrm{D}_W \;\Longrightarrow\; \mathrm{D}_{W_0}\,,
\tag*{(4.4.4)}
\]
et si $W'_0 \subset W_0$, on aura
\[
\mathrm{D}_{W_0} \;\Longrightarrow\; \mathrm{D}_{W'_0}\,.
\tag*{(4.4.5)}
\]

\textbf{Corollaire 4.4.5.} \textit{Désignons par} A, B, C \textit{respectivement les conditions} Loc (7 a), Loc (7 b), Loc (7 c). \textit{On a alors les implications}
\[
\mathrm{A} \;\Longrightarrow\; \mathrm{B} \;\Longrightarrow\; \mathrm{C}\,.
\tag*{(4.4.5.1)}
\]
\textit{De plus, si on désigne} (\textit{pour un ensemble $\cF$ de flèches de} Cat) \textit{par} $\Comp(\cF)$ \textit{la condition sur $\cF$ que $\cF$ soit stable par composition, on a les équivalences}
\[
\mathrm{A} \iff (\mathrm{C} + \Comp(\cF_W)) \iff (\mathrm{C} + \Comp(\Fibf_W)).
\tag*{(4.4.5.2)}
\]

\printedpage{27}
% head: A. Grothendieck : Dérivateurs, Ch. XVI, Édition des Universités de Montpellier II et Paris VII
% foot: Version 18 mars 2019

\textsc{Démonstration}. La première assertion est contenue dans le corollaire 4.4.2, appliqué au cas $\cF = \cF_W$. D'autre part, si $\cF_W$ satisfait à $\Comp(\cF_W)$, il satisfait aux conditions 1) 2) 3) du lemme de Brown 4.4.3, donc joint à C, il implique A. Il reste, pour la première équivalence, à prouver que
\[
A \;\Longrightarrow\; \Comp(\cF_W).
\]

\msp{page 46}

Mais on a par définition
\[
\cF_W = \Catfib \cap \Fibf_W
\tag*{$(*)$}
\]
et on sait que Catfib est stable par composition. Donc $\cF_W$ l'est aussi, pourvu que le composé $gf$ de deux flèches dans $\cF_W$ soit dans $\Fibf_W$. Or par hypothèse A, qui signifie $\cF_W \subset \Fib_W$, on a $f, g \in \Fib_W$, donc $gf \in \Fib_W$, puisque $\Fib_W$ est stable par composition, a fortiori $gf \in \Fibf_W$.

D'autre part, la formule précédente $(*)$ montre que $\Comp(\Fibf_W) \Longrightarrow \Comp(\cF_W)$, donc on a gratis l'implication $\Longleftarrow$ dans la deuxième équivalence (4.4.5.2). L'implication en sens inverse revient alors à
\[
\mathrm{A} \;\Longrightarrow\; \Comp(\Fibf_W),
\]

\msp{page 47}

ce qui sera évident grâce au résultat (donné plus bas)
\[
\mathrm{A} = \mathrm{A}(\cF_W) \iff \mathrm{A}(\Fibf_W) \qquad \text{et de même pour B, C,}
\]
qui signifie donc que
\[
\cF_W \subset \Fib_W \qquad \text{implique} \qquad \Fibf_W \subset \Fib_W \qquad (\text{donc } \Fibf_W = \Fib_W),
\]
or $\Fib_W$ étant stable par composition, il en est de même pour $\Fibf_W$, q.e.d.

\textbf{Corollaire 4.4.6.} \textit{Supposons que $\cF_{W_0}$ soit stable par composition, ce qui est le cas notamment} (\textit{par} 4.4.5.2 \textit{appliqué à $W_0$ au lieu de $W$}) \textit{si $W_0$ satisfait} Loc (7). \textit{Alors les conditions} $\mathrm{A}(\cF_{W_0})$, $\mathrm{B}(\cF_{W_0})$, $\mathrm{C}(\cF_{W_0})$ (\textit{que nous notons aussi simplement} $\mathrm{A}_{W_0}$, $\mathrm{B}_{W_0}$, $\mathrm{C}_{W_0}$) \textit{sont équivalentes.} (\textbf{\textit{N.B.}} \textit{Si on ne fait pas l'hypothèse précédente sur $W_0$, on aura seulement} $\mathrm{A}_{W_0} \Longrightarrow \mathrm{B}_{W_0} \Longrightarrow \mathrm{C}_{W_0}$, \textit{en vertu de} 4.4.2.)

4.4.7. On a fait le tour à présent de toutes les conditions de type “fibration” sur $W$, ou sur le couple $(W, W_0)$, qui m'ont parues utiles. Je vais les déployer ici dans un diagramme (plus ou moins) récapitulatif : % [sic] parues -> paru, agreement of the past participle

\msp{page 48}

\printedpage{28}
% head: A. Grothendieck : Dérivateurs, Ch. XVI, Édition des Universités de Montpellier II et Paris VII
% foot: Version 18 mars 2019

\begin{center}
\begin{tikzcd}[column sep=normal, row sep=small]
 & \overbrace{\mathrm{A}_{W_0}}^{\cF_{W_0},\,\Fibf_{W_0}} \arrow[dd, Rightarrow] & & \\
\overbrace{\mathrm{A}}^{\cF_W,\,\Fibf_W} \arrow[ur, Rightarrow] \arrow[dd, Rightarrow] & & & \\
 & \mathrm{B}_{W_0} \arrow[dd, Rightarrow] & & \\
\mathrm{B} \arrow[ur, Rightarrow] \arrow[dd, Rightarrow] & & & \\
 & \mathrm{C}_{W_0} \arrow[dr, Rightarrow, "\bullet"] \arrow[uuuu, Rightarrow, dotted, bend right=40, "\text{\scriptsize si } W_0\ \text{\scriptsize (7)}"'] & & \\
\mathrm{C} \arrow[ur, Rightarrow] \arrow[dr, Rightarrow, "\bullet"] & & \overbrace{\mathrm{D}_{W_0}}^{\Parf_{W_0}} \arrow[r, Rightarrow, "\bullet"] & \mathrm{F} : (\Cat, W, \Fib_W) \text{ une catégorie à fibrations.} \\
 & \overbrace{\mathrm{D}}^{\Parf_W} \arrow[ur, Rightarrow] & &
\end{tikzcd}
\end{center}

\begin{center}
\begin{tikzcd}[row sep=large]
 & \mathrm{A} \arrow[r, Rightarrow] \arrow[d, Rightarrow] & \mathrm{A}_{W_0} \arrow[d, Rightarrow] & \\
\cF_W, \Fibf_W \smash{\left\{\vrule height 2.6cm depth 2.6cm width 0pt\right.} & \mathrm{B} \arrow[r, Rightarrow] \arrow[d, Rightarrow] & \mathrm{B}_{W_0} \arrow[d, Rightarrow] & \smash{\left.\vrule height 2.6cm depth 2.6cm width 0pt\right\}} \quad \text{\scriptsize\begin{tabular}{@{}l@{}}version revue\\ du diagramme\end{tabular}} \\
 & \mathrm{C} \arrow[r, Rightarrow] \arrow[d, Rightarrow] & \mathrm{C}_{W_0} \arrow[d, Rightarrow] & \\
 & \underbrace{\mathrm{D}}_{\Parf_W} \arrow[r, Rightarrow] & \underbrace{\mathrm{D}_{W_0}}_{\Parf_{W_0}} \arrow[r, Rightarrow] & \mathrm{F},
\end{tikzcd}
\end{center}

diagramme que je vais commenter. Toutes les implications en traits pleins du diagramme ont été déjà obtenues ou sont tautologiques, à l'exception des trois $\overset{\bullet}{\Longrightarrow}$ sur lesquels je vais revenir. L'implication en pointillés $\mathrm{C}_{W_0} \dottedimplies \mathrm{A}_{W_0}$ est une implication conditionnelle, soumise à la condition marquée sur le diagramme comme $W_0$ (7), ce qui signifie : $W_0$ satisfait à Loc (7). Cette implication, \textit{i.e.} l'équivalence de $\mathrm{A}_{W_0}$, $\mathrm{B}_{W_0}$, $\mathrm{C}_{W_0}$ sous la condition $W_0$ (7), a été vue dans 4.4.6. J'ai marqué, au dessus des deux groupes (A, B, C) et ($\mathrm{A}_{W_0}$, $\mathrm{B}_{W_0}$, $\mathrm{C}_{W_0}$) ainsi qu'au dessus de D et de D$'$ (qui correspondent chacun à un tel groupe de conditions $\mathrm{A}(\cF)$, $\mathrm{B}(\cF)$, $\mathrm{C}(\cF)$ toutes équivalentes, avec $\cF \;=\; \Parf_W$ ou $\cF = \Parf_{W_0}$) les
\msp{page 49}
ensembles de flèches qui leur donnent naissance respectivement : on notera qu'il y en a deux pour les pour les deux triplets nommés en premier, A, B, C pouvant se définir via $\cF_W$ ou $\Fib_W$ indifféremment, et de même $\mathrm{A}_{W_0}$, $\mathrm{B}_{W_0}$, $\mathrm{C}_{W_0}$ via $\cF_{W_0}$ ou $\Fib_{W_0}$, comme déjà énoncé (4.4.3) et comme nous le verrons plus bas (cor. 4.4.10, pages 54-55). Par contre, les triplets réduits D et $\mathrm{D}_{W_0}$ sont définis chacun par un $\cF$ seulement, savoir par $\Parf_W$ et % [sic] pour les pour les -> pour les, words repeated; [sic] Fib_W / Fib_{W_0} -> Fibf_W / Fibf_{W_0}, the diagram and (4.4.3) give Fibf

% ---- from der-ch16_p031-035.tex
\setcounter{footnote}{30}

\printedpage{29}
% head: A. Grothendieck : Dérivateurs, Ch. XVI, Édition des Universités de Montpellier II et Paris VII
% foot: Version 18 mars 2019
$\Parf_{W_0}$ respectivement. (En tout, cela fait intervenir six ensembles $\mathcal{F}$ distincts, à savoir les six ensembles du diagramme (4.2.1) autres que $\Fib_W$ et $\Fib_{W_0}$.)

L'implication non tautologique $\mathrm{C} \dotimp \mathrm{D}$ est un cas particulier de l'implication $\mathrm{C}_{W_0} \dotimp \mathrm{D}_{W_0}$, laquelle est contenue dans ce qui a été annoncé à l'instant, à savoir $\mathrm{C}_{W_0}(\mathcal{F}_{W_0}) \Longrightarrow \mathrm{C}_{W_0}(\Fib_{W_0})$, compte tenu que $\Parf_{W_0} \subset \Fibf_{W_0}$ (\textit{cf.} $(*)$, p. 42). Faisant
\msp{page 50}
$W_0 = W$, on trouve donc l'implication $\mathrm{A} \Longrightarrow \mathrm{D}$, qui est cependant contenue dans les autres implications du diagramme, comme composée dans la chaîne d'implications
\[
\mathrm{A} \Longrightarrow \mathrm{B} \Longrightarrow \mathrm{C} \Longrightarrow \mathrm{D} \Longrightarrow \mathrm{D}_{W_0} \dotimp \mathrm{F}.
\]

Il reste, dans ces commentaires généraux, à dire un mot de cette dernière implication
\[
\mathrm{D}_{W_0} \dotimp \mathrm{F},
\]
\textit{i.e.} du fait que la plus faible de toutes les conditions A à $\mathrm{D}_{W_0}$ envisagées ici implique déjà que Cat a assez de $W$-fibrations, \textit{i.e.} est une catégorie à cofibrations pour $(W, \Fib_W)$. C'est clair quand on prend $\mathrm{D}_{W_0}$ sous la forme $\mathrm{A}(\Parf_{W_0})$, \textit{i.e.}
\[
\Parf_{W_0} \subset W.
\]
Comme toute flèche $f$ de Cat se factorise en $f = pi$, avec
\[
i \in W_0 \subset W \quad\text{et}\quad f \in \Parf_{W_0} \subset \Fib_W\, , % [sic] f \in \Parf_{W_0} -> p \in \Parf_{W_0}, the factor of f = pi lying in Parf is p
\]
on trouve la factorisation avec $i \in W$, $f \in \Fib_W$. % [sic] f \in \Fib_W -> p \in \Fib_W, the factorisation f = pi requires p \in \Fib_W

\msp{page 51}
4.4.8. Pour terminer d'établir les implications du diagramme p. 48, il reste seulement à établir les équivalences
\[
\left\{
\begin{aligned}
&\mathrm{A}(\mathcal{F}_{W_0}) &&\Longleftrightarrow &&\mathrm{A}(\Fibf_{W_0})\\
&\mathrm{B}(\mathcal{F}_{W_0}) &&\Longleftrightarrow &&\mathrm{B}(\Fibf_{W_0})\\
&\mathrm{C}(\mathcal{F}_{W_0}) &&\Longleftrightarrow &&\mathrm{C}(\Fibf_{W_0}),
\end{aligned}
\right.
\tag*{(4.4.8.1)}
\]
qui donneront les équivalences similaires pour $W$, en prenant $W_0 = W$.

La première équivalence (4.4.8.1) sera conséquence du lemme ($\pm$ tautologique) suivant :

\textbf{Lemme 4.4.9.} \textit{Soient} $(\mathcal{M}, W)$ \textit{une catégorie de modèles,} $\mathcal{F}, \mathcal{G} \subset \Fl(\mathcal{M})$, \textit{tels que tout} $f \in \mathcal{F}$, $f : X \longrightarrow S$, \textit{s'écrive} $f = pi$, \textit{avec} $i \in W^u_S$, $p \in \mathcal{G}$. \textit{Alors} $\mathrm{A}(\mathcal{G}) \Longrightarrow \mathrm{A}(\mathcal{F})$, i.e.
\[
\mathcal{G} \subset \Fib_W \;\Longrightarrow\; \mathcal{F} \subset \Fib_W,
\tag*{(4.4.9.1)}
\]
\textit{et de même} $\mathrm{C}(\mathcal{G}) \Longrightarrow \mathrm{C}(\mathcal{F})$, i.e.
\[
\mathcal{G} \cap W \subset W^u \;\Longrightarrow\; \mathcal{F} \cap W \subset W^u.
\tag*{(4.4.9.2)}
\]

\textsc{Démonstration}. A) En effet, si $\mathcal{G} \subset \Fib_W$, la factorisation $f = pi$ de $f$ est telle que $p \in \Fib_W$. Or un composé $pi$, avec $p \in \Fib_W$ et $i \in W^{\mathrm{u}}_S$, est dans $\Fib_W$, comme on vérifie aussitôt sur le diagramme de carrés cartésiens
\printedpage{30}
% head: A. Grothendieck : Dérivateurs, Ch. XVI, Édition des Universités de Montpellier II et Paris VII
% foot: Version 18 mars 2019
\[
\begin{tikzcd}
X \arrow[d, "W^{\mathrm{u}}_S"'] & X' \arrow[l] \arrow[d, "(W)"'] & X'' \arrow[l] \arrow[d, "(W)"'] \\
\overline{X} \arrow[d, "\Fib_W"'] & \overline{X'} \arrow[l] \arrow[d] & \overline{X''} \arrow[l, "(W)"'] \arrow[d] \\
S & S' \arrow[l] & S'' \arrow[l, "W"']
\end{tikzcd}
\]
sur lequel on lit que $X'' \longrightarrow X'$ est dans $W$, q.e.d.

\msp{page 52}
Ce lemme s'applique notamment au cas où $\mathcal{F} = \Fibf_{W_0}$, $\mathcal{G} = \mathcal{F}_{W_0}$, car $f$ se factorise en
\[
\begin{tikzcd}
X \arrow[dr, "i"] \arrow[dd, "f"'] & \\
& \Phi_S(X) \arrow[dl, "p"] \\
S, &
\end{tikzcd}
\]
et on sait que $i \in W^{\mathrm{u}}_{0\,S} \subset W^{\mathrm{u}}_S$, et $p \in \mathcal{F}_{W_0} \subset \mathcal{F}_W$, ce qui sont les conditions voulues.

C) Démonstration similaire : Soit $f \in \mathcal{F} \cap W$, factorisant en $f = pi$, avec $i \in W^{\mathrm{u}}_S$, $g \in \mathcal{G}$, on aura alors $g \in W$, donc $g \in \mathcal{G} \cap W \subset W^{\mathrm{u}}$ (par hypothèse $\mathrm{C}(\mathcal{G})$), a fortiori on a $g \in W^{\mathrm{u}}_S$, donc $f \in W^{\mathrm{u}}_S$ comme composé de deux flèches de $W^{\mathrm{u}}_S$, et comme le but de $f$ est $S$ lui-même, cela signifie $f \in W^{\mathrm{u}}$, q.e.d.
\[
\begin{tikzcd}
X \arrow[dr, "i\in W^{\mathrm{u}}_S"] \arrow[dd, "\mathcal{F}\cap W"'] & \\
& \bar{X} \arrow[dl, "g"] \\
S &
\end{tikzcd}
\]

Je ne vois pas de façon pour prouver également
\[
\mathrm{B}(\mathcal{G}) \qquad\quad \Longrightarrow \qquad\quad \mathrm{B}(\mathcal{F}),
\]
avec les seules hypothèses du lemme précédent sur le couple $\mathcal{F}$, $\mathcal{G}$. Il faudrait que pour tout $S$ dans $\mathcal{M}$, et toute flèche
\msp{page 53}
$u : X \longrightarrow Y$ dans $\mathcal{M}/S$, on puisse insérer le diagramme obtenu dans $\mathcal{M}$
\[
\begin{tikzcd}[column sep=small]
X \arrow[rr, "u"] \arrow[dr, "f"'] & & Y \arrow[dl, "g"] \\
& S &
\end{tikzcd}
\]
\printedpage{31}
% head: A. Grothendieck : Dérivateurs, Ch. XVI, Édition des Universités de Montpellier II et Paris VII
% foot: Version 18 mars 2019
dans un diagramme commutatif ci-dessous
\[
\begin{tikzcd}[column sep=small]
X \arrow[rr, "u"] \arrow[d, "i"'] & & Y \arrow[d, "j"] \\
\bar{X} \arrow[rr, "\bar{u}"] \arrow[dr, "p"'] & & \bar{Y} \arrow[dl, "q"] \\
& S &
\end{tikzcd}
\tag*{(4.4.9.3)}
\]
avec
\[
f = pi, \qquad\quad g = qj, \qquad\quad i, j \in W^{\mathrm{u}}_S, \qquad\quad p, q \in \mathcal{G}.
\]
Ce sera le cas en tous cas si pour tout $S$, et un objet variable $X$ de $\mathcal{M}/S$, on peut trouver une factorisation \textit{fonctorielle en} $X$ de son morphisme structural $f_X$ en $p_X i_X$, avec $i_X \in W^{\mathrm{u}}_S$ et $p_X \in \mathcal{G}$, \textit{i.e.} si $i$, $p$ se déduisent d'un foncteur
\[
\varphi : \mathcal{F}_S \longrightarrow \mathcal{M}/S \ \text{\footnotemark},
\]
\footnotetext{$\mathcal{F}_S$ est la sous-catégorie pleine de $\mathcal{M}/S$ formée des $X$ tels que $f_X \in \mathcal{F}$.}
et d'un morphisme de foncteurs
\[
i : \underbrace{\inc_{\mathcal{F}_S}}_{\substack{\text{foncteur}\\ \text{d'inclusion}}} \longrightarrow \varphi.
\]
Admettant l'existence d'un diagramme ayant les propriétés dites, pour toute flèche $u$ dans $\mathcal{F}_S$, ou seulement toute flèche $u$ qui soit aussi dans $W_S$, nous pouvons prouver maintenant que $u \in W_S \Longrightarrow u \in W^{\mathrm{u}}_S$, pourvu que la propriété similaire soit vrai pour $\mathcal{G}$ à la place de $\mathcal{F}$.

\msp{page 54}
En effet, l'hypothèse sur $u \in W$ implique $\bar{u} \in W$, et l'hypothèse sur $\mathcal{G}$ implique dès lors $\bar{u} \in W^{\mathrm{u}}_S$. Comme $W^{\mathrm{u}}_S$ est un localiseur dans $\mathcal{M}/S$, on en conclut bien $u \in W^{\mathrm{u}}_S$, q.e.d.

Ainsi on a obtenu

\textbf{Corollaire 4.4.10.} \textit{Supposons, sous les conditions du lemme} 4.4.9, \textit{que tout diagramme commutatif triangulaire} $\begin{tikzcd}[column sep=small, row sep=small, ampersand replacement=\&]
X \arrow[rr, "u"] \arrow[dr, "f"'] \& \& Y \arrow[dl, "g"] \\
\& S \&
\end{tikzcd}$ \textit{s'insère dans un diagramme du type} (4.4.9.3) \textit{comme ci-dessus} (\textit{avec les relations explicites sur les flèches qui y figurent}). \textit{Alors}
\[
\mathrm{B}(\mathcal{G}) \;\Longrightarrow\; \mathrm{B}(\mathcal{F}).
\]
\textit{Donc si de plus} $\mathcal{G} \subset \mathcal{F}$, \textit{de sorte que} $\mathrm{A}(\mathcal{F}) \Longrightarrow \mathrm{A}(\mathcal{G})$ \textit{tautologiquement, et de même pour} B \textit{et pour} C, \textit{on trouve}
\[
\begin{aligned}
\mathrm{A}(\mathcal{F}) &\;\Longleftrightarrow\; \mathrm{A}(\mathcal{G})\\
\mathrm{B}(\mathcal{F}) &\;\Longleftrightarrow\; \mathrm{B}(\mathcal{G})\\
\mathrm{C}(\mathcal{F}) &\;\Longleftrightarrow\; \mathrm{C}(\mathcal{G}).
\end{aligned}
\]

Nous pouvons appliquer ce corollaire dans le cas qui nous occupe ici, avec
\[
\mathcal{G} \overset{\text{déf}}{=} \mathcal{F}_{W_0} \subset \mathcal{F} \overset{\text{déf}}{=} \Fibf_{W_0},
\]
\printedpage{32}
% head: A. Grothendieck : Dérivateurs, Ch. XVI, Édition des Universités de Montpellier II et Paris VII
% foot: Version 18 mars 2019
\msp{page 55}
vu que dans ce cas on a bien une factorisation \textit{fonctorielle} d'un $f \in \Fibf_{W_0}$, $f : X \longrightarrow S$, en $f = pi$, $i \in W^{\mathrm{u}}_{0\,S} \subset W^{\mathrm{u}}_S$, $p \in \mathcal{F}_{W_0} = \mathcal{G}$, savoir celle rappelée page 43 :
\[
\begin{tikzcd}[column sep=small]
X \arrow[rr, "i_X"] \arrow[dr, "f=f_X"'] & & \bar{X} = \Phi_S(X) \arrow[dl, "p_X"] \\
& S. &
\end{tikzcd}
\]

Cela prouve donc les équivalences annoncées (4.4.8.1) et achève d'établir la validité des implications du diagramme de la page 48.

Ces implications, à la seule exception de l'implication en pointillés $\mathrm{C}_{W_0} \dottedimpl \mathrm{A}_{W_0}$ (qui suppose qu'on ait $W_0$ (7)), ne font appel à aucune hypothèse particulière sur les localiseurs fondamentaux $W$ et $W_0$ (à part l'inclusion $W \supset W_0$). Mais nous pouvons, en utilisant le lemme 4.4.10, donner des implications conditionnelles remarquables, sous la seule condition que $W$, ou $W_0$, satisfassent à l'anodin \textit{axiome}
\msp{page 56}
\textit{des limites} Loc (8), et même seulement sous la forme faible de stabilité des localiseurs par $\varinjlim$ de \textit{suites}. Donnons-les d'abord sous forme technique, avant de déployer les conséquences dans un théorème scholie :

\textbf{Corollaire 4.4.11.} \textit{Supposons que} $W_0$ \textit{satisfasse à l'axiome} Loc (8), \textit{ou seulement} Loc (8 bis) (\textit{stabilité par} $\varinjlim$ \textit{de suites}). \textit{Alors} $\mathrm{D}_{W_0} \Longrightarrow \mathrm{A}_{W_0}$ (\textit{donc les conditions} $\mathrm{A}_{W_0}$, $\mathrm{B}_{W_0}$, $\mathrm{C}_{W_0}$, $\mathrm{D}_{W_0}$ \textit{sont équivalentes}). \textit{En particulier} (\textit{faisant} $W_0 = W$), \textit{si} $W$ \textit{satisfait à} Loc (8 bis), \textit{alors on a} $\mathrm{D}_W \Longrightarrow \mathrm{A}_W$ (\textit{donc les conditions} $\mathrm{A}_W$, $\mathrm{B}_W$, $\mathrm{C}_W$, $\mathrm{D}_W$ \textit{sont équivalentes}).

En effet, l'implication à prouver s'écrit aussi
\[
\mathrm{A}(\Parf_{W_0}) \;\Longleftrightarrow\; \mathrm{A}(\Fibf_{W_0}).
\]

L'implication non tautologique $\Longrightarrow$ est justiciable du lemme 4.4.9. Il faut vérifier que moyennant Loc (8 bis) pour $W_0$, tout $f \in \Fib_{W_0}$, $f : X \longrightarrow S$, se factorise en $f = pi$, avec $i \in W^{\mathrm{u}}_0$, $p \in \Parf_{W_0}$. Or on sait\footnote{et c'est là une trouvaille tout à fait inattendue, un peu comme le lemme de K. S. \textsc{Brown} \ldots} (VI, p. 193 ff.) % [sic] f \in \Fib_{W_0} -> f \in \Fibf_{W_0}, the implication to be proved concerns A(\Fibf_{W_0}) and the next line has \Fibf_{W_0}
\msp{page 57}
que l'axiome Loc (8 bis) implique que $f \in \Fibf_{W_0}$ se factorise en
\[
\begin{tikzcd}[column sep=small]
X \arrow[rr, "i_\infty"] \arrow[dr, "f"'] & & \bar{X} = \Phi_\infty(X) \arrow[dl, "p_\infty"] \\
& S, &
\end{tikzcd}
\qquad
\text{\parbox[t]{0.33\textwidth}{\footnotesize (catégorie de chemins infinis d'un côté, relativisée sur $S$, \textit{i.e.} des couples $(x,c)$, avec $x \in \Ob X$ et $c : f(x) \longrightarrow s$ un chemin infini dans $S$, de source $f(x)$.)}}
\]
avec $p_\infty \in \Parf_{W_0}$ (comme on veut), et $i_\infty \in W^{\mathrm{u}}_{0\,S}$, et comme $W^{\mathrm{u}}_{0\,S} \subset W^{\mathrm{u}}_S$, on a bien $i_\infty \in W^{\mathrm{u}}_S$.

\textbf{Corollaire 4.4.12.} \textit{Supposons que} $W$ et $W_0$ \textit{satisfassent à} Loc (8 bis). \textit{Alors on a} $\mathrm{D}_{W_0} \Longrightarrow \mathrm{A}$ (\textit{donc les huit conditions} A, B, C, D, $\mathrm{A}_{W_0}$, $\mathrm{B}_{W_0}$, $\mathrm{C}_{W_0}$, $\mathrm{D}_{W_0}$ \textit{sont équivalentes}).

\printedpage{33}
% head: A. Grothendieck : Dérivateurs, Ch. XVI, Édition des Universités de Montpellier II et Paris VII
% foot: Version 18 mars 2019
Cela signifie en effet que
\[
\mathrm{A}(\Parf_{W_0}) \;\Longrightarrow\; \mathrm{A}(\Fibf_W),
\]
et c'est encore justiciable du lemme 4.4.9. En effet, la factorisation ci-dessus d'une flèche $f$ ne dépend pas du choix d'un localiseur, et $p_\infty$ est $W$-parfait pour \textit{tout} localiseur satisfaisant Loc (8 bis), donc en particulier pour $W_0$. Et la flèche $i_\infty$ est dans $W^{\mathrm{u}}_S$, pour tout localiseur satisfaisant Loc (8 bis), donc pour $W_0$.

\msp{page 58}
Je vais récapituler les résultats obtenus dans cette section dans un théorème, qui constituera le résultat principal du présent paragraphe sur les axiomes de $W$-fibrations.

\textbf{Théorème 4.4.13.} \textit{Soit} $W$ \textit{un localiseur fondamental dans} Cat. \textit{Pour tout localiseur} $W_0$ \textit{contenu dans} $W$, \textit{on a explicité un ensemble de quatre conditions sur la paire} $W_0$, $W$, \textit{qui pour} $W$ \textit{fixé, dépendent de} $W_0$ \textit{seul, notées}
\[
\mathrm{A}_{W_0},\ \mathrm{B}_{W_0},\ \mathrm{C}_{W_0},\ \mathrm{D}_{W_0},
\tag*{(4.4.13.1)}
\]
\textit{et définies respectivement comme}
\[
\mathrm{A}(\mathcal{F}_{W_0}),\ \mathrm{B}(\mathcal{F}_{W_0}),\ \mathrm{C}(\mathcal{F}_{W_0}), \qquad\quad \textit{et}
\]
\[
\mathrm{A}(\Parf_{W_0}) \Bigl(\Longleftrightarrow\; \mathrm{B}(\Parf_{W_0}) \;\Longleftrightarrow\; \mathrm{C}(\Parf_{W_0})\Bigr),
\]
\textit{avec les notations} $\mathrm{A}(\cdot)$, $\mathrm{B}(\cdot)$, $\mathrm{C}(\cdot)$ \textit{introduites dans} (4.4.1). \textit{Faisant} $W_0 = W$, \textit{on trouve les conditions sur} $W$
\[
\mathrm{A},\ \mathrm{B},\ \mathrm{C},\ \mathrm{D}
\tag*{(4.4.13.2)}
\]
\textit{avec} $\mathrm{A} = \mathrm{A}_W$, $\mathrm{B} = \mathrm{B}_W$, $\mathrm{C} = \mathrm{C}_W$, $\mathrm{D} = \mathrm{D}_W$, \textit{dont les trois premières ne sont autres que} Loc (7 a), Loc (7 b), Loc (7 c) (cf. 4.3, \textit{page} 38).

\circled{1} \textit{On a pour ces conditions les implications suivantes.}

\begin{enumerate}[leftmargin=2.5em]
\item[a)] \textit{On a les implications}
\[
\mathrm{A}_{W_0} \;\Longrightarrow\; \mathrm{B}_{W_0} \;\Longrightarrow\; \mathrm{C}_{W_0} \;\Longrightarrow\; \mathrm{D}_{W_0},
\tag*{(4.4.13.3)}
\]
(\textit{d'où en particulier}

\msp{page 59}
\[
\mathrm{A}_W \;\Longrightarrow\; \mathrm{B}_W \;\Longrightarrow\; \mathrm{C}_W \;\Longrightarrow\; \mathrm{D}_W\ ).
\tag*{(4.4.13.4)}
\]
\item[b)] \textit{Si} $W'_0 \subset W_0$ \textit{est un localiseur fondamental contenu dans} $W_0$, \textit{on a} (\textit{tautologiquement})
\[
\mathrm{A}_{W_0} \;\Longrightarrow\; \mathrm{A}_{W'_0}, \qquad\quad \mathrm{B}_{W_0} \;\Longrightarrow\; \mathrm{B}_{W'_0}, \qquad\quad \mathrm{C}_{W_0} \;\Longrightarrow\; \mathrm{C}_{W'_0}.
\]
\textit{On trouve donc, par} (a) \textit{et} (b) (\textit{y faisant pour le couple} $(W, W_0)$ \textit{comme} $(W_0, W'_0)$),
\end{enumerate}

% ---- from der-ch16_p036-040.tex
\begin{enumerate}[leftmargin=2.5em]
\item[] \printedpage{34}%
% head: A. Grothendieck : Dérivateurs, Ch. XVI, Édition des Universités de Montpellier II et Paris VII
\textit{le rectangle d'implications}
\[
\begin{tikzcd}[row sep=large, column sep=large, execute at end picture={
\draw[decorate,decoration={brace,amplitude=6pt}] ([xshift=-14pt]A.west |- C.south) -- ([xshift=-14pt]A.north west) node[midway,left=8pt]{$\mathcal{F}_W,\mathrm{Fibf}_W$};
\draw[decorate,decoration={brace,amplitude=6pt}] ([xshift=14pt]AW.north east) -- ([xshift=14pt]AW.east |- CW.south) node[midway,right=8pt]{$\mathcal{F}_{W_0},\mathrm{Fibf}_{W_0}$};
}]
|[alias=A]| \mathrm{A} \arrow[r,Rightarrow] \arrow[d,Rightarrow] & |[alias=AW]| \mathrm{A}_{W_0} \arrow[d,Rightarrow] \\
\mathrm{B} \arrow[r,Rightarrow] \arrow[d,Rightarrow] & \mathrm{B}_{W_0} \arrow[d,Rightarrow] \\
|[alias=C]| \mathrm{C} \arrow[r,Rightarrow] \arrow[d,Rightarrow] & |[alias=CW]| \mathrm{C}_{W_0} \arrow[d,Rightarrow] \\
\underbrace{\mathrm{D}}_{\mathrm{Parf}_W} \arrow[r,Rightarrow] & \underbrace{\mathrm{D}_{W_0}}_{\mathrm{Parf}_{W_0}}\,,
\end{tikzcd}
\tag*{(4.4.13.5)}
\]
\textit{où on a rappelé, à côté des accolades, quels sont l'ensemble, ou les ensembles $\mathcal{F}\subset\mathrm{Fl}(\mathrm{Cat})$ qui définissent les conditions envisagées (comme l'une des conditions $\mathrm{A}(\mathcal{F})$, $\mathrm{B}(\mathcal{F})$, $\mathrm{C}(\mathcal{F})$ ou $\mathrm{D}(\mathcal{F})=\mathrm{A}(\mathcal{F})$ dans le cas où les trois conditions sont équivalentes). La plus faible des huit conditions envisagées est $\mathrm{D}_{W_0}$, la plus forte} A. (\textit{Et pour $W_0$ variable, on trouve la condition la plus faible de toutes en prenant $D_{W_\omega}$, où $W_\omega$ est le plus petit des localiseurs fondamentaux dans} Cat.)

\item[c)] \textit{On a de plus, pour la condition la plus faible $\mathrm{D}_{W_0}$, l'implication}
\[
\mathrm{D}_{W_0} \implies \mathrm{F},
\]
\msp{page 60}%
\textit{où} $\mathrm{F}=\text{Loc (7 F)}$ \textit{est la condition que} Cat \textit{possède suffisamment de $W$-fibrations,} i.e. \textit{que c'est une catégorie de fibrations pour} $(W,\mathrm{Fib}_W)$.
\end{enumerate}

\noindent\cercle{2}\enspace{\itshape Si $W$ et $W_0$ satisfont à \textup{Loc (8 bis)} (axiome des limites dénombrables), alors \textup{$\mathrm{D}_{W_0}\Longrightarrow\mathrm{A}$, i.e.} les 8 conditions envisagées sont équivalentes. Donc si $W$ satisfait \textup{Loc (8 bis)}, considérant le plus petit localiseur fondamental $W_L$ satisfaisant à \textup{Loc (8 bis)} (de sorte que $W\supset W_L$), on voit que les huit conditions
\[
\mathrm{A},\ \mathrm{B},\ \mathrm{C},\ \mathrm{D},\ \mathrm{A}_{W_L},\ \mathrm{B}_{W_L},\ \mathrm{C}_{W_L},\ \mathrm{D}_{W_L}
\]
sont équivalentes~\footnote[33]{N.B. La forme la plus faible, $\mathrm{D}_{W_{\mathrm{L}}}$, de ces conditions, est la relation \fbox{$\mathrm{Parf}_{W_{\mathrm{L}}}\cap W\subset W^{\mathrm{u}}$.}}. En particulier, sous la même hypothèse (fort anodine) sur $W$, les conditions \textup{Loc (7) $=$ Loc (7 a), Loc (7 b)} et \textup{Loc (7 c)} sont équivalentes, et elles équivalent aussi à la condition \textup{$\mathrm{D}=\mathrm{A}(\mathrm{Parf}_W)$, i.e.} à
\[
\mathrm{Parf}_W\subset\mathrm{Fib}_W \qquad (\textit{cas particulier de } \text{Loc (7 a)}),
\]
et à ses deux variantes d'apparence plus faible \textup{$\mathrm{B}(\mathrm{Parf}_W)$, $\mathrm{C}(\mathrm{Parf}_W)$} (cas particulier de \textup{Loc (7 b) et de Loc (7 c)} respectivement).}

% foot: Version 18 mars 2019

\noindent\cercle{3}\enspace{\itshape Supposons que $W_0$ satisfasse à \textup{Loc (8 bis)} (sans hypothèse particulière sur $W$). Alors
\printedpage{35}%
% head: A. Grothendieck : Dérivateurs, Ch. XVI, Édition des Universités de Montpellier II et Paris VII
\msp{page 61}%
les quatre conditions $\mathrm{A}_{W_0}$, $\mathrm{B}_{W_0}$, $\mathrm{C}_{W_0}$, $\mathrm{D}_{W_0}$ sont équivalentes, et le diagramme d'implications de \textup{1°)} (p. \textup{59}) se réduit donc au diagramme linéaire
\[
\mathrm{A} \implies \mathrm{B} \implies \mathrm{C} \implies \mathrm{D} \implies \mathrm{D}_{W_0}\ (\Longleftrightarrow \mathrm{A}_{W_0}\ldots),
\]
dont le premier terme est \textup{Loc (7 a) $=$ Loc (7)}, et le dernier équivaut à l'inclusion
\[
\mathcal{F}_{W_0}\subset\mathrm{Fib}_W,
\]
cas particulier de \textup{$\mathrm{A}=$ Loc (7 a)}, ou encore à l'inclusion (en apparence plus forte)
\[
\mathrm{Fibf}_{W_0}\subset\mathrm{Fib}_W,
\]
(cas particulier de la forme équivalente
\[
\mathrm{Fibf}_W\subset\mathrm{Fib}_W,
\]
de l'axiome} Loc (6 a)).

\noindent\cercle{4}\enspace\textit{Enfin, on a les implications plus techniques conditionnelles suivantes.}
\begin{enumerate}[leftmargin=2.5em]
\item[a)] \textit{Si $W_0$ satisfait à Loc (7 F), alors les conditions $\mathrm{A}_{W_0}$, $\mathrm{B}_{W_0}$, $\mathrm{C}_{W_0}$ sont équivalentes, donc le diagramme (4.4.13.5) se réduit au diagramme}
\[
\begin{tikzcd}[row sep=large, column sep=large]
\mathrm{A} \arrow[d,Rightarrow] & \\
\mathrm{B} \arrow[d,Rightarrow] & \\
\mathrm{C} \arrow[r,Rightarrow] \arrow[d,Rightarrow] & \mathrm{C}_{W_0} \arrow[d,Rightarrow] \\
\mathrm{D} \arrow[r,Rightarrow] & \mathrm{D}_{W_0}\ .
\end{tikzcd}
\]

\item[b)] \textit{Si $W$ est fortement saturé} (\textit{axiome} Loc (1 ter), \textit{qui}
\msp{page 62}%
\textit{est satisfait dans la grande majorité des cas, et est même impératif pour construire un prédérivateur satisfaisant à l'axiome de localisation} Der 2) \textit{et satisfait à} Loc (3 bis) (\textit{axiome de localisation fort,} cf. §1), \textit{alors on a} $\mathrm{B}\Longrightarrow\mathrm{A}$, \textit{donc}
\[
\mathrm{A} \iff \mathrm{B},
\]
i.e. Loc (7 a) $=$ Loc (7) \textit{équivaut à} Loc (7 b).
\end{enumerate}

\noindent\textsc{Démonstration}. Tout ce qui est énoncé a été vu dans les pages précédentes, à la seule exception de 4°, b. Mais cette implication $\mathrm{B}\Longrightarrow\mathrm{A}$ pour $W$ fortement saturé est prouvé dans XII (prop. 1, p. 54). Je rappelle que la démonstration utilise de façon essentielle le % [sic] prouvé -> prouvée, accord avec « implication »
% foot: Version 18 mars 2019
\printedpage{36}%
% head: A. Grothendieck : Dérivateurs, Ch. XVI, Édition des Universités de Montpellier II et Paris VII
formalisme du prédérivateur $\mathrm{HOT}_W$ (construit et étudié surtout dans VI), et la théorie des foncteurs $f_!$, et notamment le résultat (valable si $W$ satisfait à Loc (3 bis, 7 b)) que pour une flèche $f:S'\longrightarrow S$ dans $W$, la flèche correspondante
\[
f^*_{\mathrm{lc}}:\mathrm{HOT}^{\mathrm{lc}}_W(S)\longrightarrow\mathrm{HOT}^{\mathrm{lc}}_W(S')
\]
est une équivalence de catégories. Donc l'implication $\mathrm{B}\Longrightarrow\mathrm{A}$, elle aussi, est un résultat qui n'a rien d'évident.

\noindent\msp{page 63}%
4.4.14. Je résume les implications obtenues jusqu'à présent, y compris les implications conditionnelles, dans le diagramme récapitulatif suivant :
\[
\begin{tikzcd}[row sep=large, column sep=large]
|[draw,circle,inner sep=1pt]| \mathrm{A} \arrow[r,Rightarrow] \arrow[d,Rightarrow] & \mathrm{A}_{W_0} \arrow[d,Rightarrow] & \\
\mathrm{B} \arrow[r,Rightarrow] \arrow[d,Rightarrow] \arrow[u, dotted, bend left=40, "\text{$W$ (1 ter, 3 bis)}"] & \mathrm{B}_{W_0} \arrow[d,Rightarrow] & \\
\mathrm{C} \arrow[r,Rightarrow] \arrow[d,Rightarrow] & \mathrm{C}_{W_0} \arrow[d,Rightarrow] \arrow[uu, dotted, bend right=25, "\text{$W_0$ (7)}"'] & \\
\mathrm{D} \arrow[r,Rightarrow] \arrow[dr, dotted, "\text{$W$ (8 bis)}"'] & \mathrm{D}_{W_0} \arrow[r, Rightarrow, "{\text{\begin{tabular}{@{}l@{}}N.B. on verra sur l'exemple de $\mathbf{W}_0$\\ que cette implication est \textit{stricte,}\\ même sous l'hypothèse\\ Loc (1 ter, 2 bis, 4, 5, 6 a, 8, 9).\end{tabular}}}"] \arrow[d, dotted, "\text{$W$ (8 bis) $+$ $W_0$ (8 bis)}"] \arrow[uuu, dotted, bend right=50, "\text{$W_0$ (8 bis)}"'] &[12em] \mathrm{F} \\
& |[draw,circle,inner sep=1pt, label={right:{$\;,$}}]| \mathrm{A} &
\end{tikzcd}
\]
où j'ai cerclé la condition la plus forte et la plus importante de toutes, $\mathrm{A}=$ Loc (7), qui apparaît en deux extrémités opposées du diagramme.

Je voudrais encore faire un commentaire sur les ingrédients des résultats prouvés ici, et notamment de l'équivalence des huit conditions A à $\mathrm{D}_{W_0}$, quand $W$ et $W_0$ satisfont Loc (8 bis) :

\msp{page 64}%
\begin{enumerate}[leftmargin=2.5em]
\item[1)] La théorie des \textit{catégories de chemins} \underline{Ch}$(X)$ (VII), s'appuyant sur celle des \textit{morphismes lisses, propres, parfaits.}
\item[2)] Le lemme de K. S. \textsc{Brown} (3.2.1, page 20).
\item[3)] Pour l'implication $\mathrm{B}\implies\mathrm{A}$, \textit{i.e.} Loc (7 b) $\implies$ Loc (7) (sous l'hypothèse $W$ (1 ter, 3 bis)), j'utilise la théorie du prédérivateur $\mathrm{Hot}_W$ (VI), et plus particulièrement celle des foncteurs $f_!$ (définis sous la seule hypothèse $W$ (3 bis), renforçant la condition automatique $W$ (3) (faisant partie de la définition d'un localiseur fondamental)).
\end{enumerate}
% foot: Version 18 mars 2019

\leavevmode\printedpage{37}%
% head: A. Grothendieck : Dérivateurs, Ch. XVI, Édition des Universités de Montpellier II et Paris VII
Il est clair que je ne serais pas arrivé à y voir clair, sans le travail fondamental de K. S. \textsc{Brown}. Il reste pourtant encore des questions ouvertes. Ainsi, si on ne suppose pas que $W$ satisfasse à Loc (8 bis), mais au besoin à Loc (3 ter) (et je n'ai pas eu l'occasion encore de travailler avec un localiseur qui n'y satisfasse \ldots), est-il vrai néanmoins que l'on a encore l'implication
\[
\mathrm{D} \qquad \Longrightarrow \qquad \mathrm{A},
\]
de façon que les conditions A, B, C, D soient équivalentes (et en particulier Loc (7 c) implique déjà Loc (7) dans toute sa force) ? Il est clair dans tout cela que
\msp{page 65}%
je ne me suis pas suffisamment attaché à développer des exemples et des contre-exemples. Sans compter qu'il me semble que tous les localiseurs fondamentaux sur Cat qu'on rencontre “en pratique” satisfont déjà à Loc (8) et a fortiori à Loc (8 bis), donc pour eux on a bien l'implication escomptée. La même remarque s'applique à l'implication hypothétique
\[
\mathrm{D}_{W_0} \overset{?}{\implies} \mathrm{A},
\]
sous hypothèse Loc (8 bis) sur $W_0$ et sur $W$. Un candidat évident pour un contre-exemple déjà à la première question serait $W_\omega=$ plus petit localiseur fondamental sur Cat, un qu'il est presque évident qu'il ne satisfait \textit{pas} à Loc (8) ni même Loc (8 bis). Mais il n'y a aucune raison qu'il satisfasse non plus à F, de telle façon que (s'il n'y satisfait pas, donc non plus à D) l'implication $\mathrm{D}\Longrightarrow\mathrm{A}$ est bel et bien valable. Cela nous amène à la question, plus substantielle, si on n'aurait pas l'implication % [sic] un qu'il est -> dont il est (?), construction défectueuse
\msp{page 66}%
\[
\mathrm{F} \overset{?}{\implies} \mathrm{A}\ \text{\footnotemark[34]},
\]
\footnotetext[34]{C'est faux, \textit{cf.} plus bas, l'implication Loc (7 $\mathrm{D}_{W_\omega}$) $\Longrightarrow$ Loc (7 F) est \textit{stricte.}}%
du moins sous des conditions convenables assez anodines sur $W$, genre Loc (1 ter, 3 bis, 8 bis).

C'est sans doute là la question la plus substantielle et la plus intéressante qui est soulevée par le diagramme de la page 63, et qui ne fait plus appel à l'introduction d'un localiseur fondamental $W_0$ auxiliaire. Vu le lemme de Brown, il ne me semble pas absolument exclu qu'un tel miracle ait bel et bien lieu !

4.4.15 (Digression sur l'axiome Loc (6) des cofibrations.) Cela attire mon attention sur la question en quelque sorte duale : Si Cat a assez de $W$-\textit{co}fibrations, \textit{i.e.} si muni de $(W,\mathrm{Cof}_W)$ c'est une catégorie à cofibrations, cela implique-t-il déjà l'axiome Loc (6 a) (qui avait été dégagé, essentiellement, pour assurer qu'il en soit bien ainsi). La proposition 3.6. (p. 25) nous dit qu'il revient au même de prouver que, si $W$ satisfait Cof, alors les carrés de M. V. ouverts sont
\msp{page 67}%
$W$-cocartésien. Mais il me semble à présent que nous pouvons le prouver, sous la seule condition supplémentaire Loc (3 bis)~\footnote[35]{et peut-être, de plus, Loc (1 ter), \textit{i.e.} $W$ fortement saturée.}, laquelle assure l'existence des foncteurs $f_!$, et % [sic] cocartésien -> cocartésiens, accord avec « les carrés » ; [sic] (note 35) saturée -> saturé, accord avec « W » (localiseur)
% foot: Version 18 mars 2019
\printedpage{38}%
% head: A. Grothendieck : Dérivateurs, Ch. XVI, Édition des Universités de Montpellier II et Paris VII
en particulier la construction de Holim d'un diagramme $\fourche$ \quad comme $\int\fourche$ \quad .
Donc dire qu'un diagramme de $\mathcal{M}$ ouvert $\carre$ est $W$-cocartésien, doit revenir à dire que la flèche canonique
\[
\int\fourche \qquad \longrightarrow \qquad X
\]
est dans $W$ ; du moins si $W$ est supposé fortement saturé. Or cela a été bel et bien vu dans XII 3.4 (théorème 4 et son \textit{corollaire} 3, p. 149). Dans loc. cit. j'énonce les conclusions sous l'hypothèse que j'appelais W (7 bis) (l'actuel Loc (6)), mais je viens de vérifier que pour sa partie \textit{directe,} donc aussi pour ses corollaires 1 à 3, on n'a pas besoin
\msp{page 68}%
de cette hypothèse (comme je l'indique en annotations marginales).

4.16. Le résultat que je viens de prouver (sans vérifier avec soin la démonstration) me paraît aussi délicat, qu'il est satisfaisant. Cela rend la “conjecture” hésitante de 4.4.14, savoir que
\[
\mathrm{F} \overset{?}{\Longrightarrow} \mathrm{A} \qquad \text{donc} \qquad \mathrm{F}\Longleftrightarrow\mathrm{A}
\]
en supposant que $W$ satisfait Loc (1 ter, 3 bis, 8 bis), nettement plus engageante. Car on commence à s'habituer aux miracles les plus époustouflants \ldots

Pourtant, je viens de vérifier que cette conjecture est fausse, sur l'exemple de
\[
\mathbf{W}_0=\{f\in\mathrm{Fl}(\mathrm{Cat})\mid\pi_0(f)\ \mathrm{iso}\}\,,
\]
qui satisfait à F, mais pas à A, ni même à $\mathrm{D}_{W_\omega}$. Ainsi l'implication
\[
\mathrm{D}_{W_\omega} \implies \mathrm{F}
\]
et a fortiori les implications
\[
\left\{
\begin{aligned}
&\text{Loc (7 a)} \implies \mathrm{F},\ \text{Loc (7 b)} \implies \mathrm{F},\ \text{Loc (7 c)} \implies \mathrm{F}\\
&\mathrm{D}_{W_0} \implies \mathrm{F}
\end{aligned}
\right.
\]
sont \textit{strictes,} même en se limitant à des localiseurs fondamentaux $W$ satisfaisant (comme $\mathbf{W}_0$) à
\[
\text{Loc (1 ter, 2, 3 bis, 4, 5, 6 a, 8, 9)}\ \text{\footnotemark[36]}.
\]
\footnotetext[36]{\textbf{N.B. $\mathbf{W}_0$} ne satisfait \textit{pas} Loc (6 b), ni Loc (7 $\mathrm{D}_{W_\omega}$) (\textit{cf.} plus bas), par contre Loc (7 F).}
% foot: Version 18 mars 2019

% ---- from der-ch16_p041-045.tex
\printedpage{39}
% head: A. Grothendieck : Dérivateurs, Ch. XVI, Édition des Universités de Montpellier II et Paris VII

\msp{page 69}
Que $\mathbf{W}_0$ satisfasse à cet ensemble d'axiomes se vérifie par AQT \texttt{[âne qui trotte]}. (\textbf{N.B.} Qu'il satisfasse à Loc (1 ter, 2, 3 bis) est déjà conséquence du fait que $\mathbf{W}_0$ est de la forme $W_{\mathbf{D}}$, où $\mathbf{D}$ ($=\mathbf{D}_{\mathbf{W}_0}$ en l'occurrence) est un prédérivateur satisfaisant Der (3, 4, 5), ce qui est bien le cas pour $\mathbf{D}_{\mathbf{W}_0}$, qui est un vrai \textit{dérivateur.}) Il faut vérifier que tout $f$ dans Cat se factorise en $f = pi$, avec $i \in \mathbf{W}_0$, $p$ une $\mathbf{W}_0$-fibration. Or si $f : X \longrightarrow Y$, $I = \pi_0(X)$, on prend $\bar{X} = \coprod_{\alpha\in I} Y_{f(\alpha)}$, d'où une factorisation $f = pi$, où $\pi_0(i)$ est un isomorphisme, \textit{i.e.} $i \in \mathbf{W}_0$, et où $p : \bar{X} \longrightarrow Y$ est, au dessus de la composante connexe $Y_j$ de $Y$, donné par le produit $\bar{X}|Y_j \simeq Y_j \times \underbrace{\pi_0(f)^{-1}(j)}_{\text{catégorie discrète}}$ . Le morphisme $p$ est bien une $\mathbf{W}_0$-fibration, car grâce à Loc (4), il suffit de le vérifier sur chaque $Y_j$, et là c'est déduit d'une flèche $I_j \longrightarrow e$ (une $\mathbf{W}_0$-fibration) par changement de base $Y_j \longrightarrow e$.

D'autre part, que $\mathbf{W}_0$ ne satisfait \textit{pas} même à $\mathrm{D}_{W_\omega}$ a déjà été vu (I, p. 87),
\msp{page 70}
en réalisant l'application du $n$-revêtement $f : S^1 \longrightarrow S^1$ dans Cat par un morphisme $W_\omega$-parfait (en fait, par une Cat-fibration localement triviale), chose très facile\footnote[37]{Autre façon de le voir : on \textit{sait} (XII) que $\mathrm{W}$ (5, 7 $\mathrm{D}_W$) $\Longrightarrow W \subset \mathbf{W}_\infty$, or on a $\mathbf{W}_\infty \subset \mathbf{W}_0$, et l'inclusion est \textit{stricte\,!}}.

Je pense qu'à présent l'ensemble des différentes variantes de l'axiome des $W$-fibrations est à peu près compris, à cela près qu'il n'est toujours pas clair, en dehors de l'axiome Loc (8 bis), si les implications du diagramme
\[
\begin{tikzcd}[row sep=large, column sep=large]
\mathrm{A} \arrow[r, Rightarrow] \arrow[d, Rightarrow] & \mathrm{A}_{W_0} \arrow[d, Rightarrow] \\
\mathrm{B} \arrow[r, Rightarrow] \arrow[d, Rightarrow] & \mathrm{B}_{W_0} \arrow[d, Rightarrow] \\
\mathrm{C} \arrow[r, Rightarrow] \arrow[d, Rightarrow] & \mathrm{C}_{W_0} \arrow[d, Rightarrow] \\
\mathrm{D} \arrow[r, Rightarrow] & \mathrm{D}_{W_0}
\end{tikzcd}
\]
sont strictes ou non. Mais il me semble que tous les localiseurs fondamentaux intéressants rencontrés jusqu'à présent satisfont Loc (8), et même Loc (1 ter, 2, 3 bis, 4, 5, 6 a, 8, 9) \marginpar{?}(peut-être même Loc (7 F) - il faudrait que je teste sur les $\mathbf{W}_n$), donc la perplexité non résolu me paraît sans grande conséquence. % [sic] résolu -> résolue, accord avec « la perplexité » ; « ? » imprimé dans la marge gauche à la hauteur de cette ligne

\msp{page 71}

\printedpage{40}
% head: A. Grothendieck : Dérivateurs, Ch. XVI, Édition des Universités de Montpellier II et Paris VII

\section*{5\quad Axiome des limites Loc (8), et l'axiome Loc (9) d'accessibilité}

C'est l'axiome :

\textbf{Loc (8)}\quad $W$ est stable par $\varinjlim$ filtrantes.

Une variante affaiblie utile est la suivante :

\textbf{Loc (8 bis)}\quad $W$ est stable par $\varinjlim$ dénombrables.

Explicitement, si on a deux systèmes inductifs filtrants dans Cat
\[
\begin{aligned}
\underline{X} &= (X_i)_{i\in I}\\
\underline{Y} &= (Y_i)_{i\in I},
\end{aligned}
\]
et une flèche
\[
\underline{f} = (f_i)_{i\in I} : \underline{X} \longrightarrow \underline{Y},
\]
donnée par des $f_i : X_i \longrightarrow Y_i$, posant $X = \varinjlim X_i$, $Y = \varinjlim Y_i$, $f = \varinjlim f_i : X \longrightarrow Y$, Loc (8) nous dit que si les $f_i \in W$, alors $f \in W$. Ici $I$ est
\msp{page 72}
une (petite) catégorie filtrante quelconque, et quitte à la remplacer par une catégorie cofinale, on peut supposer $I$ un ensemble ordonné. L'axiome Loc (8 bis) correspond au cas où $I = \mathbf{N}$.

Notons que grâce à la construction de la catégorie des chemins, nous avons une factorisation, fonctorielle en $f$, de toute flèche $f : X \longrightarrow Y$ de Cat en $f = pi$ :
\[
X \xrightarrow{\quad i = i_f\quad} \bar{X}(f) = \bar{X} \xrightarrow{\quad p = p_f\quad} Y,
\]
avec $i \in W_\omega$, $p \in \Parf_{W_\omega}$, a fortiori $i \in W$, $p \in \Parf_W$. Ce foncteur $f \longmapsto (\bar{X}(f), p_f)$ commute aux limites inductives. Ainsi, si on applique ceci à chacun des $f_\alpha : X_\alpha \longrightarrow Y_\alpha$, le factorisant en
\[
X_\alpha \xrightarrow{i_\alpha} \bar{X}_\alpha \xrightarrow{p_\alpha} Y_\alpha,
\]
on trouve un système inductif de diagrammes de type $\boldsymbol{\Delta}^2$ (factorisant les $f_i$), et en passant à la $\varinjlim$ on trouve
\[
X \xrightarrow{i} \bar{X} \xrightarrow{p} Y,
\]
qui n'est autre que la factorisation canonique
\msp{page 73}
de $f = \varinjlim f_i$ en $pi$, avec $i \in W_\omega$, $p \in \Parf_{W_\omega}$ (N.B. Alors que ni $W_\omega$, ni $\Parf_{W_\omega}$ ne sont stables par $\varinjlim$ filtrantes, à coup sûr\,!). On a donc $i \in W$, et la condition $f \in W$, \textit{i.e.} $pi \in W$, équivaut donc à $p \in W$. Cela nous ramène, pour vérifier l'axiome Loc (8) (resp. Loc (8 bis)) au cas où les $f_i$ sont dans $\Parf_{W_\omega}$.

Supposons maintenant que $W$ satisfasse à $\mathrm{D}_{W_\omega}$ (la plus faible de toutes les variantes envisagées ici de l'axiome des $W$-fibrations Loc (7), à la seule exception de F). Alors on a
\[
p \in W \qquad\quad \Longleftrightarrow \qquad\quad \text{les fibres de $p$ sont $W$-asphériques,}
\]
\printedpage{41}
% head: A. Grothendieck : Dérivateurs, Ch. XVI, Édition des Universités de Montpellier II et Paris VII
l'implication $\Longrightarrow$ résultant de l'hypothèse $\mathrm{D}_{W_\omega}$ $\Parf_{W_\omega} \cap W \subset W^{\mathrm{u}}$, et $\Longleftarrow$ résultant du sorite des morphismes $W$-lisses \textit{etc.}, via Loc (3). De même, l'hypothèse que les $f_i \in W$ équivaut à celle que les fibres des $p_i$ sont $W$-asphériques. D'autre part, il est immédiat que
\msp{page 74}
toute fibre de $p$ est isomorphe à la $\varinjlim\limits_{I}$ des fibres correspondantes des $p_\alpha$. De cette façon, on trouve :

\textbf{Proposition 5.1.} \textit{Supposons que $W$ satisfasse l'axiome des $W$-fibrations} Loc (7 $\mathrm{D}_{W_\omega}$) (\textit{ce qui est donc le cas s'il satisfait un des axiomes plus forts} Loc (7 a \textit{ou} b \textit{ou} c)). \textit{Pour que $W$ satisfasse à} Loc (8) (resp. \textit{à} Loc (8 bis)), \textit{il faut et il suffit que l'ensemble des objets dans} Cat \textit{qui sont $W$-asphériques soit stable par $\varinjlim$ filtrante} (resp. \textit{par $\varinjlim$ de suites dénombrables}).

5.2. Notons que la catégorie Cat est un paratopos, a fortiori elle est \textit{accessible} (\textit{cf.} XVIII, §8), et il en est donc de même de $\uFl(\mathrm{Cat})$. Ainsi on a dans $\mathrm{Fl}(\mathrm{Cat}) = \mathrm{Ob}\,\uFl(\mathrm{Cat})$ la notion de \textit{partie accessible.} On dira donc qu'une partie $W$ de $\mathrm{Fl}\,\mathrm{Cat}$ (plus généralement d'une catégorie $\mathrm{Fl}(\mathcal{M})$, où $\mathcal{M}$ est une catégorie accessible) est \textit{accessible,} si elle l'est par rapport à la catégorie ambiante
\msp{page 75}
$\uFl\,\mathrm{Cat}$. Ceci posé, l'axiome Loc (9) est le suivant.

\textbf{Loc (9)}\quad $W$ est une partie \textit{accessible} de $\mathrm{Fl}(\mathrm{Cat}) = \mathrm{Ob}\,\uFl(\mathrm{Cat})$.

5.3. Rappelons que cette notion un peu technique se simplifie considérablement, si on suppose que $W$ satisfait l'axiome de limites Loc (8). Alors $W$ est accessible si et seulement si il existe une \textit{petite} partie $W_0$ de Fl(Cat), telle que $W$ soit la ``L$_1$-enveloppe'' de $W_0$, \textit{i.e.} la plus petite partie de Fl(Cat) contenant $W_0$ et L$_1$-stable (\textit{i.e.} stable par $\varinjlim$ filtrantes). S'il en est ainsi, il est clair qu'on peut prendre $W_0$ de la forme
\[
W_\pi = W \cap \mathrm{Fl}_\pi(\mathrm{Cat}),
\]
où $\pi$ est un cardinal infini convenable, et où $\mathrm{Fl}_\pi(\mathrm{Cat}) \overset{\text{déf}}{=} \mathrm{Fl}(\uCat_\pi)$, $\uCat_\pi$ désignant la sous-catégorie strictement pleine de Cat formée des objets $X$ tels que $\mathrm{card}\,X \leq \pi$ (où par définition, $\mathrm{card}\,X = \mathrm{card}\,\mathrm{Fl}\,X$).

\msp{page 76}
Notons que la filtration qu'on vient de décrire sur Cat (indexée par les petits cardinaux infinis) n'est autre que sa filtration canonique, \textit{i.e.} $\uCat_\pi$ est la sous-catégorie pleine de Cat formée des objets $\pi$-accessibles de Cat, et de même dans $\uFl(\mathrm{Cat})$. On trouve ainsi, grâce aux développements de XVIII, §13, la caractérisation suivante, que je rappelle ici pour mémoire (en me bornant pour simplifier au cas où $W$ satisfait à Loc (8)) :

\textbf{Proposition 5.3.1.}\footnote[38]{\textit{cf. loc. cit.} 13.36.} \textit{Soit $W$ une partie de} Fl(Cat) \textit{stable par $\varinjlim$ filtrantes, et soient pour tout cardinal infini $\pi$, $W_\pi = W \cap \mathrm{Fl}_\pi(\mathrm{Cat})$, et $\underline{W}$, $\underline{W}_\pi$ les sous-catégories pleines correspondantes de} $\uFl(\mathrm{Cat})$. \textit{On dit que $W$ est} $\pi$-accessible \textit{si $W$ est la L-enveloppe de $W_\pi$ (donc $W$ est accessible si et seulement s'il existe un cardinal infini $\pi$ tel que $W$ soit $\pi$-accessible). Ceci posé, les conditions suivantes sont équivalentes :}

\printedpage{42}
% head: A. Grothendieck : Dérivateurs, Ch. XVI, Édition des Universités de Montpellier II et Paris VII

\begin{enumerate}[label={}]
\item[a)] \textit{$W$ est $\pi$-accessible.}
\item[b)] \textit{Pour tout système inductif $(f_i)_{i\in I}$,}
\msp{page 77}
\textit{$L_1$-$\pi$-adapté dans} $\uFl\,\mathrm{Cat}$ (cf. \textit{ci-dessous}), \textit{de limite inductive $f$, on a $f \in W$ si et seulement si la sous-catégorie pleine $I_W$ des $i \in \mathit{Ob}I$ tels que $f_i \in W$ est cofinale dans $I$.}
\item[b$'$)] (\textit{Quand on a choisi, pour tout $f \in$} Fl(Cat)\textit{, une représentation $L_1$-$\pi$-adaptée de $f$ comme $\varinjlim$}) : \textit{Si $f \in W$, alors la sous-catégorie pleine $I_W$ définie dans} b) \textit{est cofinale dans $I$.}
\end{enumerate}

Rappelons ici la notion de système inductif L$_1$-$\pi$-adapté. Il s'agit des conditions :
\begin{enumerate}[label={}]
\item[1°)] $I$ est grand devant $\pi$, et stable par $\varinjlim$ filtrantes de cardinal $\leq \pi$.
\item[2°)] Le foncteur $i \longmapsto x_i : I \longrightarrow \uFl(\mathrm{Cat})$ commute aux $\varinjlim$ en question.
\item[3°)] Les $x_i$ sont dans $\uFl(\mathrm{Cat})_\pi \overset{\text{déf}}{=} \uFl_\pi(\mathrm{Cat})$.
\end{enumerate}
Notons que dans le cas d'un ``paratopos parabélien'' tel que Cat ou $\uFl\,\mathrm{Cat}$, il y a une façon canonique, $\pi$ étant donné, de représenter un objet $f$ comme $\varinjlim$ d'un système inductif L$_1$-$\pi$-adapté. Ici,
\[
f : X \longrightarrow Y
\]
étant donné, on considère l'ensemble $I$ des
\msp{page 78}
couples $(X_i \subset X, Y_i \subset Y)$ de sous-catégories de $X$ et de $Y$, telles que l'on ait $X_i, Y_i \in \mathrm{Ob}\,\mathrm{Cat}_\pi$ (\textit{i.e.} $\mathrm{card}\,(X_i), \mathrm{card}\,(Y_i) \leq \pi$) et que $f(X_i) \subset Y_i$. Il est immédiat que $I$ est filtrant, pour la relation d'ordre évidente
\[
(X_i, Y_i) \subset (X_j, Y_j) \qquad\quad \Longleftrightarrow \qquad\quad X_i \subset Y_i \ \text{ et } \ X_j \subset Y_j, % [sic] X_i \subset Y_i et X_j \subset Y_j -> X_i \subset X_j et Y_i \subset Y_j, la relation d'ordre compare les couples (X_i,Y_i) et (X_j,Y_j)
\]
et même grand devant $\pi$, que $I$ est stable par $\varinjlim$ filtrantes de cardinal $\leq \pi$, et que le foncteur
\[
i \longmapsto (f_i : X_i \longrightarrow Y_i) : I \longrightarrow \uFl\,\mathrm{Cat}
\]
commute auxdites $\varinjlim$ , enfin les $f_i$ sont dans $\uFl(\mathrm{Cat})_\pi$ par construction. Ainsi on a un système inductif L$_1$-$\pi$-adapté, et il est clair que sa limite inductive est canoniquement isomorphe à $f$. Ceci vu :

\textbf{Corollaire 5.3.2.} \textit{La condition que $W$ soit $\pi$-accessible signifie donc que pour tout $f : X \longrightarrow Y \in W$, et pour toute sous-catégorie $X_i$ de $X$ et $Y_i$ de $Y$, de cardinal $\leq \pi$ (avec $f(X_i) \subset Y_i$), mais en fait c'est inutile, quitte à agrandir $Y_i$), il existe des sous-catégories} % [sic] (avec f(X_i) \subset Y_i), mais ... agrandir Y_i) -> (avec f(X_i) \subset Y_i, mais ... agrandir Y_i), parenthèse fermante surnuméraire
\msp{page 79}
\textit{$X_j \supset X_i$ et $Y_j \supset Y_i$ de cardinal $\leq \pi$, telles que $f$ induise $f_j : X_j \longrightarrow Y_j$ et que l'on ait $f_j \in W$.}

La notion d'accessibilité de $W$ me semble aussi bien comprise que celle de limites, je n'y vois plus aucun mystère. D'autre part, il me semble que tous les localiseurs fondamentaux utilisés à ce jour y satisfont (tout comme à l'axiome Loc (8)). Ici, je vais me borner à établir :

\printedpage{43}
% head: A. Grothendieck : Dérivateurs, Ch. XVI, Édition des Universités de Montpellier II et Paris VII

\textbf{Théorème 5.4.}\footnote[39]{Mérite bien le nom de \textit{théorème,} surtout si on y inclut toute la liste des axiomes Loc mentionnés ci-contre\,! Ce théorème n'a rien d'évident, à partir de la définition cohomologique de $\mathbf{W}_\infty$ - c'est un résultat (``bien connu'') qui me paraît profond.} \textit{Le localiseur fondamental $\mathbf{W}_\infty$ (équivalences faibles ``habituelles'') satisfait à l'axiome} Loc (9) (\textit{en fait, il satisfait à} tous \textit{les axiomes vu jusqu'à présent, savoir} Loc (1 ter, 2, 3 bis, 4, 5, 6, 7, 8, 9)). \textit{Plus précisément, $\mathbf{W}_\infty$ est $\aleph_0$-accessible (où $\aleph_0$, ou $\pi_0$, est le premier cardinal infini).} % [sic] les axiomes vu -> les axiomes vus, accord du participe

Je ne reviens pas ici sur la vérification des autres axiomes, et me borne à prouver Loc (9), l'accessibilité, en utilisant ce qui peut m'être utile. Je me
\msp{page 80}
borne à indiquer la démonstration (assez technique) dans ses grandes lignes. Pour le moment, $W$ est à nouveau un localiseur fondamental quelconque, sauf qu'il satisfait à Loc (8).

\textbf{Lemme 5.5.}\footnote[40]{On suppose que $W$ satisfasse à Loc (8).}
\begin{enumerate}[label={}]
\item[a)] \textit{Pour que le critère} 5.3.1 b) \textit{de $\pi$-accessibilité de $W$ soit satisfait, il faut et il suffit qu'il le soit lorsque les $f_i$ sont $W_\omega$-parfaits.}
\item[b)] \textit{Supposons que $\mathit{Parf}_{W_\omega} \cap W \subset W^u$,} i.e. \textit{que $W$ satisfait à} Loc (7 $\mathrm{D}_{W_\omega}$) (\textit{la plus faible des variantes de l'axiome des $W$-fibrations, à l'exception seulement de} Loc (7 F)). \textit{Supposons de plus (pour être tranquille) que $W$ satisfait} \texttt{[phrase non terminée]}. \texttt{[?Alors]} \textit{$W$ est $\pi$-accessible si et seulement si l'ensemble} $\mathrm{Asph}_W \subset \mathrm{Ob}\,\mathrm{Cat}$ \textit{des objets $W$-asphériques de} Cat \textit{est $\pi$-accessible. Ou encore, que la partie} $\mathrm{Asph}_W \subset \mathrm{Ob}\,\mathrm{Cat}_{\textit{0-connexe}}$ \textit{soit $\pi$-accessible, où} $\mathrm{Cat}_{\textit{0-connexe}}$ \textit{désigne la sous-catégorie pleine de} Cat \textit{formée des objets} 0-\textit{connexes.}
\end{enumerate}
Pour l'essentiel, la démonstration est la même que celle de 5.1, en un peu plus technique. L'outil essentiel est toujours la factorisation fonctorielle $f = pi$ d'une flèche de Cat, avec $i \in W_\omega$, $p \in \Parf_{W_\omega}$, et la propriété de
\msp{page 81}
commutation aux $\varinjlim$ filtrantes de ladite factorisation.

Revenons maintenant au cas de $W = \mathbf{W}_\infty$. On doit prouver que $\mathrm{Asph}_{\mathbf{W}_\infty} \subset \mathrm{Ob}\,\mathrm{Cat}_{\text{0-connexe}}$ est $\pi_0$-accessible. Donc à prouver le critère correspondant à 5.3.1 b), en se plaçant maintenant dans $\mathrm{Cat}_{\text{0-connexe}}$ au lieu de $\mathrm{Fl}\,\mathrm{Cat}$. On a donc
\[
X = \varinjlim_{i\in I} X_i,
\]
représentation L$_1$-$\pi$-admissible de $X$ $\mathbf{W}_\infty$-asphérique, à prouver que pour tout $i$ dans $I$, $\exists j$ dans $I$ et $i \longrightarrow j$, avec $X_j$ $\mathbf{W}_\infty$-asphérique.

J'admets ici le lemme AQT \texttt{[âne qui trotte]} :

\textbf{Lemme 5.6.} \textit{Soit $\mathbf{W}_0 = \{f \in \mathrm{Fl}\,\mathrm{Cat} \mid \pi_0(f) \ \ \textit{isomorphisme}\}$. Alors $\mathbf{W}_0$ est $\pi_0$-accessible (et $L_1$-stable). A fortiori, l'ensemble des objets} 0-\textit{connexes de} Cat \textit{est $\pi_0$-accessible.}

Ceci nous ramène, quitte à remplacer $I$ par une catégorie cofinale, à supposer que les $X_i$ sont 0-connexes. (Mais j'ai oublié que cette réduction est déjà faite\,! Donc il faut voir 5.6 comme ingrédient

% ---- from der-ch16_p046-050.tex
\printedpage{44}
% head: A. Grothendieck : Dérivateurs, Ch. XVI, Édition des Universités de Montpellier II et Paris VII
% foot: Version 18 mars 2019

\msp{page 82}
de la démonstration de 5.5 b).)

Soit $x$ un objet (point) de $X$, $x_i$ son image dans $X_i$. Considérons (via ma construction standard des catégories de chemins) les $\Omega^1(X,x)$, $\Omega^1(X_i,x_i)$. On trouve
\[ \Omega^1(X,x) \simeq \varinjlim_{I} \Omega^1(X_i,x_i), \]
et on vérifie que c'est là encore un système inductif $\mathrm{L}_1$-$\pi_0$-adapté. En itérant, on trouve pour $n \in \mathbf{N}^*$
\[ \Omega^n(X,x) \simeq \varinjlim_{I} \Omega^n(X_i,x_i), \]
limite inductive $\mathrm{L}_1$-$\pi_0$-adaptée. Utilisant 5.6, on en conclut que pour tout $n$, le sous-ensemble (supposons $I$ un ensemble ordonné, que diable ça suffit !)
\[ I_n = \{ i \in I \mid \Omega^n(X_i,x_i) \ \text{ est 0-connexe} \} \]
est cofinal dans $I$, de plus il est stable par $\varinjlim$ filtrantes de cardinal $\leq \pi_0$. Donc % [sic] $\leq \pi_0$ -> $\leq \aleph_0$ (?), a cardinal bound is apparently intended
\[ I_\infty \overset{\text{déf}}{=} \bigcap I_n \]
est aussi cofinal dans $I$ (\textit{cf.} XVIII, 13.10.2, p. 302).

\msp{page 83}
Or $i \in I_\infty$ signifie que tous les $\Omega^n(X_i,x_i)$ sont 0-connexes, ce qui signifie aussi que $X_i$ est $\mathbf{W}_\infty$-asphérique (puisqu'il est déjà 0-connexe) (\textit{cf.} XII, cor. p. 190). Cela achève la démonstration du théorème 5.4.

\section*{6\quad La condition Loc (5, 7 D$_{W_\omega}$) implique $W \subset \mathbf{W}_\infty$}

Je vais redonner ici, pour mémoire, la démonstration du théorème 1 de XII :

\textbf{Théorème 6.1.} \textit{Soit $W$ un localiseur fondamental satisfaisant à} D$_{W_\omega}$ \textit{(où $W_\omega$ est le plus petit de tous les localiseurs fondamentaux). Alors on a $W = \mathrm{Fl}\,\mathrm{Cat}$,} ou \textit{$W \subset \mathbf{W}_\infty$ (suivant que $W$ ne satisfait pas à l'axiome de connexité} Loc (5), \textit{ou qu'il y satisfait.})

Je vais faire la démonstration en supposant que $W$ satisfait à Loc (5). Pour la démonstration du fait que dans le cas contraire, on a $W = \mathrm{Fl}(\mathrm{Cat})$, je renvoie à XII, th. 2, p. 20. (Voir aussi le résultat très voisin, plus général et plus faible, rappelé dans prop. 2.7, page 12.)

\textbf{Lemme 6.2.} \textit{Soient $W$, $W'$ deux localiseurs fondamentaux, $W$ satisfaisant à} Loc (7 D$_{W_\omega}$). \textit{Pour qu'on ait $W \subset W'$, il faut et il suffit qu'on ait} $\mathrm{Asph}_W \subset \mathrm{Asph}_{W'}$.

\msp{page 84}
Il suffit bien sûr de prouver “il suffit”, donc supposons que tout objet $W$-asphérique est $W'$-asphérique, et prouvons que $W \subset W'$. Soit donc $f : X \longrightarrow Y$ dans $W$, prouvons $f \in W'$. Pour ceci, j'utilise la factorisation canonique $f = pi$,
\[ X \xrightarrow{\ i \in W_\omega\ } \bar{X} \xrightarrow{\ p \in \mathrm{Parf}_{W_\omega}\ } Y. \]
\printedpage{45}
% head: A. Grothendieck : Dérivateurs, Ch. XVI, Édition des Universités de Montpellier II et Paris VII
% foot: Version 18 mars 2019
On a donc aussi $i \in W$, et comme $f = pi \in W$, on a $p \in W$, donc $p \in \mathrm{Parf}_{W_\omega} \cap W$. Or D$_{W_\omega}$ signifie que cet ensemble est contenu dans $W^{\mathrm{u}}$, donc on a $p \in W^{\mathrm{u}}$, et en particulier les fibres de $p$ sont $W$-asphériques. Elles sont donc $W'$-asphériques, et comme $p$ est $W_\omega$-parfait et a fortiori $W'$-parfait, il s'ensuit que $p \in W'$. On a aussi $i \in W_\omega \subset W'$, donc $i \in W'$, donc aussi $f = pi \in W'$, q.e.d.

Revenons alors à la démonstration de 6.1. En vertu du lemme, il suffit de prouver que si $X$ dans Cat est $W$-asphérique, il est $\mathbf{W}_\infty$-asphérique. Mais par hypothèse Loc (5) on sait déjà que $X$ est 0-connexe, et choisissant un point $x$ dans $X$, qu'il soit $\mathbf{W}_\infty$-asphérique s'exprime par la condition
\msp{page 85}
que les $\Omega^n(X,x)$ (pour $n \geq 1$) sont 0-connexes (\textit{cf.} XII, cor. p. 190 déjà cité). Pour ceci, à cause de Loc (5) il suffit de prouver qu'ils sont $W$-asphériques, et par induction il suffit de prouver que $\Omega^1(X,x)$ l'est. Or par définition de $\Omega^1(X,x)$, c'est là la fibre en $x$ d'une flèche dans Cat
\[ E = \underline{\mathrm{Ch}}(X,x,-) \xrightarrow{\ p \in \mathrm{Parf}_{W_\omega}\ } X, \]
où la source $E$ est $W_\omega$-asphérique, a fortiori $W$-asphérique. Comme $X$ l'est également, on a donc $p \in W$, donc $p \in \mathrm{Parf}_{W_\omega} \cap W \subset W^{\mathrm{u}}$, donc $p$ est à fibres $W$-asphériques, et en particulier $\Omega^1(X,x)$ est $W$-asphérique, q.e.d.

\textbf{Corollaire 6.3.} (Caractérisation de $\mathbf{W}_\infty$.) \textit{$\mathbf{W}_\infty$ est le plus grand localiseur fondamental satisfaisant à} Loc (5, 7 D$_{W_\omega}$) \textit{- donc aussi le plus grand localiseur fondamental satisfaisant à} Loc (5, 7).

Je pense que nous trouverons plus bas une caractérisation “interne” de $\mathbf{W}_\infty$, via un “axiome de descente” (très fort) Loc (10),
\msp{page 86}
comme l'unique localiseur fondamental satisfaisant Loc (5, 7 D$_{W_\omega}$, 10) (ou encore Loc (5, 7, 10)).

6.4. \textit{Commentaire sur les axiomes des $W$-fibrations du types} Loc (7 \ldots), \textit{et notamment sur} Loc (7 D$_{W_\omega}$). \textit{Relation conjecturale avec les} “types d'homotopie simples” \textit{de} \textsc{Whitehead}. % [sic] du types -> du type (or des types), grammatical agreement

a) Il me semble que \textit{tous} les localiseurs fondamentaux avec lesquels on a travaillé tacitement jusqu'à présent (les $\mathbf{W}_n$ ($n \in \mathbf{N}^*$), équivalence homologique et cohomologique, équivalence rationnelle \ldots) correspondent à des localiseurs $W \supset \mathbf{W}_\infty$. Par le théorème 6.1, il s'ensuit donc qu'\textit{aucun} d'eux, à part $\mathbf{W}_\infty$ lui-même, ne satisfait à la plus faible version D$_{W_\omega}$ de toutes de l'axiome Loc (7 \ldots), à la seule exception de Loc (7 F). Ainsi, les localiseurs $W$ qui y satisfont sont à l'exception de D$_{W_\omega}$ lui-même, d'un type tout nouveau - des localiseurs qui sont strictement contenus dans $\mathbf{W}_\infty$.

b) Ceci soulève, bien sûr, la question de l'existence de localiseurs fondamentaux contenus dans $\mathbf{W}_\infty$ et
\msp{page 87}
distincts de $\mathbf{W}_\infty$. Je n'ai pas même prouvé que $W_\omega \neq \mathbf{W}_\infty$, alors que je n'ai pas le moindre doute que tel est bien le cas. \texttt{[Cisinski a démontré que} $W_\omega = \mathbf{W}_\infty$\texttt{.]} A fortiori, j'ignore si le plus petit localiseur fondamental satisfaisant Loc (7 D$_{W_\omega}$) est différent de $\mathbf{W}_\infty$. Je serais tout particulièrement intéressé par le plus petit localiseur $\tilde{\mathbf{W}}_\infty \subset \mathbf{W}_\infty$ satisfaisant tous les axiomes de stabilité Loc (1 ter, 2, 3 bis, 4, 6, 7, 8) (\textbf{N.B.} il satisfait
\printedpage{46}
% head: A. Grothendieck : Dérivateurs, Ch. XVI, Édition des Universités de Montpellier II et Paris VII
% foot: Version 18 mars 2019
nécessairement Loc (5) car $\tilde{\mathbf{W}}_\infty \subset \mathbf{W}_\infty \subset \mathbf{W}_0$, de plus Loc (7) équivaut, en présence de Loc (8), à la forme très faible Loc (7 D$_{W_\omega}$) : $\mathrm{Parf}_{W_\omega} \cap W \subset W^{\mathrm{u}}$. Enfin, on montre que Loc (4) est conséquence des autres axiomes, \textit{cf.} plus bas.) Je conjecture à présent que $\tilde{\mathbf{W}}_\infty$ (qui donne lieu à un formalisme homotopique tout aussi complet que la sempiternelle “équivalence faible”) est bel et bien différent de $\mathbf{W}_\infty$ \texttt{[non par Cisinski]}. Et je suspecte
\msp{page 88}
que la catégorie $\mathrm{Hot}_{\tilde{\mathbf{W}}_\infty}$ mérite d'être appelée catégorie des \textit{“types d'homotopie simples”} (à la \textsc{Whitehead})\footnote[41]{Canulé (?), \textit{cf.} commentaire page 92.}. Un premier test décisif, c'est de vérifier que l'invariant cohomologique principal de \textsc{Whitehead} d'un type d'homotopie simple, savoir un élément dans un $K(\mathbf{Z}[\pi_1])$ (si je me rappelle bien\footnote[42]{non, c'est un invariant plus fin \ldots}) - ou quelque chose d'approchant, peut se définir déjà en termes des objets dans $\mathrm{Hot}_{\bullet\,\tilde{\mathbf{W}}_\infty}$ (le Hot pointé). Et deuxio, si on associe à une “maquette simpliciale” finie (décrivant combinatoirement une triangulation finie) l'ensemble ordonné de ses simplexes, vu comme un objet de Cat, et l'objet correspondant dans $\mathrm{Hot}_{\tilde{\mathbf{W}}_\infty}$, celui-ci ne change pas\footnote[43]{C'est évident même pour $W_\omega$, puisque le type d'homotopie \textit{stricte ne change pas !} - Non, il y a ici confusion entre homotopie topologique et homotopie combinatoire, \textit{cf.} p. 92 \ldots}, à isomorphisme canonique près, quand on effectue sur ce “polyèdre” envisagé une opération élé\msp{page 89}mentaire à la \textsc{Whitehead}. Cela suggère d'ailleurs d'essayer de définir de telles opérations élémentaires dans Cat tout entier, ou tout au moins dans la sous-catégorie des (petits) ensembles ordonnés, dans l'espoir de parvenir (peut-être) à une construction plus directe du localiseur pertinent pour le type d'homotopie simple.

c) Je termine par un commentaire sur le sens de l'axiome Loc (7 D$_{W_\omega}$), par opposition avec l'axiome plus faible Loc (7 F) qui assure seulement l'existence d'une “structure à fibrations” $(W, \mathrm{Fib}_W)$ sur Cat. Je signale aussi en passant que Loc (7 F), tout comme Loc (5) et comme Loc (9), et à la différence de tous les autres axiomes passés en revue jusqu'à présent, n'est \textit{pas} un axiome de stabilité : il affirme, non que certaines flèches (sous des hypothèses convenables) sont dans $W$, mais l'existence d'une factorisation. Aussi il n'est
\msp{page 90}
nullement clair que l'intersection de deux localiseurs fondamentaux satisfaisant à Loc (7 F) y satisfasse.

Loc (7 D$_{W_\omega}$) est exactement ce qu'il faut pour assurer que la factorisation canonique d'une flèche $f : X \longrightarrow Y$ de Cat en $f = pi$, avec $i \in W_\omega$, $p \in \mathrm{Parf}_{W_\omega}$, est aussi une factorisation du type stipulé par l'axiome de factorisation Loc (7 F) : on a (bien sûr) $i \in W$, mais aussi $p \in \mathrm{Fib}_W$ (car D$_{W_\omega}$ équivaut à sa forme forte $\mathrm{Parf}_{W_\omega} \subset \mathrm{Fib}_W$). Cela signifie aussi que la factorisation canonique de $f$ peut être utilisée pour “calculer” ses fibres homotopiques, non seulement en les points de $Y$, mais aussi par rapport à un morphisme de changement de base $Y' \longrightarrow Y$.

En l'absence de cet axiome,
\msp{page 91}
\printedpage{47}
% head: A. Grothendieck : Dérivateurs, Ch. XVI, Édition des Universités de Montpellier II et Paris VII
% foot: Version 18 mars 2019
admettant seulement Loc (7 F), on sait qu'il existe une théorie des fibres $W$-homotopiques, mais on n'a a priori aucun moyen pratique pour “calculer” lesdites fibres, faute de disposer d'un critère maniable suffisant pour assurer $p \in \mathrm{Fib}_W$. J'ai donc l'impression que l'axiome Loc (7 F) va être d'un maigre usage. (Il est vrai que l'exemple de $\mathbf{W}_0$ devrait me rendre prudent ! Il me faudrait regarder le cas des $\mathbf{W}_n$, s'ils satisfont Loc (7 F) \ldots). À mon sens, l'axiome Loc (7 D$_{W_\omega}$) est l'exact “minimum vital” pour être à l'aise dans le maniement des fibres $W$-homotopiques. Il est remarquable que ceci implique qu'on doive avoir $W \subset \mathbf{W}_\infty$, et que cela exclut du même coup \textit{toutes} les variantes considérées à ce jour de la notion de type d'homotopie !

\msp{page 92}
6.5. Finalement, la relation entrevue avec les “types d'homotopie simples” de \textsc{Whitehead} s'évanouit\footnote[44]{Mais peut-être quand même pas - la question reste ouverte.} : en effet, ces derniers sont strictement plus fin que les types d'homotopie (\textit{stricts}) ordinaires, et l'invariant de Whitehead n'est \textit{pas} un invariant du type d'homotopie ordinaire. Il est vrai que lorsqu'on considère le squelette simplicial d'un polyèdre triangulé comme objet de Cat, l'homotopie (continue) ordinaire entre deux tels polyèdres s'interprète, dans Cat, non pas (me semble-t-il) par l'existence d'un homotopisme entre les objets correspondants de Cat, mais par leur isomorphie dans $\mathrm{Hot}_{\mathbf{W}_\infty} \overset{\text{déf}}{=} \mathrm{Cat}(\mathbf{W}_\infty)^{-1}$. Mais le procédé d'approximations simpliciales de \textsc{Brouwer} montre que si les polyèdres $X$ et $Y$ sont homotopes, alors il existe des subdivisions simpliciales $X'$ et $Y'$ de $X$ et de $Y$ qui sont “combinatoirement homotopes” (dans un sens convenable). Donc je présume que les maquettes simpliciales de $X'$ et de $Y'$ sont alors homotopes dans Cat. Il s'ensuivrait que % [sic] plus fin -> plus fins, agreement with "ces derniers"
\msp{page 93}
$X$ et $Y$, j'entends leurs maquettes, sont bien homotopes dans Cat, s'il est vrai qu'une subdivision simpliciale (il suffirait la subdivision barycentrique, en itérant \ldots) $X'$ de $X$ a une maquette dans Cat qui lui est homotope. La question soulevée ici me paraît intéressante, il faut visiblement en trouver la réponse. Elle se pose déjà pour tout ensemble ordonnée $I$, qu'on va supposer fini pour simplifier, en prenant l'ensemble (ordonné par inclusion) $I' = \mathrm{sd}(I)$ des parties totalement ordonnées de $I$. On sait que les réalisations topologiques de $I$ et de $\mathrm{sd}(I)$ sont canoniquement isomorphes, mais $I$ et $\mathrm{sd}(I)$ sont-ils homotopes dans Cat ? On a bien une application croissante canonique % [sic] ensemble ordonnée -> ensemble ordonné, agreement
\[ f : \mathrm{sd}(I) \longrightarrow I, \]
qui associe à toute partie totalement ordonnée de $I$ son plus grand élément. Mais existe-t-il une application croissante en sens inverse, qui soit un homotopisme inverse du précédent à homotopie près ? Je viens de vérifier dans le cas

\msp{page 94} manque

\msp{page 95}
du bord du 2-simplexe

\begin{center}
\begin{tikzpicture}[baseline=(b.base)]
\node (b) at (0.95,0.3) {};
\draw (0,0) -- (1.9,0) -- (0.95,0.95) -- cycle;
\fill (0,0) circle (2.2pt);
\fill (1.9,0) circle (2.2pt);
\fill (0.95,0.95) circle (2.2pt);
\end{tikzpicture}
\qquad\qquad\qquad ou \qquad\qquad\qquad
\begin{tikzpicture}[baseline=(c.base)]
\node (c) at (0,-0.1) {};
\draw (0,0) circle (0.48);
\fill (90:0.48) circle (2.2pt);
\fill (210:0.48) circle (2.2pt);
\fill (330:0.48) circle (2.2pt);
\end{tikzpicture}
\qquad\quad ,
\end{center}

subdivision barycentrique
\printedpage{48}
% head: A. Grothendieck : Dérivateurs, Ch. XVI, Édition des Universités de Montpellier II et Paris VII
% foot: Version 18 mars 2019

\begin{center}
\begin{tikzpicture}[baseline=(b.base)]
\node (b) at (0.95,0.3) {};
\draw (0,0) -- (1.9,0) -- (0.95,0.95) -- cycle;
\fill (0,0) circle (2.2pt);
\fill (0.95,0) circle (2.2pt);
\fill (1.9,0) circle (2.2pt);
\fill (0.475,0.475) circle (2.2pt);
\fill (1.425,0.475) circle (2.2pt);
\fill (0.95,0.95) circle (2.2pt);
\end{tikzpicture}
\qquad\qquad\qquad ou \qquad\qquad\qquad
\begin{tikzpicture}[baseline=(c.base)]
\node (c) at (0,-0.1) {};
\draw (0,0) circle (0.48);
\fill (30:0.48) circle (2.2pt);
\fill (90:0.48) circle (2.2pt);
\fill (150:0.48) circle (2.2pt);
\fill (210:0.48) circle (2.2pt);
\fill (270:0.48) circle (2.2pt);
\fill (330:0.48) circle (2.2pt);
\end{tikzpicture}
\end{center}

(hexagone), que pour toute application croissante $g : I \longrightarrow \mathrm{sd}(I)$, sa réalisation topologique est homotope à une application constante, donc $g$ ne peut être un homotopisme dans Cat (car cela impliquerait que sa réalisation topologique est un homotopisme).

Ceci montre donc que la notion d'homotopisme dans Cat est \textit{strictement plus fine} que celle pour les réalisations topologiques, ce qui redonne à nouveau substance à l'espoir que l'invariant de Whitehead pourrait se définir comme invariant d'homotopie sur Cat, ou du moins sur une sous-catégorie pleine convenable, telle celle des ensembles ordonnés (éventuellement restreints à être finis), ou des maquettes simpliciales. Et même, que l'ensemble des Wh-types d'homotopie (types d'homotopie simples) peut se
\msp{page 96}
plonger canoniquement dans un $\mathrm{Hot}_W$ convenable, où $W$ serait un localiseur fondamental strictement plus fin que $\mathbf{W}_\infty$, p.ex. (qui sait ?) $\tilde{\mathbf{W}}_\infty$. Le premier travail qu'on aurait envie de faire, c'est de construire un tel $W$ \textit{ad hoc}, pour que l'invariant de Whitehead puisse s'écrire comme un foncteur sur $\mathrm{Hot}_{\bullet W}$. Pour commencer, il faudrait déjà s'assurer que cet invariant peut être vu comme un invariant du type d'homotopie, \textit{i.e.} qu'il s'interprète comme un foncteur sur $\mathrm{Hot}_{\bullet}\mathrm{Cat} \simeq \mathrm{Cat}(\,\underbrace{\text{homotop}}_{\substack{\text{ensemble des}\\ \text{homotopismes}}}\,)^{-1}$. Ce travail suggérera peut-être aussi la définition des “invariants de Whitehead supérieurs” (impliquant les $\pi_i$, non seulement le $\pi_1$), lesquels pourraient peut-être servir (avec celui de dimension 1) à décrire le localiseur hypothétique $W$ (“de Whitehead”).

\msp{page 97}

\section*{7\quad Commentaires sur les axiomes}

7.1. La principale raison d'être d'un localiseur fondamental $W$, c'est de définir une théorie d'homotopie’, et plus précisément un prédérivateur $\mathrm{HOT}_W$, dont on espère que ce sera même un dérivateur, ou qu'il satisfasse du moins à certains des axiomes des dérivateurs. Les axiomes Loc (1) à Loc (9) sont en majeure partie orientés vers cet objectif, à cela près que Loc (1, 2, 3) constitue un “minimum vital” absolu, pour pouvoir commencer à travailler avec $W$. Les principales questions qui se posent me paraissent les suivantes : % [sic] d'homotopie’ -> d'homotopie (or “théorie d'homotopie”), closing quote without opening quote
\begin{enumerate}[label={}, leftmargin=2.5em]
\item[1)] Quand Cat munie de $(W, \mathrm{Cof}_W)$, resp. de $(W, \mathrm{Fib}_W)$, est-elle une catégorie à cofibrations resp. à fibrations, de façon que les résultats de K. S. \textsc{Brown} soient applicables ?
\item[2)] Sous quelles conditions, sur $W$ et sur $f$, $f_!$ et $f_*$ existent-ils ? Et quand a-t-on, de plus, exactitude des carrés cocartésiens resp. cartésiens ?

\msp{page 98}
\item[3)] Plus généralement, en termes des axiomes imposés à $W$, quelles sont les axiomes Der (\ldots) que sont vérifiés pour le prédérivateur $\mathrm{HOT}_W$ ? (Ce qui pourrait aussi dépendre du choix du domaine Diag $\subset$ Cat. Le prédérivateur est défini a priori sur Cat tout entier). En particulier, % [sic] quelles sont les axiomes -> quels sont les axiomes, agreement; [sic] que sont vérifiés -> qui sont vérifiés, relative pronoun
\end{enumerate}

% ---- from der-ch16_p051-055.tex
\printedpage{49}
% head: A. Grothendieck : Dérivateurs, Ch. XVI, Édition des Universités de Montpellier II et Paris VII
% foot: Version 18 mars 2019

\begin{enumerate}[label=4)]
\item quand $\mathrm{HOT}_W$ est-il bel et bien un dérivateur sur Cat tout entier, \textit{i.e.} quand satisfait-il à tous les axiomes ?
\end{enumerate}

Une question qui me paraît moins importante (notamment que 1°) est la suivante :

\begin{enumerate}[label=1 bis)]
\item Quand $W$ s'insère-t-il dans un triple de Quillen clos $(W,C,F)$ (propre de préférence).
\end{enumerate}

Je vais passer en revue ces questions une à une.

7.2. On a vu (mais sans que j'aie fait une vérification soigneuse pour la réciproque) que $(W,\mathrm{Cof}_W)$ définit une structure à cofibrations si l'axiome \textit{Loc (6 a)} est satisfait, et que cette condition est aussi nécessaire si Loc (3 bis) est satisfait\footnote[45]{Vérifier quand même si Loc (3 bis) n'est pas déjà utilisé dans la partie directe ! \textbf{O.K.}}.

\msp{page 99}
Alors les immersions ouvertes de Dwyer (ou au choix les immersions fermées de Dwyer) sont dans $\mathrm{Cof}_W$, et la construction standard du mapping-cone définit, pour toute flèche $f$ dans Cat, une factorisation \textit{fonctorielle}
\[
\begin{aligned}
f = pi \quad & i \text{ immersion ouverte de Dwyer (donc dans } \mathrm{Cof}_W\text{),}\\
& p \in W.
\end{aligned}
\]
Cela, plus la construction du foncteur $f_!$, et en particulier des $W$-sommes amalgamées
\begin{tikzcd}[row sep=tiny, column sep=small, baseline=(\tikzcdmatrixname-2-1.base)]
 & X_1 \\
\displaystyle\int X_0 \arrow[ur] \arrow[dr] & \\
 & X_2
\end{tikzcd}, données dans VI (théorie de $\mathrm{HOT}_W$) et XII, nous donne une maîtrise parfaite (il me semble) sur la “structure homotopique gauche” définie par $W$, sous le seul axiome Loc (6 a) (Axiome du cube, forme directe). Je rappelle que les foncteurs $f_!$ sont construits dans VI sous la seule hypothèse \textit{Loc (3 bis)} ($f$ flèche quelconque de Cat), mais sauf erreur, la construction
\msp{page 100}
et l'interprétation $W$-homotopique de
\begin{tikzcd}[row sep=tiny, column sep=small, baseline=(\tikzcdmatrixname-2-1.base)]
 & X_1 \\
\displaystyle\int X_0 \arrow[ur] \arrow[dr] & \\
 & X_2
\end{tikzcd}
donnée dans XII ne dépend pas même de cette hypothèse\footnote[46]{\textit{cf.} XII, page 264} - mais je ne peux le garantir. Il me semble que XII (et notamment le long §5 sur les sommes amalgamées) est écrit sous l'hypothèse que W (1, 2, 3), \textit{i.e.} Loc (1, 2, 3bis), sont satisfaits.

7.3. C'est une chose assez extraordinaire que les foncteurs $f_!$ pour $\mathrm{HOT}_W$ existent, sous la seule hypothèse \textit{Loc (3 bis)} (axiome de localisation, forme forte) - et ils se “calculent” d'une façon féeriquement simple ! Je suis à peu près sûr que sous cette seule hypothèse, \textit{l'axiome d'exactitude pour les carrés cocartésiens} n'est \textit{pas} satisfait\footnote[47]{Il faudrait un exemple !}. Je ne vois aucun moyen de le vérifier, sauf via la théorie de \textsc{Brown}, sous l'hypothèse \textit{Loc (6 a)} (\textit{cf.} 7.2).

Mais sous la seule condition Loc (3 bis) le calcul standard des fibres d'un $f_!(\xi)$
\printedpage{50}
% head: A. Grothendieck : Dérivateurs, Ch. XVI, Édition des Universités de Montpellier II et Paris VII
% foot: Version 18 mars 2019
\msp{page 101}
est valable, en d'autre termes le prédérivateur $\mathrm{HOT}_W$ satisfait Der (4$^*$, 5$^*$) - et de plus Der (6$^*$) (exactitude “à gauche”) si $W$ satisfait Loc (6 a). % [sic] d'autre termes -> d'autres termes, missing plural s

7.4. Venons-en à la question si Cat muni de $(W,\mathrm{Fib}_W)$ est une catégorie à fibrations. C'est l'axiome \textit{Loc (7 F),} mais je ne connais guère de moyen de le vérifier (en dehors du cas très, très particulier $W = \mathbf{W}_0$) que moyennant tout au moins l'axiome \textit{Loc (7 $D_{W_\omega}$),} ou l'une ou l'autre des variantes plus fortes de Loc (7). De Loc (7 $\mathrm{D}_{W_\omega}$) il a été abondamment question précédemment, \textit{cf.} notamment 6.4 (pages 89-91). En mon sens, cet axiome assure une maîtrise comparable sur la structure $W$-homotopique droite (fibres $W$-homotopiques, suites exactes de cosuspension) que Loc (6 a) pour la
\msp{page 102}
structure gauche. En particulier, on a la factorisation \textit{fonctorielle} de toute flèche $f$ de Cat en
\[
f = pi, \qquad i \in W_\omega \subset W,\ p \in \mathrm{Parf}_{W_\omega} \subset \mathrm{Fib}_W .
\]
Mais \textit{l'axiome Loc (7 $D_{W_\omega}$) me paraît de nature considérablement plus forte que Loc (6 a),} lequel (si je ne m'abuse) est vérifié pour tous les localiseurs qui ont été envisagés jusqu'à présent. (Mais j'avoue n'avoir pas regardé le cas des $\mathbf{W}_n$, $n \in \mathbf{N}$, sauf $\mathbf{W}_0$.) Par contre, l'axiome Loc (7 $\mathrm{D}_{W_\omega}$) implique $W \subset \mathbf{W}_\infty$ (théorème 6.1), relation vraiment draconienne (\textit{cf.} commentaire 6.4 a).

Par ailleurs, le seul axiome Loc (7 $\mathrm{D}_{W_\omega}$), pour autant que je sache, n'implique pas l'existence de $f_*$ pour \textit{toute} flèche dans Cat, mais seulement pour des
\msp{page 103}
flèches de la forme $f \times \mathrm{id}_X : S \times X \longrightarrow T \times X$, où $f : S \longrightarrow T$ est un morphisme entre ensembles ordonnés finis. Ou plus généralement, entre ensembles ordonnés $S$, $T$ tels que $S^\circ$ soit artinien, \textit{i.e.} $\mathrm{S}$ noethérien, et que $\mathrm{Esp}(S^\circ)$ soit noethérien (cette dernière hypothèse étant inutile sous l'axiome Loc (8) des limites, \textit{cf.} XVII, th. 3.1, p. 50). Donc ce n'est pas le résultat rêvé.

7.5. Le seul cas que je connaisse où $f_*$ existe pour \textit{toute} flèche dans Cat, est le cas où $W$ satisfait à
\[
\textit{Loc (6 a, 8, 9)}\ \ \text{\footnotemark[48]\,\footnotemark[49]}, \tag*{$(*)$} % [sic] marks printed (48) (49) -> the two footnotes printed below are numbered 49 and 50, numbering mismatch between marks and footnotes
\]
\footnotetext[49]{Vérifier s'il ne faut pas aussi la condition Loc (3 bis) en plus de Loc (6 a, 8, 9). Vérification faite, c'est O.K. (sans Loc (3 bis)).}
\footnotetext[50]{\textbf{N.B.} Après nouvelle réflexion, il apparaît qu'il suffit même Loc (6$'$ a, 8, 9).}
c'est à dire l'axiome des $W$-cofibrations faibles, déjà abondamment mentionné et enfin ceux des limites et d'accessibilité, \textit{i.e.} $W$ est $\mathrm{L}_1$-accessible dans \underline{Fl}(Cat). Le résultat-clef pertinent est donné dans XV 8.4 (p. 105-107), qu'il faut joindre à XV 7.4 (p. 98-99) pour en faire
\msp{page 104}
bon usage, et en déduire le théorème 10.1 (XV, p. 134-135). Mieux que l'existence des $f_*$, l'ensemble d'axiomes étonnamment faible que je viens d'expliciter suffit à assurer que $\mathrm{HOT}_W$ est un \textit{dérivateur} sur Cat tout entier, en plein sens du terme, \textit{i.e.} qu'il satisfait à Der (1, 2, 3, 4, 5, 6, 4$^*$, 5$^*$, 6$^*$). Ceci est obtenu en remplaçant (grâce à XV 8.4) Cat par
\printedpage{51}
% head: A. Grothendieck : Dérivateurs, Ch. XVI, Édition des Universités de Montpellier II et Paris VII
% foot: Version 18 mars 2019
$A^\wedge$, où $A$ est une catégorie-test convenable, $W$ par $W_A \subset \mathrm{Fl}\, A^\wedge$, et en prouvant que $W_A$ s'insère dans un triple de Quillen clos $(W,C,F)$, avec $C = \mathrm{Mono}(A^\wedge)$, ce qui détermine $F$ comme $(C \cap W_A)_*$. On applique alors XV 7.4 à cette situation (résultats résumés dans XIX, th. 2.3 (p. 19), 2.6 (p. 25), 2.9 (p. 28). On retrouve l'existence des $f_!$ (avec un mode de calcul bien différent de celui dans Cat !), et en plus on trouve l'existence de $f_*$ (plus l'exactitude dans cor. [?] 5
\msp{page 105}
des carrés $W$-cartésiens) - ce qui pour moi, pendant un bon moment, était \textit{le} gros problème pour la théorie de $\mathrm{HOT}_W$ !).

Je trouve époustouflant qu'on parvienne à un résultat aussi fort avec des axiomes aussi ridiculement faibles. Plus particulièrement, qu'on ait toutes les propriétés désirables de la structure homotopique droite, sans avoir à faire sur $W$ aucune hypothèse du type $W$-fibration ! Ces axiomes n'impliquent \textit{pas} l'axiome Loc (7) (équivalent à Loc (7 $\mathrm{D}_{W_\omega}$) en présence de Loc (8)), comme on a vu sur l'exemple de $\mathbf{W}_0$. Impliquent-ils tout au moins Loc (7 F), \textit{i.e.} l'existence d'assez de $W$-fibrations ? Je ne hasarderais aucune conjecture à ce sujet ! Mais il faut bien se dire qu'en l'absence de Loc (7) (et même
\msp{page 106}
en admettant Loc (7 F)), tout en sachant que la structure homotopique droite associée à $W$ est parfaite à tous égards, on n'a a priori aucune maîtrise sur celle-ci, notamment pour le calcul des fibres $W$-homotopiques et la construction des carrés $W$-cartésiens. Même en ramenant toutes ces constructions à $A^\wedge$ (prenant, si on y tient, $A = \boldsymbol{\Delta}$, \textit{i.e.} se ramenant au contexte semisimplicial), quoi qu'on sache qu'il y a une factorisation \textit{fonctorielle} $f = pi$ avec $i \in \mathrm{TC} \subset W_A$, $p \in F \subset \mathrm{Fib}_{W_A}$, je ne vois aucune façon concrète de l'expliciter. À moins que la considération des chemins infinis, qui (en l'absence de Loc (7)) s'avère inopérante dans Cat sauve la mise dans $\boldsymbol{\Delta}^\wedge$ ? ?

\msp{page 107}
Question qui devrait être facile à élucider, mais sur laquelle je ne veux pas m'attarder à présent.

Notons en passant que le résultat de saturation forte de \textsc{Quillen}, pour les triples de Quillen clos, montre que l'ensemble d'axiomes $(*)$ implique l'axiome Loc (1 ter), \textit{i.e.} que $W$ est \textit{fortement saturé.}

Notons aussi que contrairement à ce qui a lieu pour les $f_!$, je ne sais prouver l'existence de $f_*$, et notamment des suites de cosuspension, que sous des conditions (soit Loc (7 $\mathrm{D}_{W_\omega}$), soit Loc (6a, 8, 9)) qui assurent déjà que les carrés $W$-cartésiens sont exacts, et par suite les suites de cosuspension exactes.

Je pense que dans les pages précédentes, j'ai dit tout ce que je sais dire de pertinent, concernant les questions 1, 2, 4 soulevées dans 7.1. Ces réponses me parais\msp{page 108}sent satisfaisantes - plus satisfaisantes que je ne l'aurais osé espérer pendant les premiers deux mois de mes cogitations ! Il reste les questions 3 et 1 bis, que je vais à présent examiner.

\printedpage{52}
% head: A. Grothendieck : Dérivateurs, Ch. XVI, Édition des Universités de Montpellier II et Paris VII
% foot: Version 18 mars 2019
7.6. Sur les différents axiomes des dérivateurs : sous quelles conditions sur $W$ un de ces axiomes est-il satisfait pour $\mathrm{HOT}_W$\footnote[51]{Pour les axiomes Der (1) à Der (5$^*$), \textit{cf.} I, p. 67-68. Les axiomes Der (6, 6$^*$) sont les axiomes d'exactitude, dont la nécessité m'est apparue seulement ultérieurement.} ? Pour ce qui est des axiomes Der (4, 5, 6, 4$^*$, 5$^*$, 6$^*$), je crois avoir dit tout ce que je sais en dire, et avoir des critères ce validité parfaitement maniables, dans les pages précédentes. Résumons pourtant : % [sic] ce validité -> de validité, misprint
\[
\left\{
\begin{array}{lcl}
\text{structure à cofibrations} & : & \text{Loc (6) (+ Loc (3 bis) ?)} \\[1ex]
\text{structure à fibrations} & : & \text{Loc (7 }\mathrm{D}_{W_\omega}\text{)} \\[1ex]
\text{Der (4}^*\text{, 5}^*\text{)} & : & \text{Loc (3 bis)} \\[1ex]
\text{Der (4}^*\text{, 5}^*\text{, 6}^*\text{)} & : & \text{Loc (3 bis, 6a) ou Loc (6}'\text{ a, 8, 9)} \\[1ex]
\text{Der (4, 5, 6) (sur Cat) (et tout le reste !)} & : & \text{Loc (6}'\text{ a, 8, 9)} \\[1ex]
\text{Der (4, 5, 6) (sur Diag $=$ ens. ordonnés finis)} & : & \text{Loc (7 }\mathrm{D}_{W_\omega}\text{) (ou seulement Loc (7 F)).}
\end{array}
\right.
\]

Il reste les axiomes Der (1) (décomposition), Der (2) (localisation), Der (3) (objet unité), que je vais passer à présent en revue. On vérifie
\msp{page 109}
de façon quasi-tautologique que l'axiome Der 1 est satisfait sans condition sur $W$ (pas même Loc (2, 3) !) - en fait, il est satisfait pour tout prédérivateur $\mathbf{D}_{(\mathcal{M},W)}$ défini par une catégorie de modèles.

L'axiome de localisation Der (2) est par contre de nature délicate. Il est valable, plus généralement, pour une catégorie de modèles $(\mathcal{M},W)$ (définissant un prédérivateur $\mathbf{D} = \mathbf{D}_{(\mathcal{M},W)}$) et un objet $S$ de Cat, lorsque $(\mathcal{M}(S),W(S))$ est associé à une structure à fibrations ou à cofibrations, ce qui par K. S. \textsc{Brown} implique que toute flèche de $\mathrm{HO}_W(S) = \mathcal{M}(S)W(S)^{-1}$ se “réalise” par une flèche de $\mathcal{M}(S)$ - et alors le résultat est évident. La condition que je viens de dire est satisfaite lorsque $(\mathcal{M},W)$ lui-même est associé à une structure de cofibrations ou de fibrations,
\msp{page 110}
et si de plus on suppose que l'on est dans l'un de deux cas suivants : % [sic] l'un de deux cas -> l'un des deux cas, apparent misprint
\begin{enumerate}[label=\alph*)]
\item $S$ est un ensemble ordonné fini (ce qui nous amène à nous restreindre, à défaut de mieux, à Diag $=$ sous-catégorie de Cat formée des ensembles ordonnés finis).
\item La factorisation stipulée dans l'axiome de factorisation $f = pi$ (avec $i \in \mathrm{Cof}$, $p \in W$ (cas de cofibrations) ou $i \in W$, $p \in \mathrm{Fib}$ (cas de fibrations)) peut être choisie fonctorielle en $f$.
\end{enumerate}
Or, dans le cas de Cat, dans les cas où on sait prouver que $W$ est associé à une structure à cofibrations (axiome \textit{Loc (6 a)}) où à fibrations (axiome \textit{Loc (7 $D_{W_\omega}$)}), on trouve bel et bien une factorisation fonctorielle. Donc chacun de ces axiomes implique Der (2). % [sic] où à fibrations -> ou à fibrations, accent misprint

\msp{page 111}
Reste l'axiome de contractibilité : si $S$ dans Cat est contractile, est-il $\mathrm{HOT}_W$-asphérique, \textit{i.e.} le foncteur
\[
\mathrm{HOT}_W(e) = \mathrm{Hot}_W \longrightarrow \mathrm{HOT}_W(S)
\]
\printedpage{53}
% head: A. Grothendieck : Dérivateurs, Ch. XVI, Édition des Universités de Montpellier II et Paris VII
% foot: Version 18 mars 2019
est-il pleinement fidèle ? Quand on dispose de la théorie de $f_!$ “version Cat” (donc sous le seul axiome \textit{Loc (3 bis)}), on trouve ceci (noté dans VI, p. 202)\footnote[52]{moyennant Loc (3 bis).} :
\begin{quote}
Soit $f : S \longrightarrow T$ une flèche dans Cat. Pour que ce soit une $\mathrm{HOT}_W$-équivalence, \textit{i.e.} pour que $f^* : \mathrm{HOT}_W(T) \longrightarrow \mathrm{HOT}_W(S)$ induise un foncteur \textit{pleinement fidèle} sur la sous-catégorie des objets constants de $\mathrm{HOT}_W(T)$, il faut et il suffit que $\mathrm{hot}_W(f) : \mathrm{hot}_W(S) \longrightarrow \mathrm{hot}_W(T)$ soit un isomorphisme. Donc il \textit{suffit} que $f \in W$, et cette condition est aussi nécessaire si $W$ est fortement saturé (axiome Loc (1 ter)).
\end{quote}
Par acquit de conscience, j'explicite ici la démonstration. L'existence des foncteurs
\msp{page 112}
$f_!$, en particulier des objets d'homologie $H_\bullet^{\mathrm{HOT}_W}(S,\xi)$ (où simplement $H_\bullet^W(S,\xi)$) pour $\xi \in \mathrm{HOT}_W(e) = \mathrm{Hot}_W$, permet d'interpréter homologiquement la condition de $\mathrm{HOT}_W$-équivalence, comme signifiant que pour tout $\xi = \mathrm{hot}_W(X)$ dans $\mathrm{Hot}_W$, la flèche canonique dans $\mathrm{Hot}_W$ % [sic] (où simplement -> (ou simplement, accent misprint
\[
H_\bullet^W(S,\xi) \longrightarrow H_\bullet^W(T,\xi) \tag*{$(*)$}
\]
est un isomorphisme. Or le calcul des $f_!$ (dans le cas de $p : S \longrightarrow e$ et de $T \longrightarrow e$) nous montre que les deux membres de $(*)$ ne sont autres (à isomorphisme près) que $\mathrm{hot}_W(X \times S)$ et $\mathrm{hot}_W(X \times T)$, la flèche étant induite par
\[
\mathrm{id}_X \times f : X \times S \longrightarrow X \times T. \tag*{$(*)$}
\]
Donc la condition envisagée signifie que cette flèche est un isomorphisme dans $\mathrm{Hot}_W$ pour tout $X$ dans Cat. Prenant $X = e$, on voit que cela implique que $\mathrm{hot}_W(f)$ est un isomorphisme, et cela suffit, vu que le
\msp{page 113}
foncteur $\mathrm{hot}_W : \mathrm{Cat} \longrightarrow \mathrm{Hot}_W$ commute aux produits finis (il suffit de Loc (1, 2, 3)) \texttt{[plutôt Loc (1, 2, 3 bis)]} et que la flèche $(*)$ dans $\mathrm{Hot}_W$ s'identifie à
\[
\mathrm{id}_{\mathrm{hot}_W(X)} \times \mathrm{hot}_W(f)\,.
\]
Cela montre donc que si $S$ est un objet $W$-asphérique, \textit{i.e.} tel que $S \longrightarrow e$ soit dans $W$ (et a fortiori, si $S$ est contractile), alors $S \longrightarrow e$ est $\mathrm{HOT}_W$-asphérique.

Pour résumer, on trouve le tableau
\[
\left\{
\begin{array}{lcl}
\text{Der (1)} & : & \text{sans hypothèse sur } W \\[1ex]
\text{Der (2)} & : & \text{Loc (6 a) \textit{ou} Loc (7 }\mathrm{D}_{W_\omega}\text{)} \\[1ex]
\text{Der (3)} & : & \text{Loc (3 bis) \textit{ou} Loc (6}'\text{ a, 8, 9)\ \ \footnotemark[53].}
\end{array}
\right.
\]
\footnotetext[53]{\textbf{N.B.} On verra plus bas que Loc (6 a, 8, 9) implique déjà Loc (3 bis).}

7.7. Il faut seulement justifier encore que Loc (6$'$ a, 8, 9) implique Der (3). Mais alors, on l'a vu, $\mathrm{HOT}_W$ peut être défini par un triple de Quillen très spécial. Mais même pour un triple de Quillen spécial sans plus $(\mathcal{M},W)$, le prédérivateur $\mathbf{D}$ qu'il définit satisfait (en plus de Der (4$^*$, 5$^*$, 6$^*$)) aux conditions Der (1) (trivialement), Der (2)

% ---- from der-ch16_p056-059.tex
\printedpage{54}
% head: A. Grothendieck : Dérivateurs, Ch. XVI, Édition des Universités de Montpellier II et Paris VII
\msp{page 114}
(\textit{cf.} p. 110, cas b)) et \textit{de plus} Der (3). (Donc dans le cas d'un triple de Quillen \textit{très} spécial, le prédérivateur obtenu est bel et bien un dérivateur.) Je vais expliciter ici la démonstration.

Je vais désigner par $(\mathcal{M},\Sigma)$ la catégorie de modèles donnée ($\Sigma\subset\operatorname{Fl}\mathcal{M}$), par $\mathbf{D}=\mathbf{D}_{\mathcal{M},W}$ le prédérivateur correspondant, et par
\[ W=W_D \] % [sic] W_D -> W_{\mathbf{D}}, index printed as italic D, elsewhere bold D
le localiseur dans Cat associé, formé des ``$\mathbf{D}$-équivalences''. On voit tout de suite que $W$ est fortement saturé (vrai pour tout $W_{\mathbf{D}}$, $\mathbf{D}$ un prédérivateur), \textit{i.e.} satisfait Loc (1 ter), a fortiori Loc (1), d'autre part il satisfait à Loc (3 bis$^*$), comme conséquence du fait que $\mathbf{D}$ satisfait à Loc (4$^*$, 5$^*$). A fortiori il satisfait à Loc (3$^*$), et pour voir que tout $X$ contractile dans Cat est $W$-asphérique, \textit{i.e.} $X\longrightarrow e$ dans $W$, on sait (\textit{cf.} §1)
\msp{page 115}
qu'il suffit de prouver qu'il en est ainsi si $X$ a un objet initial (axiome Loc (2$^*$) dual de l'axiome Loc (2)). À ce moment, on saura que $W$ est un localiseur fondamental, et de plus il satisfait Loc (3$^*$ bis), donc aussi Loc (3 bis) (forme forte de l'axiome de localisation).

Donc \textit{on peut supposer} que si $S$ a un objet initial, alors $S$ est $W$-asphérique, \textit{i.e.} $S\longrightarrow e$ est dans $W$, \textit{i.e.} $p^*:\mathbf{D}(e)\longrightarrow\mathbf{D}(S)$ est pleinement fidèle. Il revient au même de dire que la flèche d'adjonction
\[ p_!p^*\xi\longrightarrow\xi \]
est un isomorphisme, pour tout $\xi$ dans $\mathbf{D}(e)$. Or le calcul de $p_!(\eta)$ (pour $\eta\in\operatorname{Ob}\mathbf{D}(S)$, \textit{i.e.} $\eta\in\operatorname{Ob}\mathcal{M}(S)$) se fait ainsi : on choisit une flèche
\[ \eta'\xrightarrow{\Sigma_S}\eta \]
avec $\eta'$ ``cofibrant'' (pour la structure à cofibrations pertinente sur $\mathcal{M}(S)$), et on
\msp{page 116}
pose \footnote[54]{opération $p_!$ au niveau de $\mathcal{M}(S)$.}
\[ p_!(\operatorname{ho}_{W,S}(\eta))=\operatorname{ho}_W(p_!(\eta')). \]
Quand $\eta=p^*(\xi)$, la flèche d'adjonction pertinente dans $\operatorname{Ho}_\Sigma=\mathbf{D}(e)$ s'obtient par composition
\[ p_!(\eta')\longrightarrow p_!(p^*(\xi))\xrightarrow{\text{adj.}}\xi, \tag*{$(*)$} \]
où ici $p_!$ est encore pris dans le sens ordinaire ($\mathcal{M}(S)\longrightarrow\mathcal{M}(e)=\mathcal{M}$). Or $S$ ayant un objet initial 0, on a
\[ \begin{aligned}
p_!(\eta') &\ \simeq\ \eta'(0)\\
p_!(p^*(\xi)) &\ \simeq\ p^*(\xi)(0)\ \simeq\ \xi,
\end{aligned} \]
et la flèche d'adjonction dans $(*)$ est l'identité, la première flèche est la flèche $\eta'(0)\longrightarrow\eta(0)$ induite par $\eta'\longrightarrow\eta$, elle est dans $\Sigma$, donc $(*)$ est dans $\Sigma$, donc c'est un isomorphisme dans $\operatorname{Ho}_\Sigma=\mathbf{D}(e)$, q.e.d.

Le même argument marche d'ailleurs si $\mathcal{M}$ est seulement une catégorie à cofibrations ou à fibrations, mais en se bornant alors au cas où $S$ est un ensemble ordonné fini.
% foot: Version 18 mars 2019

\newpage
\printedpage{55}
% head: A. Grothendieck : Dérivateurs, Ch. XVI, Édition des Universités de Montpellier II et Paris VII
\msp{page 117}
Je vais dégager à part ce qu'on vient d'obtenir :

\textbf{Proposition 7.7.1.} \textit{Soit} $(\mathcal{M},W)$ \textit{une catégorie de modèles, associée à un triple de Quillen clos} spécial $(W,C,F)$\textit{. Alors le prédérivateur} $\mathbf{D}=\mathbf{D}_{\mathcal{M},W}$ \textit{sur} Cat \textit{satisfait aux axiomes} Der (1, 2, 3, 4$^*$, 5$^*$, 6$^*$)\textit{. Si de plus le triple de Quillen est} très \textit{spécial} (\textit{donc en particulier} $C=\operatorname{Mono}(\mathcal{M})$), \textit{alors il satisfait de plus à} Der (4, 5, 6), i.e. \textit{c'est un dérivateur au plein sens du terme. D'autre part, le localiseur} $W=W_{\mathbf{D}}$ \textit{dans} Cat \textit{associé à} $\mathbf{D}$ (\textit{dans le cas d'un triple de Quillen spécial sans plus}) \textit{est un} localiseur fondamental, \textit{satisfaisant de plus} Loc (1 ter, 3 bis).

\textbf{N.B.} Ce qui est dit ici pour le cas d'un triple de Quillen très spécial (qui ne fait pas appel à la théorie des localiseurs
\msp{page 118}
fondamentaux dans Cat) montre à nouveau que si $W$ localiseur fondamental dans Cat satisfait à Loc (6$'$ a, 8, 9), alors $\mathrm{HOT}_W$ est un vrai dérivateur. On a vu de plus dans 7.6 (page 111) que l'on a alors $W=W_{\mathrm{HOT}_W}$ \footnote[55]{Mais c'était sous l'hypothèse Loc (3 bis) ! Gare au cercle vicieux.}, donc il s'ensuit :

\textbf{Corollaire 7.7.2 :} \textit{On a l'implication}
\[ \text{Loc } (6'\text{ a, 8, 9})\ \Longrightarrow\ \text{Loc(1 ter, 3 bis)}\ \text{\footnotemark[56]}. \]
\footnotetext[56]{Loc (3 bis) pas prouvé, probablement faux.}

Notons aussi, pour mémoire :

\textbf{Corollaire 7.7.3.} \textit{Soit} $(\mathcal{M},W)$ \textit{une catégorie de modèles, associée à une catégorie à cofibrations. Alors sur le domaine} Diag \textit{formé des ensembles ordonnés finis, le prédérivateur} $\mathbf{D}=\mathbf{D}_{\mathcal{M},W}$ \textit{satisfait aux axiomes} Der (1, 2, 3, 4$^*$, 5$^*$, 6$^*$), \textit{et dualement. Donc si} $W$ \textit{est associé à la fois à une structure de cofibrations et de fibrations,} $\mathbf{D}$ \textit{est un dérivateur sur} Diag.

\msp{page 119}
\texttt{[Les six lignes qui suivent semblent hors contexte.]}

\textit{alors la flèche de changement de base}
\[ p^*f_!\longrightarrow f'_!q^* \]
\textit{entre foncteurs} $\mathrm{HOT}_W(X)\longrightarrow\mathrm{HOT}_W(S')$ \textit{est un isomorphisme.}

C'est pratiquement trivial, compte tenu, dans le cas $f$ $W$-lisse, que si $Z$ sur $X$ est $W$-lisse sur $X$, il l'est sur $S$ ; et dans le cas $p$ $W$-propre, en utilisant le dernier énoncé (``oubli'') dans \circled{6}, p. 199.

7.8. Je termine en mettant en vedette ce qui, à présent, me paraît la substance principale de tous ces ``rappels'' et comme la quintessence de la théorie des localiseurs fondamentaux
% margin: three vertical rules printed in the left margin, beside the footnotes and down to the foot of the page
% foot: Version 18 mars 2019
\printedpage{56}
% head: A. Grothendieck : Dérivateurs, Ch. XVI, Édition des Universités de Montpellier II et Paris VII
développée jusqu'à présent :
\[ \left\{ \begin{array}{rcl}
\text{Loc (3 bis)} & \Longrightarrow & \text{Der (4$^*$, 5$^*$)}\\[1ex]
\text{Loc (3 bis, 6 a)} & \Longrightarrow & \text{Der (1, 2, 3, 4$^*$, 5$^*$, 6$^*$)}\\[1ex]
\text{Loc (6$'$ a, 8, 9)} & \Longrightarrow & \text{tous les axiomes des dérivateurs (+ Loc (1 ter)).}
\end{array} \right. \]

Je considère les axiomes Loc (3 bis), Loc (8), Loc (9) comme des axiomes ``anodins'' - tous les localiseurs qu'on définit par voie cohomologique ou homotopique ordinaire, y satisfont. Il reste l'axiome Loc (6 a) ou Loc (6$'$ a) des $W$-cofibrations (forme faible),
\msp{page 120}
qui visiblement joue un rôle crucial dans la théorie de $\mathrm{HOT}_W$. Par contre, les axiomes Loc (4) (décomposition), Loc (5) (connexité) et Loc (7) et ses variantes (et bien sûr aussi Loc (6 b)) n'interviennent absolument pas ! C'est particulièrement impressionnant, en ce qui concerne l'axiome des $W$-fibrations Loc (7) et ses variantes (toutes équivalentes à Loc (7) en présence de Loc (8), à la seule exception de Loc (7 F)). Comme on a vu, Loc (7) n'est nullement impliqué par Loc (6 a, 8, 9) ($\mathbf{W}_0$ n'y satisfait pas), tout au plus se pourrait-il que Loc (7 F) le soit. Il en est de même de l'axiome de connexité Loc (5), comme on voit sur l'exemple de $W=\mathrm{Fl}(\mathrm{Cat})$, qui satisfait à tous les axiomes vus jusqu'à présent, à la seule exception
\msp{page 121}
de Loc (5). Quant à la question de l'indépendance de l'axiome Loc (4) par rapport à Loc (6 a, 8, 9), nous allons examiner ce genre de questions au paragraphe suivant.

L'ensemble d'axiomes Loc (6$'$ a, 8, 9) me paraît à présent si important, que j'ai envie de lui donner un nom, et d'appeler \textit{localiseur} (fondamental) \textit{spécial} dans Cat \footnote[57]{J'hésite si je ne dois pas rajouter l'axiome Loc (3 bis), qui assure que $W=W_{\mathbf{D}}$, $\mathbf{D}=\mathrm{HOT}_W$.}, un localiseur fondamental qui y satisfait. Ce sont donc aussi ceux pour lesquels nous savons que $\mathrm{HOT}_W$ est un vrai \textit{dérivateur.} Si on ne peut écarter comme ``pathologique'' une structure aussi riche et délicate qu'un dérivateur, il s'ensuit aussitôt que l'on ne peut non plus écarter les localiseurs spéciaux dans Cat (même si on affecte de regarder de haut les
\msp{page 122}
localiseurs fondamentaux généraux !). En effet, $W$ se récupère de façon unique à l'aide de $\mathbf{D}=\mathrm{HOT}_W$, ou de tout dérivateur équivalent à celui-ci, comme l'ensemble des $\mathbf{D}$-équivalences dans Cat \footnote[58]{Mais ceci n'est prouvé que moyennant Loc (3 bis).} : c'est le localiseur fondamental (correspondant à une notion d'équivalence ``universelle'', \textit{i.e.} dans la catégorie Cat de tous les ``espaces'' élémentaires) canoniquement associé au dérivateur. Ce n'est donc nullement un localiseur de référence fixe assez fin, tel $\mathbf{W}_\infty$, qui est apte à jouer un rôle !

Cela nous amène à nouveau à une question du genre déjà soulevé : si $W\subset\mathbf{W}_\infty$ est le plus petit localiseur qui satisfait à Loc (6$'$ a, 8) (ou Loc (3 bis, 6$'$ a, 8)), peut-on dire a) que $W$ est même \textit{spécial}, \textit{i.e.}
\msp{page 123}
satisfait également Loc (9) (accessibilité) ?, et de plus b) a-t-on $W\neq\mathbf{W}_\infty$ ? Je suis persuadé que c'est le cas (il y a moins de mérite pour ce $W$-là d'être différent de $\mathbf{W}_\infty$,
% foot: Version 18 mars 2019
\printedpage{57}
% head: A. Grothendieck : Dérivateurs, Ch. XVI, Édition des Universités de Montpellier II et Paris VII
que pour le $\tilde{\mathbf{W}}_\infty$ envisagé dans 6.4 (p. 87)).

7.9. Oubli : la question 1 bis (p. 98), sur l'existence d'une structure de Quillen close sur Cat, correspondant à $W$. Le seul résultat que je connaisse dans ce sens est celui de \textsc{Thomason}, construisant explicitement une telle structure de Quillen (de plus propre), dans le cas où $W=\mathbf{W}_\infty$. Je n'ai pas vraiment regardé la démonstration de \textsc{Thomason}, et ignore si celle-ci peut s'étendre à des localiseurs fondamentaux plus généraux, par exemple à ceux satisfaisant Loc (6$'$ a, 8, 9). S'il en était ainsi, cela prouverait d'ailleurs que ces localiseurs
\msp{page 124}
ont assez de cofibrations et assez de fibrations, donc qu'ils satisfont Loc (6 a) et Loc (7 F). (Pour ce qui est de Loc (6 a), il semble qu'on puisse le prouver directement très simplement, \textit{cf.} le paragraphe suivant. Mais ici encore, il faudrait regarder quelques exemples (équivalences homologiques ou cohomologiques, $\mathbf{W}_n$, \textit{i.e.} les $n$-équivalences d'homotopie), et regarder déjà si Loc (7 F) est satisfait.)

Il ne me semble pas que la théorie de \textsc{Thomason} ait été jusqu'à présent \texttt{[mots illisibles]} être très utile, techniquement parlant. Psychologiquement, je crois qu'il a été important, pour donner confiance que Cat est vraiment une excellente catégorie de modèles pour Hot. J'ai été vraiment enchanté et stimulé d'apprendre ce théorème. Mais la notion d'application de Dwyer (que j'ai apprise en 1983 dans l'article de \textsc{Thomason}) a été, techniquement, beaucoup plus importante pour moi : c'était exactement ce qui manquait à ma panoplie de notions et d'intuitions ``homotopiques'' dans Cat \dots
% foot: Version 18 mars 2019
\end{document}
